\documentclass[11pt]{amsart}
\usepackage[margin=1.1in]{geometry}
\usepackage{amsmath,amssymb,amsthm,mathtools}
\usepackage{booktabs}
\usepackage{microtype}
\usepackage{xcolor}
\usepackage[colorlinks=true,linkcolor=blue!50!black,citecolor=blue!50!black,urlcolor=blue!50!black]{hyperref}

\newtheorem{theorem}{Theorem}[section]
\newtheorem{proposition}[theorem]{Proposition}
\newtheorem{lemma}[theorem]{Lemma}
\newtheorem{corollary}[theorem]{Corollary}
\newtheorem{conjecture}[theorem]{Conjecture}
\theoremstyle{definition}
\newtheorem{definition}[theorem]{Definition}
\newtheorem{question}[theorem]{Question}
\theoremstyle{remark}
\newtheorem{remark}[theorem]{Remark}
\newtheorem{example}[theorem]{Example}

\newcommand{\Z}{\mathbb{Z}}
\newcommand{\F}{\mathbb{F}}
\newcommand{\A}{\mathcal{A}}
\newcommand{\Aut}{\operatorname{Aut}}
\newcommand{\End}{\operatorname{End}}
\newcommand{\Sym}{\operatorname{Sym}}
\newcommand{\Cent}{\mathcal{C}}
\newcommand{\sh}{\sigma}
\newcommand{\Lam}{\Lambda}
\newcommand{\Occ}{\operatorname{Occ}}

\title[Every finite group from a wall]{Every finite group from a wall: large automorphism groups of second-order cellular automata}
\author{Michael Dayah}
\email{michael@dayah.com}
\date{\today}

\begin{document}

\begin{abstract}
For an elementary rule $f$, the second-order (Fredkin) automaton $F_f$ on $(\{0,1\}^2)^{\Z}$ is reversible, and the automorphism group of its spacetime subshift is the centralizer of $F_f$ in the automorphism group of the full $4$-shift. Many of these automata have walls, patterns that stay fixed or periodic at a fixed position whatever surrounds them; they were listed by Takesue (1989), who observed that no information passes through them. We determine what walls imply for the automorphism group. If a reversible automaton of radius $1$ has a word that it maps to itself in every context, and at least two symbols can sit between two copies of it, then permuting the contents of trains of identical segments embeds $\bigoplus_mS_m$, and hence every finite group, in the automorphism group. With Takesue's fixed walls this applies to $32$ of the $88$ conjugacy classes of second-order elementary automata, including $54$R and $37$R, and we prove that no other class has a fixed word of any length. Variants of the construction, for walls that are periodic in time or that carry a bit invisible to their neighbours, give the same conclusion in seven further classes. Three affine classes, among them that of $90$R, which has no wall and no blocking word, reach it through marker permutations; altogether $42$ classes ($102$ rules) have every finite group in their automorphism group. For four other affine classes the isomorphism rigidity theorem of Kitchens and Schmidt shows that every automorphism is affine, and the automorphism group is $\Z^2$, $\Z^2\times\Z/2$ or, for rule $150$, $\Z^3$. So for affine rules the absence of a blocking word does not force the automorphism group to be virtually $\Z^2$.
\end{abstract}


\maketitle

\section{Introduction}\label{sec:intro}

Let $\A$ be a finite alphabet and $\sh$ the left shift on $\A^{\Z}$. By the Curtis--Hedlund--Lyndon theorem \cite{Hedlund1969} the endomorphisms of the full shift are the block maps, and the automorphisms are the invertible block maps. The group $\Aut(\A^{\Z})$ is enormous: it contains every finite group and free groups \cite{BLR1988}, and every countable, locally finite, residually finite group \cite{KimRoush1990}, and its center consists of the powers of the shift \cite{Ryan1972}. A reversible cellular automaton $F\in\Aut(\A^{\Z})$ determines a two-dimensional subshift, its \emph{spacetime} $X_F\subset\A^{\Z^2}$, whose rows are the configurations $F^n(c)$, $n\in\Z$ \cite{CFK2019}. Its automorphism group is the centralizer of $F$ in $\Aut(\A^{\Z})$ (Lemma~\ref{lem:spacetime}), and it always contains $\langle\sh,F\rangle$. This paper is about automata for which it is much larger.

The automata we consider are the second-order elementary rules, built by the second-order technique of Fredkin \cite{TM1990,Wolfram2002}; as a family they were introduced and studied systematically by Takesue \cite{Takesue1987,Takesue1989}, under the name of elementary reversible cellular automata. For an elementary rule $f\colon\{0,1\}^3\to\{0,1\}$ define, on the alphabet $\A=\{0,1\}^2$ with cells written $(x,y)$,
\begin{equation}\label{eq:F}
F_f(c)_i=\bigl(y_i,\;x_i\oplus f(y_{i-1},y_i,y_{i+1})\bigr).
\end{equation}
Reading $y$ as the current row and $x$ as the previous one, this is the recursion $y_{t+1}=y_{t-1}\oplus f(y_t)$, which is reversible for every $f$, with $F_f^{-1}(c)_i=\bigl(y_i\oplus f(x_{i-1},x_i,x_{i+1}),\,x_i\bigr)$. Following Takesue \cite{Takesue1987}, one writes $f$R for this automaton; $37$R is Wolfram's standard example of a reversible automaton with complex behaviour \cite[pp.~440 and 452--457]{Wolfram2002}. Under reflection and the cellwise complement the $256$ rules fall into $88$ conjugacy classes, the classes of Takesue \cite[Table~I]{Takesue1989}, and $\Aut(X_{F_f})$ depends only on the class (Section~\ref{sec:prelim}).

For some of these automata the automorphism group is as small as possible. For rule $30$, for all non-affine rules permutive in an outer variable, and for the classes of $63$ and $187$, $\Aut(X_{F_f})=\langle\sh,F_f\rangle\cong\Z^2$ \cite{PaperB}; the same holds for the classes of $110$, $14$, $46$ and $79$ \cite{PaperA}. The classification of all $88$ classes is the subject of \cite{PaperD}. Here we study the opposite behaviour. The spacetime $X_F$ has zero $\Z^2$-entropy, so Hochman's theorem that positive entropy forces every finite group into the automorphism group of a $\Z^2$ shift of finite type \cite{Hochman2010} never applies to it, and something else must supply the room. In every case we find, it is a local structure.

For most of the classes treated here the local structure is a \emph{wall} in the sense of Takesue \cite[Section~4]{Takesue1989}: a pattern, fixed or periodic in time, that appears at a fixed position irrespective of its surroundings. Takesue relates walls to local conservation laws, exhibits walls in $40$ of the $88$ classes (the $38$ classes of \cite[Fig.~2, Table~III]{Takesue1989} and those of $27$ and $59$), and observes that the two sides of a wall cannot communicate, so that a large system with walls is divided into parts that do not interact and is not mixing. This observation precedes the formalism of blocking words and equicontinuity points \cite{Kurka1997,BT2000}, of which walls are instances (Remark~\ref{hw:kurka}). The walls of the second-order elementary automata are therefore known, and we refer to \cite{Takesue1989} for them. What this paper adds is their consequence for the automorphism group: the embeddings of $\bigoplus_mS_m$ by permutations along trains of walls; a proof that the list of fixed walls is complete for words of every length; a bit hidden in words built from the periodic walls, with the flips that it gives; the census of rule $59$R at radius $2$; and the results on rule $54$. Where there are no walls, it adds the results on the linear and affine rules.

\subsection*{Fixed walls}
The basic mechanism is a word that the automaton fixes in every context.

\begin{theorem}\label{thm:walls}
Let $F$ be a reversible cellular automaton of radius $1$ on $\A^{\Z}$ and $W$ a nonempty word such that $F(c)$ contains $W$ at every position at which $c$ contains $W$. Let $C\subseteq\A$ be the set of symbols $s$ for which $W\,s\,W$ contains $W$ only at its two ends. If $|C|\ge2$, then $\Aut(X_F)$ contains a copy of $\bigoplus_{m\ge1}S_m$, hence of every finite group.
\end{theorem}

The copies of $W$ act as walls. A symbol between two walls evolves by a fixed permutation of $C$, whatever lies beyond the walls, so the contents of a run of such cells can be permuted without disturbing the dynamics. For the second-order elementary automata such words are known: the shortest fixed words of the $31$ classes of Table~\ref{tab:walls}, and the word $0330$ for the class of $133$, are the fixed walls (i)--(ix) of Takesue \cite[Fig.~2, Table~III and p.~391]{Takesue1989}, who examined the words of length at most $6$. We verify the hypothesis $|C|\ge2$ for them, and we prove that no other class has a fixed word of any length (Proposition~\ref{cl:prop:133}).

\begin{corollary}\label{cor:48}
For the $81$ elementary rules $f$ in the $31$ classes listed in Table~\ref{tab:walls}, $\Aut(X_{F_f})$ contains every finite group. These include $54$, $37$, $22$ and $73$. For the class $\{133,161\}$ the same holds, through the fixed word $0330$ of length $4$. No other elementary rule has a word of any length fixed by $F_f$ in every context.
\end{corollary}

\subsection*{Rule 54}
Rule $54$ is, with rule $110$, one of the two elementary rules for which Wolfram reports class~4 behaviour, rule $110$ as the clear example and rule $54$ as a possible one \cite[pp.~693--697]{Wolfram2002}. The two second-order automata behave in opposite ways: $\Aut(X_{F_{110}})=\langle\sh,F_{110}\rangle$ \cite{PaperA}, while for $54$R the pattern $33$ (writing $3$ for the cell $(1,1)$) is a wall, the wall (iv) of \cite[Table~III]{Takesue1989}, and the walls give the following.

\begin{theorem}\label{thm:finite}
For every $m\ge1$ there is an injective homomorphism $\Psi_m\colon S_m\to\Aut(X_{F_{54}})$ whose image consists of block maps of radius at most $3m+3$, and the images for different $m$ commute, so $\bigoplus_{m\ge1}S_m$ embeds in $\Aut(X_{F_{54}})$.
\end{theorem}

The relabellings of the segments between walls are automorphisms of smaller radius: a permutation of the contents of the segments of a given length $L$ is an automorphism whenever it commutes with the permutation that $F$ induces on these contents, and these maps give an embedding of the direct sum over $L$ of the centralizers of the induced permutations (Theorem~\ref{thm:54}). The smallest of them is $D$, of radius $2$ and order $3$, which acts as the $3$-cycle $0\to1\to2\to0$ on every cell that is a single-cell segment between two walls and as the identity elsewhere. At radius $2$ an exhaustive computation shows that the walls account for everything: the automorphisms of radius at most $2$ are exactly the $13$ maps $\sh^aF^b$ with $|a|+|b|\le2$ and the $10$ maps $F^bD^k$ with $|b|\le 2$, $k\in\{1,2\}$ (Theorem~\ref{thm:54census}). At radius $3$ this is no longer so: there is an involution $\tau$ of radius $3$ that commutes with $F$, moves the boundary between a long wall and a single cell next to it, and does not commute with $D$ (Theorem~\ref{thm:tau}). A partial enumeration of the automorphisms of radius $3$ (Section~\ref{sec:54r3}) finds that every automorphism it produces differs from some $\sh^aF^b$ only on windows that contain a wall, and that most of them, like $\tau$, do not preserve walls.

\subsection*{Walls that are not fixed}
A wall need not be fixed in every context. Takesue's list contains walls that are periodic in time and not fixed, in the classes of $27$, $59$, $179$, $127$, $223$, $126$ and $219$ \cite[Section~4, (4.1)--(4.3), Fig.~2 and Table~III, walls (xii), (xiii), (xv), (xvi)]{Takesue1989}. None of these seven classes has a word fixed in every context, so Theorem~\ref{thm:walls} does not apply to them. In Section~\ref{sec:hidden} we give a general mechanism (Theorem~\ref{hw:general}) for walls whose image never depends on the context, except at cells that do not matter, and which carry one bit that $F$ cannot see. Permuting these hidden bits along chains of walls gives every finite group again. The mechanism applies to words built from Takesue's periodic walls in the classes of $59$, $27$, $127$, $179$ and $223$ (Theorems~\ref{hw:59} and~\ref{hw:27}, Proposition~\ref{hw:tp}). For $219$R and $126$R the walls rotate through the alphabet in every context, and permuting the cells between them again embeds every finite group (Theorem~\ref{hw:219}, Proposition~\ref{hw:126}).

\subsection*{The affine rules 90 and 165}
Rule $90$ has no blocking word \cite{Kurka1997,BT2000}, hence no wall of any of the kinds above (Corollary~\ref{lin:noblock}), and the automorphisms of $X_{F_{90}}$ of radius at most $2$ are all affine (Section~\ref{sec:linear}). Nevertheless $F_{90}$ has two factor maps onto the binary full shift that carry $F$ to $\sh^{\pm1}$, with matching embeddings (Theorem~\ref{thm:90}), and these carry marker automorphisms of the binary full shift into $X_{F_{90}}$: by Theorem~\ref{lin:movers}, $(\bigoplus_mS_m)^2$ embeds in $\Aut(X_{F_{90}})$ and in $\Aut(X_{F_{165}})$. Theorem~\ref{thm:90} also describes the block maps of radius at most $3$ that commute with $F_{90}$, invertible or not: they are the sums of an affine map and of maps that factor through the two factor maps.

\subsection*{The other affine rules}
For the affine rules whose characteristic polynomial $t^2+pt+1$ is irreducible and whose spacetime action is mixing, the isomorphism rigidity theorem of Kitchens and Schmidt \cite{KS2000} forces every automorphism to be affine (Lemma~\ref{aff:lem:R}), and the automorphism group becomes a unit group, which we compute. Write $\kappa$ for the cellwise complement.

\begin{theorem}\label{thm:affine}
Every automorphism of $X_{F_f}$ is affine for the rules $f$ of the four classes below, and:
\begin{enumerate}
\item for $f\in\{170,240\}$, $\Aut(X_{F_f})=\langle F,h\rangle\cong\Z^2$ with $h=\mathrm{id}+F$, and $h^2=\sh^{\pm1}F$, so $\langle\sh,F\rangle$ has index $2$; for $f\in\{15,85\}$ the same holds with $h(c)=c+F(c)+(1,1)^{\Z}$;
\item for $f\in\{60,102\}$, $\Aut(X_{F_f})=\langle\sh,F\rangle\times\langle\kappa\rangle\cong\Z^2\times\Z/2$;
\item for $f\in\{153,195\}$, $\Aut(X_{F_f})=\langle\sh,F\rangle\times\langle\kappa\rangle\cong\Z^2\times\Z/2$;
\item for $f=150$, $\Aut(X_{F_f})=\langle\sh,F,h\rangle\cong\Z^3$ with $h=\mathrm{id}+\sh F$; for $f=105$ the same holds with $h(c)=c+\sh F(c)+(1,1)^{\Z}$.
\end{enumerate}
\end{theorem}

For $170$ the spacetime, with the action of $F$ and $\mathrm{id}+F$, is conjugate to Ledrappier's three-dot shift, with $\langle\sh,F\rangle$ acting as a subgroup of index $2$ of its shift action (Remark~\ref{aff:rem:ledrappier}). The hypotheses fail exactly for the other affine classes, those of $0$, $51$, $90$, $165$ and $255$, and each of them has non-affine automorphisms (Remark~\ref{aff:rem:consist}). For $255$, $F$ adds $1$ modulo $4$ in every cell, and $\Aut(X_{F_{255}})\cong\Z/4\times\mathcal G_4$, where $\mathcal G_4$ is a group of automorphisms of the full shift on $\Z/4$ that contains $\bigoplus_mS_m$ (Proposition~\ref{aff:prop:255}).

\subsection*{Summary}
Collecting these results, the following $42$ classes of second-order elementary automata, comprising $102$ rules, have the property that $\bigoplus_{m\ge1}S_m$ embeds in $\Aut(X_{F_f})$, so that $\Aut(X_{F_f})$ contains every finite group and is not virtually $\Z^n$ for any $n$:
\begin{itemize}
\item the $31$ classes ($81$ rules) of Table~\ref{tab:walls}, with a word of length at most $3$ fixed in every context, and the class of $133$ ($2$ rules), with one of length $4$ (Theorem~\ref{thm:walls}, Corollary~\ref{cor:48});
\item the classes of $59$ and $27$ ($8$ rules), with hidden-bit walls (Theorems~\ref{hw:59} and~\ref{hw:27});
\item the classes of $127$, $179$ and $223$ ($6$ rules), with time-periodic walls (Proposition~\ref{hw:tp});
\item the classes of $219$ and $126$ ($2$ rules), with rotating walls (Theorem~\ref{hw:219}, Proposition~\ref{hw:126});
\item the affine classes of $90$ and $165$ ($2$ rules), where $(\bigoplus_mS_m)^2$ embeds (Theorem~\ref{lin:movers});
\item the affine class of $255$ ($1$ rule), through marker permutations of the difference sequence (Proposition~\ref{aff:prop:255}).
\end{itemize}
The first four items comprise $39$ classes. Their walls are those of Takesue or are built from them, and with the class of $255$, which also has walls in his list, they are the $40$ classes in which he exhibits walls \cite[Section~4]{Takesue1989}.
Apart from the description $\Aut(X_{F_{255}})\cong\Z/4\times\mathcal G_4$, we do not determine the whole automorphism group in these cases. For the other $46$ classes, the status of each is recorded in \cite{PaperD}: for the class of $95$ the automorphism group is determined in \cite{PaperE}, and it also contains $\bigoplus_{m\ge1}S_m$, through automorphisms that permute moving particles and not through a wall; $11$ classes have $\Aut(X_{F_f})=\langle\sh,F_f\rangle\cong\Z^2$ \cite{PaperA,PaperB,PaperE}, the four affine classes of Theorem~\ref{thm:affine} have the groups given there, two further classes have the cellwise symmetry $\kappa$, and for $28$ classes an exhaustive search finds no automorphism outside $\langle\sh,F_f\rangle$ of radius at most $3$.

\subsection*{Blocking words and linearity}
In the nonlinear examples of this paper the automorphisms outside $\langle\sh,F\rangle$ all come from blocking words with a degree of freedom next to the blocking column (Remark~\ref{hw:kurka}). The following remark records what the affine rules say about the converse.

\begin{remark}\label{rem:dichotomy}
Without a hypothesis on $f$, the absence of a blocking word does not force $\Aut(X_{F_f})$ to be virtually $\Z^2$. The affine rules $90$ and $165$ have no blocking word (Corollary~\ref{lin:noblock}), yet $(\bigoplus_mS_m)^2$ embeds in their automorphism groups (Theorem~\ref{lin:movers}). Rule $150$ is a counterexample of a different kind. It has no blocking word, and its commuting block maps of radius at most $3$ are all linear (Proposition~\ref{prop:150}); but $\mathrm{id}+\sh F_{150}$ is an automorphism outside $\langle\sh,F_{150}\rangle$, and $\langle\sh,F_{150},\mathrm{id}+\sh F_{150}\rangle\cong\Z^3$ (Proposition~\ref{prop:150Z3}), so $\Aut(X_{F_{150}})$ is not virtually $\Z^2$; in fact $\Aut(X_{F_{150}})=\langle\sh,F_{150},\mathrm{id}+\sh F_{150}\rangle$ (Theorem~\ref{thm:affine}). In all three the room comes from linearity, not from a local structure.
\end{remark}

Nor does the absence of a blocking word force $\Aut(X_{F_f})$ to be virtually $\Z^2$ when $f$ is not affine: the rule $95$ is not affine, $F_{95}$ has no blocking word, and $\Aut(X_{F_{95}})$ contains every finite group \cite{PaperE}. There the room comes from automorphisms that permute particles moving at light speed, and neither from a local structure nor from linearity. We do not know another example of this kind among the elementary rules.

\begin{question}\label{q:dichotomy}
Let $f$ be an elementary rule that is not affine and does not belong to the class of $95$. If the second-order automaton $F_f$ has no blocking word \cite{Kurka1997,BT2000}, is $\Aut(X_{F_{f}})$ virtually $\Z^2$?
\end{question}

Having no blocking word is stronger than having no word fixed in every context; the rules of Section~\ref{sec:hidden} show that the difference matters (Remark~\ref{hw:kurka}). For $110$R and $30$R, whose automorphism groups are $\Z^2$ \cite{PaperA,PaperB}, a single changed cell produces a difference that spreads in both directions at linear speed, whereas in $54$R it spreads only a few cells (Table~\ref{tab:spread}).

\subsection*{Related work}
\emph{Large automorphism groups.} The automorphism group of a mixing one-dimensional shift of finite type contains every finite group \cite{BLR1988}. In two dimensions the same holds under strong irreducibility \cite{Ward1994} and under positive entropy \cite{Hochman2010}. Hochman also notes that zero entropy does not prevent a large automorphism group, with the example of a shift whose columns are constant; that shift is the spacetime of the identity automaton. Theorems~\ref{thm:walls} and~\ref{thm:finite} are of the same kind, but for nonlinear automata of infinite order, where the room comes from a conserved local structure. For linear automata the algebraic theory of Kitchens and Schmidt applies \cite{KS1989}; their isomorphism rigidity theorem \cite{KS2000} is what determines the automorphism groups in Theorem~\ref{thm:affine}, and for Ledrappier's shift \cite{Ledrappier1978} the automorphism group is $\Z^2$ \cite{Ward1994}, \cite[Theorem~7]{BRY2018}, \cite[Example~8.20]{Schmidt2002}.

\emph{Centralizers and blocking words.} Ryan's theorem identifies the centralizer of the whole automorphism group of a full shift with the powers of the shift \cite{Ryan1972}, and finitary versions find finite sets of automorphisms with the same centralizer \cite{Salo2019Ryan,Kopra2022}; a spacetime automorphism group is the centralizer of a single automaton, which these results do not cover. The term blocking word is due to Blanchard and Tisseur \cite{BT2000}; the relation of blocking words to equicontinuity is in the proof of \cite[Theorem~3]{Kurka1997}. For the second-order elementary automata, that a wall separates the two sides of a configuration was observed by Takesue \cite{Takesue1987}, \cite[Section~4]{Takesue1989}, before these notions were introduced.

\emph{Rule 54.} The particles and collisions of the elementary rule $54$ were catalogued by Boccara, Nasser and Roger \cite{BNR1991} and by Mart\'inez, Adamatzky, Chen and Chua \cite{MAM2013}. A different reversible automaton, introduced by Bobenko, Bordemann, Gunn and Pinkall \cite{BBGP1993} and studied extensively as an integrable model of statistical mechanics \cite{BKP2021}, is also called Rule~54 because it uses the same local function $b\oplus(a\vee c)$, but it updates the two sublattices alternately in place. In the present notation it is one of the two sublattices of a spacetime diagram of $250$R \cite{BKP2021}, which is the conjugate of $95$R under the cellwise complement: it is the half system of the class of $95$, whose automorphism group is determined in \cite{PaperE}. It is not the second-order automaton $F_{54}$ studied here, which has the extra term from the previous row; its solitons move at speed $\pm1$, while in $F_{54}$ the structures that carry the automorphisms are the stationary walls.

\subsection*{Organisation}
Section~\ref{sec:prelim} fixes notation, proves Lemma~\ref{lem:spacetime}, and shows that $\langle\sh,F_f\rangle\cong\Z^2$ unless $f$ reads only its middle argument. Section~\ref{sec:method} describes the SAT encoding of the commutation condition, the invertibility test and the verification. Section~\ref{sec:54} treats rule $54$: it proves Theorem~\ref{thm:finite}, and it states and proves Theorems~\ref{thm:54} and~\ref{thm:tau}. Section~\ref{sec:walls} proves Theorem~\ref{thm:walls} and Corollary~\ref{cor:48}, including the class of $133$. Section~\ref{sec:hidden} treats walls with hidden bits, rules $59$ and $27$, the time-periodic walls of rules $127$, $179$ and $223$, the rotating walls of rules $219$ and $126$, and blocking words. Section~\ref{sec:linear} treats the linear rules $90$ and $150$ (Theorem~\ref{thm:90}, Proposition~\ref{prop:150}) and proves that $(\bigoplus_mS_m)^2$ embeds in the automorphism groups of $90$R and $165$R. Section~\ref{sec:affine} proves Theorem~\ref{thm:affine} and treats rule $255$. Section~\ref{sec:questions} lists open questions. The appendix describes the computations.

\section{Preliminaries}\label{sec:prelim}

Throughout, $\A=\{0,1\}^2$, and a cell $(x,y)\in\A$ is also written as the integer $2x+y\in\{0,1,2,3\}$; thus $3=(1,1)$ and $0=(0,0)$. The left shift is $\sh(c)_i=c_{i+1}$. A \emph{block map of radius $R$} is a map $g\colon\A^{\Z}\to\A^{\Z}$ of the form $g(c)_i=\gamma(c_{i-R},\dots,c_{i+R})$ for a local rule $\gamma\colon\A^{2R+1}\to\A$; we identify $g$ with the lookup table of $\gamma$, which has $4^{2R+1}$ entries of two bits each. A block map of radius $R$ is also a block map of every larger radius. By Curtis--Hedlund--Lyndon, $\End(\A^{\Z})$ is the monoid of block maps and $\Aut(\A^{\Z})$ the group of invertible ones. We write
\[
\Cent_R(F)=\{g\ \text{a block map of radius}\le R:\ g\circ F=F\circ g\},\qquad \Cent(F)=\bigcup_R\Cent_R(F),
\]
for the centralizer of $F$ in $\End(\A^{\Z})$, and $\Cent^{\times}_R(F)$, $\Cent^{\times}(F)$ for the invertible elements. For a configuration $c$ and integers $i<j$ we write $c_{[i,j)}=c_ic_{i+1}\cdots c_{j-1}$, and $c_{[i,j]}=c_{[i,j+1)}$.

For an elementary rule $f$, identified with its Wolfram number $\sum_{a,b,c}f(a,b,c)2^{4a+2b+c}$, the automaton $F_f$ is defined by \eqref{eq:F}. It has radius $1$, and $F_f^b$ has radius at most $|b|$, so $\sh^aF_f^b\in\Cent_{|a|+|b|}(F_f)$.

\begin{definition}
The \emph{spacetime subshift} of a reversible cellular automaton $F$ on $\A^{\Z}$ is
\[
X_F=\bigl\{(F^n(c))_{n\in\Z}:\ c\in\A^{\Z}\bigr\}\subset\A^{\Z^2},
\]
with the $\Z^2$-action generated by the horizontal shift $\sh$ and the vertical shift $\tau$, $(\tau z)_{n}=z_{n+1}$. An \emph{automorphism} of $X_F$ is a homeomorphism of $X_F$ commuting with both.
\end{definition}

\begin{lemma}\label{lem:spacetime}
For reversible $F$ the map $c\mapsto(F^n(c))_{n\in\Z}$ is a homeomorphism $\A^{\Z}\to X_F$ conjugating $(\sh,F)$ to $(\sh,\tau)$, and it induces an isomorphism
\[
\Aut(X_F)\;\cong\;\Cent^{\times}(F)=\{g\in\Aut(\A^{\Z}):\ gF=Fg\}.
\]
Under this isomorphism $\langle\sh,\tau\rangle$ corresponds to $\langle\sh,F\rangle$.
\end{lemma}

\begin{proof}
The map is continuous, injective (the $0$-th row recovers $c$), and surjective onto $X_F$ by definition, hence a homeomorphism of compact spaces; it carries $\sh$ to $\sh$ and $F$ to $\tau$ since $F^{n}(F(c))=F^{n+1}(c)$. A homeomorphism of $\A^{\Z}$ commuting with $\sh$ is a block map by Curtis--Hedlund--Lyndon, and commuting with $\tau$ means commuting with $F$.
\end{proof}

So determining $\Aut(X_F)$ means determining which invertible block maps commute with $F$. The non-invertible elements of $\Cent(F)$ are the endomorphisms of $X_F$ that are not automorphisms; the constant map onto a fixed point of $F$ is always one. For every rule with $f(0,0,0)=f(1,1,1)=0$, which includes $54$ and $90$, both $(0,0)^{\Z}$ and $(1,1)^{\Z}$ are fixed by $F_f$; for rule $150$, which has $f(1,1,1)=1$, only $(0,0)^{\Z}$ is.

\begin{sloppypar}
\emph{Conjugacies.} Let $\mathrm{rev}(c)_i=c_{-i}$ be the reflection and $\kappa(x,y)=(1-x,1-y)$ the cellwise complement. By \eqref{eq:F}, $\mathrm{rev}\,F_f\,\mathrm{rev}=F_{f'}$ with $f'(a,b,c)=f(c,b,a)$, and $\kappa F_f\kappa=F_{f''}$ with $f''(a,b,c)=f(1-a,1-b,1-c)$. Both $\mathrm{rev}$ and $\kappa$ are involutions that act cellwise up to reflection, so conjugation by them maps $\Cent_R(F_f)$ bijectively onto $\Cent_R(F_{f'})$, respectively $\Cent_R(F_{f''})$, preserves invertibility and maps $\langle\sh,F_f\rangle$ onto the corresponding group. So $\Aut(X_{F_f})$, and every statement below about block maps of bounded radius, depends only on the class of $f$ under these two operations. There are $88$ classes \cite{Takesue1987}, \cite[Table~I]{Takesue1989}; one representative of each appears in Table~\ref{tab:walls}. Wolfram's complement $1\oplus f(1-a,1-b,1-c)$ is not a conjugacy of second-order automata, so these are not the classes of the elementary rules \cite[p.~377 and footnote~3]{Takesue1989}.
\end{sloppypar}

\subsection{The group generated by the shift and the automaton}

The automorphisms constructed below lie outside $\langle\sh,F_f\rangle$. To make statements about the group they generate together with $\sh$ and $F_f$ we need the structure of $\langle\sh,F_f\rangle$ itself. For $c\in\A^{\Z}$ let $Y_c(t,i)$ be the $y$-bit of $F^t(c)_i$; then the $x$-bit of $F^t(c)_i$ is $Y_c(t-1,i)$, and by \eqref{eq:F}
\[
Y(t+1,i)\oplus Y(t-1,i)=f\bigl(Y(t,i-1),Y(t,i),Y(t,i+1)\bigr)\qquad\text{for all }(t,i).
\]
A block map $g$ \emph{depends on} a bit of the cell $c_i$ (for the output bit $y(g(c)_0)$, say) if two configurations that differ only in that bit have different outputs. For a nonconstant block map $g$ let $r^+(g)$ be the largest and $-r^-(g)$ the smallest $i$ such that $g(c)_0$ depends on $c_i$. Then $r^\pm(\sh^ag)=r^\pm(g)\pm a$, and $r^\pm(F_f^n)\le|n|$ because $F_f^n(c)_0$ depends only on $c_{[-|n|,|n|]}$.

\begin{lemma}[free staircases]\label{cl:lem:stair}
Let $f$ be any elementary rule and $n\ge1$. For every assignment of bits to the $2n$ points $(k,n-k-2)$ and $(k,n-k-1)$, $0\le k<n$, there is a configuration $c$ whose diagram $Y_c$ takes the assigned values at these points.
\end{lemma}

\begin{proof}
The points lie on the two diagonals $d_m=\{(t,i):t+i=m\}$ with $m=n-2$ and $m=n-1$. Put $K=2n+2$, $\tau(m)=m-n+1+K$, and let $\Omega$ be the set of $(t,i)$ with $n-1-K\le t+i\le n-1$ and $-1\le t\le\tau(t+i)$; so $\Omega\cap d_m=\{(t,m-t):-1\le t\le\tau(m)\}$ for $n-1-K\le m\le n-1$, and $\tau(m+1)=\tau(m)+1$. We define $Y\colon\Omega\to\{0,1\}$, one diagonal at a time, from right to left. On $d_{n-1}$ and $d_{n-2}$ let $Y$ take the assigned values at the $2n$ points, which lie in $\Omega$ because $k\le n-1<\tau(n-2)$, and the value $0$ elsewhere. For $m=n-3,n-4,\dots,n-1-K$ in turn put $Y(\tau(m),m-\tau(m))=0$ and then, for $t=\tau(m),\tau(m)-1,\dots,0$,
\begin{equation}\label{cl:eq:fill}
Y(t-1,m-t+1)=Y(t+1,m-t+1)\oplus f\bigl(Y(t,m-t),\,Y(t,m-t+1),\,Y(t,m-t+2)\bigr).
\end{equation}
The point $(t,m-t)$ of $d_m$ has just been defined, and $(t,m-t+1)\in d_{m+1}$, $(t+1,m-t+1)$ and $(t,m-t+2)\in d_{m+2}$ have been defined before, because $t\le\tau(m)<\tau(m+1)$ and $t+1\le\tau(m+2)$. So every point of $\Omega$ receives a value, and \eqref{cl:eq:fill} is the spacetime relation at the point $(t,m-t+1)$.

Conversely, suppose that the five points $(t\pm1,i)$, $(t,i\pm1)$ of the relation at $(t,i)$ all lie in $\Omega$, and put $m=t+i-1$. Then $(t,i-1)\in d_m$ and $(t,i+1)\in d_{m+2}$ give $n-1-K\le m\le n-3$ and $t\le\tau(m)$, and $(t-1,i)\in\Omega$ gives $t\ge0$. So the relation at $(t,i)$ is \eqref{cl:eq:fill} for this $m$ and $t$, and it holds. Thus $Y$ satisfies every instance of the spacetime relation whose five points lie in $\Omega$.

Let $c$ be any configuration with $c_i=(Y(-1,i),Y(0,i))$ for $n-K\le i\le n-1$; these points lie in $\Omega$. We claim that $Y_c(t,i)=Y(t,i)$ for $t\in\{-1,0\}$ and $n-K\le i\le n-1$, and for $t\ge1$ and $n-K+t\le i\le n-1-t$. For $t\le0$ this is the choice of $c$. For $t\ge1$, $Y_c(t,i)=Y_c(t-2,i)\oplus f(Y_c(t-1,i-1),Y_c(t-1,i),Y_c(t-1,i+1))$, the four points on the right lie in the ranges already covered for the rows $t-2$ and $t-1$, and the five points of the relation at $(t-1,i)$ lie in $\Omega$ (their diagonals are between $t+i-2\ge n-K$ and $t+i\le n-1$, and their rows between $t-2\ge-1$ and $t\le\tau(t+i)$, the last since $i\ge n-K+t>n-1-K$); so $Y$ satisfies the same relation there, and the claim follows by induction on $t$. The assigned points $(k,n-k-2)$, $(k,n-k-1)$ lie in these ranges: for $k=0$ in row $0$, and for $k\ge1$ because $n-K+k\le n-k-2$, that is $K\ge2k+2$. So $Y_c$ takes the assigned values there.
\end{proof}

\begin{theorem}\label{cl:thm:A}
Let $f$ be an elementary rule, $F=F_f$ and $n\ge1$.
\begin{enumerate}
\item If $f$ depends on its third argument, then $y(F^n(c)_0)$ depends on $y(c_n)$, and $x(F^{-n}(c)_0)$ depends on $x(c_n)$. So $r^+(F^n)=r^+(F^{-n})=n$.
\item If $f$ depends on its first argument, then $y(F^n(c)_0)$ depends on $y(c_{-n})$, and $x(F^{-n}(c)_0)$ depends on $x(c_{-n})$. So $r^-(F^n)=r^-(F^{-n})=n$.
\item The map $(a,b)\mapsto\sh^aF^b$ is an isomorphism $\Z^2\to\langle\sh,F\rangle$ if and only if $f$ is not a function of its middle argument alone.
\end{enumerate}
\end{theorem}

\begin{proof}
(1) Choose $(a^*,b^*)$ with $f(a^*,b^*,0)\ne f(a^*,b^*,1)$. By Lemma~\ref{cl:lem:stair} there is a configuration $c$ with $Y_c(k,n-k-2)=a^*$ and $Y_c(k,n-k-1)=b^*$ for $0\le k<n$. Let $c'$ be $c$ with $y(c_n)$ complemented, and $\Delta=Y_c\oplus Y_{c'}$. We show by induction on $k$, $0\le k\le n$, that $\Delta(k,i)=0$ for $i<n-k$, $\Delta(k,n-k)=1$, and $\Delta(k-1,i)=0$ for $i<n-k+1$. For $k=0$ this holds because $c$ and $c'$ differ only in $y(c_n)$. For the step from $k<n$ to $k+1$, the relation at $(k,i)$ gives $\Delta(k+1,i)=0$ for $i<n-k-1$, since all four inputs $\Delta(k-1,i)$, $\Delta(k,i-1)$, $\Delta(k,i)$, $\Delta(k,i+1)$ vanish. At $i=n-k-1$ the inputs $\Delta(k-1,n-k-1)$, $\Delta(k,n-k-2)$ and $\Delta(k,n-k-1)$ vanish, so both diagrams have the values $a^*,b^*$ at $(k,n-k-2)$, $(k,n-k-1)$, while $\Delta(k,n-k)=1$; hence $\Delta(k+1,n-k-1)=f(a^*,b^*,0)\oplus f(a^*,b^*,1)=1$. For $k=n$ this says $Y_c(n,0)\ne Y_{c'}(n,0)$, that is $y(F^n(c)_0)\ne y(F^n(c')_0)$. For $F^{-n}$ use the time reversal $\rho(x,y)=(y,x)$, which satisfies $\rho F\rho=F^{-1}$ for every $f$ \cite[p.~2500]{Takesue1987}: $x(F^{-n}(c)_0)=y(F^n(\rho c)_0)$ depends on $y((\rho c)_n)=x(c_n)$. With $r^+(F^{\pm n})\le n$ this gives $r^+(F^{\pm n})=n$.

(2) The reflection conjugates $F_f$ to $F_{f'}$ with $f'(a,b,c)=f(c,b,a)$, and $f'$ depends on its third argument; apply (1) to $f'$.

(3) Suppose $f$ depends on its third argument. By (1), $r^+(\sh^aF^b)=a+|b|$ for all $(a,b)$. If $\sh^aF^b=\mathrm{id}$, then also $\sh^{-a}F^{-b}=\mathrm{id}$, so $a+|b|=0=-a+|b|$ and $a=b=0$. Since $\sh$ and $F$ commute, $(a,b)\mapsto\sh^aF^b$ is a homomorphism onto $\langle\sh,F\rangle$, and it is injective. If $f$ depends on its first argument, the same argument with $r^-(\sh^aF^b)=|b|-a$ applies. If $f$ depends on neither, then $f(a,b,c)=g(b)$, and $F(c)_i=(y_i,x_i\oplus g(y_i))$ applies one permutation of $\A$ to every cell, so $F$ has finite order and the map is not injective.
\end{proof}

\begin{proposition}\label{cl:prop:kappa}
For $f\in\{126,219\}$ the cellwise complement $\kappa$ commutes with $F_f$ and $\kappa\notin\langle\sh,F_f\rangle$, so $\langle\sh,F_f,\kappa\rangle=\langle\sh,F_f\rangle\times\langle\kappa\rangle\cong\Z^2\times\Z/2$.
\end{proposition}

\begin{proof}
Since $\kappa F_f\kappa=F_{f''}$ with $f''(a,b,c)=f(1-a,1-b,1-c)$, $\kappa$ commutes with $F_f$ if and only if $f(1-a,1-b,1-c)=f(a,b,c)$ for all $a,b,c$, that is, if and only if $f$ is self-conjugate in the sense of \cite[Appendix~A]{Takesue1989}, which holds for $126$ (which vanishes exactly on $000$ and $111$) and for $219$ (which vanishes exactly on $010$ and $101$). For $126$, $f(0,0,0)=0$, so $(0,0)^{\Z}$ is fixed by $\sh$ and $F_f$, hence by every element of $\langle\sh,F_f\rangle$, while $\kappa$ maps it to $(1,1)^{\Z}$. For $219$, the orbit of $(01)^{\infty}$ under $\langle\sh,F_{219}\rangle$ consists of the four configurations of period $2$ with words $01,10,02,20$, and $\kappa$ maps $(01)^{\infty}$ to $(32)^{\infty}$, which is not among them. Since an element of $\langle\sh,F_f\rangle$ maps each configuration into its orbit, $\kappa\notin\langle\sh,F_f\rangle$. Both rules depend on their first and third arguments, so $\langle\sh,F_f\rangle\cong\Z^2$ by Theorem~\ref{cl:thm:A}(3). As $\kappa$ is an involution commuting with $\sh$ and $F_f$, and $\langle\sh,F_f\rangle\cong\Z^2$ is torsion-free, $\langle\kappa\rangle\cap\langle\sh,F_f\rangle=\{\mathrm{id}\}$, which gives the direct product.
\end{proof}

\section{The computation}\label{sec:method}

Several statements below are finite computations of $\Cent_R(F)$ or $\Cent^{\times}_R(F)$ for small $R$. This section describes the method; Appendix~\ref{app:comp} lists what was computed.

\subsection{Commutation as a finite constraint}
Let $g$ be a block map of radius $R$ with local rule $\gamma$, and $F=F_f$. Both $g\circ F$ and $F\circ g$ are block maps of radius $R+1$, so they are equal if and only if their local rules agree on every window $w\in\A^{2R+3}$. Write $n=2R+1$ and, for a window $w=(w_0,\dots,w_{2R+2})$, let $e_1,e_2,e_3$ be the three subwords of length $n$ starting at positions $0,1,2$ and let $e_L=F(w)$ be the word of length $n$ obtained by applying the local rule of $F$ to $w$. Then
\[
(g\circ F)(w)=\gamma(e_L),\qquad (F\circ g)(w)=\Bigl(\gamma(e_2)_y,\ \gamma(e_2)_x\oplus f\bigl(\gamma(e_1)_y,\gamma(e_2)_y,\gamma(e_3)_y\bigr)\Bigr),
\]
where $\gamma(e)_x,\gamma(e)_y$ are the two output bits. Writing $X_e,Y_e$ for the two unknown bits of table entry $e$, each window contributes the two constraints
\begin{equation}\label{eq:constraints}
X_{e_L}=Y_{e_2},\qquad Y_{e_L}\oplus X_{e_2}=f(Y_{e_1},Y_{e_2},Y_{e_3}).
\end{equation}
The first is two binary clauses; the second is a constraint on at most five variables, encoded by its truth table in at most sixteen clauses. Altogether there are $2\cdot4^{2R+1}$ variables and at most $18\cdot4^{2R+3}$ clauses: $17{,}680$ clauses at $R=1$ and $2.9\times10^5$ at $R=2$.

\subsection{Enumeration and verification}
All solutions are enumerated by repeated solving with blocking clauses, using PySAT \cite{PySAT} with CaDiCaL~1.5.3 \cite{CaDiCaL}. Every table returned by the solver is re-checked, independently of the CNF, by evaluating $g\circ F$ and $F\circ g$ on all $4^{2R+3}$ windows. Each solution is then classified: it is matched against the tables of the maps $\sh^aF^b$ with $|a|+|b|\le R$, which are built directly by composing local rules, and against the constant maps. The program asserts that every $\sh^aF^b$ with $|a|+|b|\le R$ occurs among the solutions, which checks the encoding against the only maps known in advance to commute.

\subsection{Invertibility}
A block map is an automorphism of $\A^{\Z}$ if and only if it is injective. Injectivity of a block map of radius $R$ is decided exactly by the pair graph: its vertices are the pairs $(u,v)$ of words of length $2R$, with an edge $(u,v)\to(u',v')$ labelled $(a,b)\in\A^2$ whenever $u'$ and $v'$ are $ua$ and $vb$ with the first letter dropped, and $\gamma(ua)=\gamma(vb)$. Then $g$ fails to be injective if and only if a bi-infinite path passes through an off-diagonal vertex, which is decided by iteratively deleting vertices without a predecessor or without a successor. The test was implemented twice, in pure Python and vectorised, and the two implementations were checked to agree at $R=1$; at $R\le2$ every map found injective was also checked to be injective on all spatially periodic configurations of period at most $2R+4$.

When $\Cent_R(F)$ is too large to enumerate, only its invertible part is enumerated. An invertible block map is surjective, and a surjective block map of radius $R$ is \emph{balanced} \cite[Theorem~5.4]{Hedlund1969}: each output symbol occurs exactly $4^{2R}$ times in the table. Balance is added to the CNF as four cardinality constraints, the balanced commuting maps are enumerated, and each is submitted to the pair-graph test. Balance is necessary and the pair-graph test is exact, so this yields $\Cent^{\times}_R(F)$ exactly.

\subsection{Linear rules}
When $f$ is linear over $\F_2$, the constraints \eqref{eq:constraints} are homogeneous linear equations in the table bits and $\Cent_R(F)$ is an $\F_2$-subspace of table space. Its dimension is computed by Gaussian elimination: the equalities $X_{e_L}=Y_{e_2}$ are merged by union--find, which reduces the $2\cdot4^{2R+1}$ bits to $4^{2R}$ classes, and the remaining equations are reduced by bit-packed elimination. For $R\le2$ the result was cross-checked against the SAT enumeration of all $2^{\dim}$ solutions.

\section{Rule 54: walls}\label{sec:54}

Rule $54$ is $f(a,b,c)=b\oplus(a\vee c)$; its values with $b=1$ are $f(0,1,0)=1$ and $f(a,1,c)=0$ whenever $a\vee c=1$. We need only
\begin{equation}\label{eq:wall}
f(\,\cdot\,,1,1)=f(1,1,\,\cdot\,)=0,\qquad f(1,0,\,\cdot\,)=f(\,\cdot\,,0,1)=1 .
\end{equation}

\subsection{Walls and segments}
Call a pair of adjacent cells $(c_i,c_{i+1})=(3,3)$ a \emph{wall}; it is the wall (iv) of Takesue \cite[Fig.~2, Table~III]{Takesue1989}, whose list for this wall contains $54$R. A \emph{segment} of a configuration $c$ is a maximal interval $[i,j]$ such that $c_{i-2}=c_{i-1}=3$, $c_{j+1}=c_{j+2}=3$, $c_i\ne3$, $c_j\ne3$, and $c$ contains no wall inside $[i,j]$; its \emph{content} is the word $c_i\cdots c_j\in\A^{L}$, $L=j-i+1$, and $S_L\subset\A^L$ denotes the set of words that can occur as contents, namely those with first and last letter different from $3$ and without $33$ as a factor. Note that a segment may contain isolated $3$'s.

\begin{lemma}\label{lem:walls}
Let $F=F_{54}$.
\begin{enumerate}
\item[(i)] $F(c)$ has a wall at $(i,i+1)$ if and only if $c$ does. The same holds for $F^{-1}$.
\item[(ii)] If $[i,j]$ is a segment of $c$ then $[i,j]$ is a segment of $F(c)$, and the content of $F(c)$ on $[i,j]$ depends only on the content of $c$ on $[i,j]$. Writing $F_L\colon S_L\to S_L$ for the induced map, $F_L$ is a permutation.
\item[(iii)] $F_1$ is the $3$-cycle $\pi\colon0\mapsto1\mapsto2\mapsto0$ on $S_1=\{0,1,2\}$.
\end{enumerate}
\end{lemma}

\begin{proof}
(i) If $c_i=c_{i+1}=3$ then $y_i=y_{i+1}=1$ and $x_i=x_{i+1}=1$, so by \eqref{eq:F} and \eqref{eq:wall}, $F(c)_i=(1,1\oplus f(y_{i-1},1,1))=3$ and $F(c)_{i+1}=(1,1\oplus f(1,1,y_{i+2}))=3$. Conversely, if $F(c)_i=F(c)_{i+1}=3$ then $y_i=y_{i+1}=1$, hence $f(y_{i-1},1,1)=f(1,1,y_{i+2})=0$, and the second coordinates force $x_i=x_{i+1}=1$. The statement for $F^{-1}$ follows by applying the forward statement to $F^{-1}(c)$.

(ii) By (i) the walls of $F(c)$ are those of $c$, so the maximal wall-free intervals between walls are the same, and the condition $c_i\ne3$, $c_j\ne3$ holds for $F(c)$ as well, since otherwise $F(c)$ would have a wall at $(i-1,i)$ or $(j,j+1)$ that $c$ lacks. The new value at a position $k\in[i,j]$ depends on $c_{k-1},c_k,c_{k+1}$, and the only positions outside $[i,j]$ involved are $i-1$ and $j+1$, where $c$ equals $3$ regardless of the configuration; so the content of $F(c)$ on $[i,j]$ is a function of the content of $c$ on $[i,j]$. This function maps $S_L$ into $S_L$ by the first part, and it is injective since $F^{-1}$ induces its inverse by the same argument.

(iii) For a segment $3\,s\,3$ with $s=(x,y)\ne3$, \eqref{eq:F} gives $F(c)$ at that position $=(y,x\oplus f(1,y,1))=(y,x\oplus y\oplus1)$, which sends $0\mapsto1\mapsto2\mapsto0$.
\end{proof}

\subsection{Relabelling automorphisms}

\begin{definition}
For $L\ge1$ and a permutation $\varphi$ of $S_L$, let $\Phi_L(\varphi)\colon\A^{\Z}\to\A^{\Z}$ be the map that replaces the content $s$ of every segment of length exactly $L$ by $\varphi(s)$ and leaves every other cell unchanged.
\end{definition}

Whether the cell at position $k$ belongs to a segment of length $L$, and if so which content and which position in it, is determined by $c_{k-L-1},\dots,c_{k+L+1}$: one needs at most $L-1$ cells on one side to reach the end of the segment and two more to see the wall. So $\Phi_L(\varphi)$ is a block map of radius at most $L+1$.

\begin{proof}[Proof of Theorem~\ref{thm:54}]
Let $\varphi\in C_L$, the centralizer of $F_L$ in $\Sym(S_L)$, and $g=\Phi_L(\varphi)$. Since $\varphi(s)\in S_L$ has no wall inside and does not begin or end with $3$, the walls of $g(c)$ are the walls of $c$, so $g(c)$ has the same segments as $c$, and $\Phi_L(\varphi^{-1})$ inverts $g$. Hence $g$ is an automorphism of $\A^{\Z}$ and $\Phi_L$ is an injective homomorphism $\Sym(S_L)\to\Aut(\A^{\Z})$; injectivity is seen on a configuration consisting of one wall and one segment of length $L$ repeated periodically.

To see that $g$ commutes with $F$, compare $F(g(c))$ and $g(F(c))$ cell by cell. Both configurations have the walls of $c$. At a position inside a segment of length $L$ with content $s$, the first is $F_L(\varphi(s))$ and the second is $\varphi(F_L(s))$ at the corresponding position, by Lemma~\ref{lem:walls}(ii) applied to $g(c)$ and to $c$ respectively, and these agree because $\varphi$ commutes with $F_L$. At a position $k$ inside a segment of another length, or inside a wall, or not in any segment, $g$ fixes $c_k$ and $F(c)_k$; moreover $F(c)_k$ and $F(g(c))_k$ agree, because the only cells where $g(c)$ differs from $c$ are inside segments of length $L$, and such a cell is adjacent to $k$ only if $k$ is a wall cell, whose image under $F$ is $3$ regardless of its neighbours by \eqref{eq:wall}.

For $L\ne L'$ the maps in $\Phi_L(C_L)$ and $\Phi_{L'}(C_{L'})$ act on disjoint sets of cells and commute, so the $\Phi_L$ assemble to a homomorphism $\Phi$ on the restricted direct sum $\bigoplus_L C_L$. It is injective: suppose $\Phi(\varphi_{L_1},\dots,\varphi_{L_m})=\mathrm{id}$ with $\varphi_{L_i}\in C_{L_i}$. Fix $i$ and $s\in S_{L_i}$, and let $c$ be the periodic configuration $(3\,3\,s)^{\infty}$. Every segment of $c$ has length $L_i$, so the components $\varphi_{L_j}$ with $j\ne i$ act on $c$ as the identity, and $\Phi(\varphi_{L_1},\dots,\varphi_{L_m})(c)=\Phi_{L_i}(\varphi_{L_i})(c)=(3\,3\,\varphi_{L_i}(s))^{\infty}$. This equals $c$ only if $\varphi_{L_i}(s)=s$; as $s$ was arbitrary, $\varphi_{L_i}=\mathrm{id}$. Every $\Phi_L(\varphi)$ commutes with $\sh$, being a block map.

Finally, $F_L\in C_L$, and $F_L\ne\mathrm{id}$ because the content $0^L$ is not fixed: by \eqref{eq:wall} its first cell becomes $(0,0\oplus f(1,0,0))=(0,1)$. So $o_L\ge2$ and $\Phi(\bigoplus_L\langle F_L\rangle)\cong\bigoplus_L\Z/o_L$ is an infinite direct sum of nontrivial finite cyclic groups; it is not finitely generated, since any finite set of elements lies in finitely many summands and generates a finite group. Every subgroup of a virtually $\Z^d$ group is finitely generated, since a subgroup $H$ meets the finite-index copy of $\Z^d$ in a finitely generated group of finite index in $H$. Hence $\Aut(X_F)$ is not virtually $\Z^d$ for any $d$. For $L=2$, both cells of a segment are adjacent to a wall, and by \eqref{eq:wall} each evolves as $(x,y)\mapsto(y,x\oplus y\oplus1)$ independently of the other, so $F_2=\pi\times\pi$ on $S_2=\{0,1,2\}^2$, a product of three $3$-cycles, and $C_2\cong\Z/3\wr S_3\cong(\Z/3)^3\rtimes S_3$, which is nonabelian. So $\Aut(X_F)$ is not abelian, and in particular is not $\langle\sh,F\rangle$. Note that $\bigoplus_L\Z/o_L$ is abelian and $\bigoplus_LC_L$ may or may not be virtually abelian; we do not know whether $\Aut(X_F)$ is virtually abelian.
\end{proof}

The element $D=\Phi_1(F_1)=\Phi_1(\pi)$ is the radius-$2$ block map that sends each window $3\,3\,s\,3\,3$ with $s\ne3$ to $\pi(s)$ and acts as the identity on every other window. Its table differs from that of the identity on exactly $3$ of the $1024$ entries. The orders $o_L$ of $F_L$ for $L=1,\dots,7$ are $3,3,6,90,2520,102702600$ and $4354384834800$; the content $0^L$ has period $3,3,6,9,18,21,18,21,30,33$ for $L=1,\dots,10$. More automorphisms come from the rest of $C_L$: for $L=2$, $S_2$ has $9$ elements on which $F_2$ acts as three $3$-cycles, so $C_2\cong(\Z/3)^3\rtimes S_3$ has order $162$, and $\Phi_2(C_2)$ consists of $162$ automorphisms of radius $3$, among them the elements of order $2$ that swap two of the three cycles.

\subsection{The census at radius 2}

\begin{theorem}\label{thm:54census}
$\Cent^{\times}_2(F_{54})$ has exactly $23$ elements: the $13$ maps $\sh^aF^b$ with $|a|+|b|\le2$ and the $10$ maps $F^bD^k$ with $|b|\le2$ and $k\in\{1,2\}$. At radius $1$, $\Cent_1(F_{54})$ has $13$ elements: the $5$ maps $\sh^aF^b$, the $2$ constants, and $6$ non-injective maps $g_{\pm,k}$, $k=0,1,2$, where $g_{+,k}(c)_i=\pi^k(c_{i-1})$ if $(c_i,c_{i+1})=(3,3)$ and $g_{+,k}(c)_i=3$ otherwise, and $g_{-,k}$ is the mirror image. These record, on each wall, the cell just outside it; they commute with $F$ because by \eqref{eq:wall} the cell adjacent to a wall evolves by $\pi$ whatever its other neighbour is. 
\end{theorem}

\begin{proof}
The radius-$1$ centralizer was enumerated completely; the $13$ tables are listed in the data accompanying the paper, and each can be checked by hand against \eqref{eq:constraints}. At radius $2$ the full centralizer exceeded the enumeration cap of $2\times10^4$ solutions (it contains, for instance, the compositions of the $g_{\pm,k}$ with the maps $\Phi_L$ and many other non-injective endomorphisms), so the balanced commuting maps were enumerated instead: there are exactly $23$, every one of them passed the pair-graph test, $13$ matched the tables of $\sh^aF^b$, and each of the remaining $10$ differs from the table of some $F^b$ on exactly the three windows $3\,3\,s\,3\,3$, $s\in\{0,1,2\}$, where it applies $\pi^{k}$ to the value that $F^b$ would produce. That is the table of $F^bD^k=D^kF^b$.
\end{proof}

The count is what the wall mechanism predicts. Maps $\Phi_L(\varphi)$ with $L\ge2$ have radius at least $3$, and $\sh^aD^k$ with $a\ne0$ has radius $3$, so the only elements of $\langle\sh,F\rangle\cdot\Phi(\bigoplus_LC_L)$ of radius at most $2$ are the $13+10$ listed. At radius $3$ the wall relabellings stop being the whole story.

\subsection{Every finite group}\label{sec:trains}

Call a maximal run of the symbol $3$ of length at least $2$ a \emph{wall of length} equal to the run length. By Lemma~\ref{lem:walls}, $F$ maps every wall onto itself and every maximal wall-free interval onto itself, so walls and their lengths are invariant. A \emph{car} is a cell $c_i\ne3$ with walls immediately on both sides, that is $c_{i-2}=c_{i-1}=c_{i+1}=c_{i+2}=3$. By \eqref{eq:wall} a car evolves by $\pi$, whatever lies beyond the walls. A \emph{train with $m$ cars} is a word
\[
W\;s_1\,3\,3\,s_2\,3\,3\cdots3\,3\,s_m\;W'
\]
in which every $s_j\ne3$ is a car, consecutive cars are separated by walls of length exactly $2$, and the outer walls $W,W'$ have length at least $3$. Trains are disjoint, and their positions, lengths and number of cars are invariant under $F$.

\begin{proof}[Proof of Theorem~\ref{thm:finite}]
For $\sigma\in S_m$ let $\Psi_m(\sigma)$ replace the cars $(s_1,\dots,s_m)$ of every train with exactly $m$ cars by $(s_{\sigma^{-1}(1)},\dots,s_{\sigma^{-1}(m)})$ and leave every other cell unchanged. Whether a cell is the $j$-th car of an $m$-car train is decided by the cells within distance $3m+3$, so $\Psi_m(\sigma)$ is a block map of radius at most $3m+3$. It maps trains to trains with the same cars in other places, so it preserves the set of trains, $\Psi_m(\sigma)\Psi_m(\rho)=\Psi_m(\sigma\rho)$, and $\Psi_m(\sigma^{-1})$ inverts $\Psi_m(\sigma)$. It commutes with $F$: both $F\Psi_m(\sigma)(c)$ and $\Psi_m(\sigma)F(c)$ carry the walls and non-train cells of $F(c)$, because $\Psi_m(\sigma)$ changes only car cells and a car's neighbours are wall cells whose images are $3$ regardless, and in each $m$-car train both carry $(\pi(s_{\sigma^{-1}(1)}),\dots,\pi(s_{\sigma^{-1}(m)}))$, since $F$ applies the same $\pi$ to every car. The homomorphism is injective: on the periodic configuration whose period is one $m$-car train with a single car equal to $1$ in position $j$ and all other cars $0$, followed by a wall of length $3$, the map $\Psi_m(\sigma)$ moves the $1$ to position $\sigma(j)$, so $\sigma$ is determined by its action. Maps for different $m$ act on disjoint sets of cells and commute. Finally every finite group $G$ embeds in $S_{|G|}$ by Cayley's theorem.
\end{proof}

Since $\bigoplus_mS_m$ is not finitely generated and contains nonabelian finite groups, Theorem~\ref{thm:finite} also implies the qualitative conclusions of Theorem~\ref{thm:54}; the relabellings $\Phi$ remain of interest because they have small radius and account exactly for the automorphisms of radius $2$.


\subsection{Trading wall cells for segment cells}\label{sec:tau}

\begin{proof}[Proof of Theorem~\ref{thm:tau}]
For a configuration $c$ let $P_1(c)$ be the set of positions $i$ with $c_{i-3}=c_{i-2}=c_{i-1}=3$, $c_i\ne3$, $c_{i+1}=c_{i+2}=3$ (a singleton whose left wall has length at least $3$), and $P_2(c)$ the set of positions $i$ with $c_{i-3}=c_{i-2}=3$, $c_{i-1}=\pi(c_i)$, $c_i\ne3$, $c_{i+1}=c_{i+2}=3$ (a two-cell segment of the form $(\pi(s),s)$). By definition $\tau(c)_{i-1}=\pi(c_i)$ for $i\in P_1(c)$, $\tau(c)_{i-1}=3$ for $i\in P_2(c)$, and $\tau(c)_j=c_j$ at every other position $j$. Whether $j$ is modified, and how, depends on $c_{j-2},\dots,c_{j+3}$, so $\tau$ is a block map of radius $3$. The sets $P_1(c)$ and $P_2(c)$ are disjoint, since $c_{i-1}=3$ in the first and $c_{i-1}\ne3$ in the second, and distinct $i$ give distinct modified positions $i-1$.

\emph{Involution.} Modifying position $i-1$ for $i\in P_1(c)\cup P_2(c)$ changes no other membership. Indeed, a position $j\ne i$ in $P_1$ or $P_2$ reads the cells $j-3,\dots,j+2$, and $i-1$ lies among them only for $j\in\{i-3,i-2,i,i+1,i+2\}$; but $i+1$ and $i+2$ are wall cells in both patterns, $i-2$ is a wall cell, and $i-3$ is a wall cell in $P_1$ and in $P_2$, while membership of $j$ in $P_1\cup P_2$ requires $c_j\ne3$. So the modifications at different positions are independent, $P_1(\tau(c))=P_2(c)$ and $P_2(\tau(c))=P_1(c)$, and the two modifications undo each other. Hence $\tau^2=\mathrm{id}$, and $\tau$ is a bijection, so $\tau\in\Aut(\A^{\Z})$.

\emph{Commutation.} By Lemma~\ref{lem:walls}, $F$ preserves walls. A cell adjacent to a wall evolves by $(x,y)\mapsto(y,x\oplus y\oplus1)$, which is $\pi$ on $\{0,1,2\}$, whatever its other neighbour is, by \eqref{eq:wall}. For $i\in P_1(c)$ the cell $c_i$ is adjacent to a wall, so $F(c)_i=\pi(c_i)\ne3$ and $i\in P_1(F(c))$. For $i\in P_2(c)$ both cells $c_{i-1}=\pi(s)$ and $c_i=s$ are adjacent to walls, so $F(c)$ carries $(\pi^2(s),\pi(s))=(\pi(\pi(s)),\pi(s))$ there and $i\in P_2(F(c))$. Thus $P_1(F(c))=P_1(c)$ and $P_2(F(c))=P_2(c)$. We compare $\tau(F(c))$ and $F(\tau(c))$ cell by cell. At $i-1$ with $i\in P_1(c)$: the first is $\pi(F(c)_i)=\pi^2(c_i)$; in $\tau(c)$ the cell $i-1$ holds $\pi(c_i)$ and has the wall $c_{i-3}c_{i-2}=3\,3$ on its left, so it evolves by $\pi$ and the second is $\pi^2(c_i)$ as well. At $i$ with $i\in P_1(c)$: both equal $\pi(c_i)$, since $c_i$ is adjacent to the wall on its right in $c$ and in $\tau(c)$. At $i-1$ with $i\in P_2(c)$: the first is $3$; in $\tau(c)$ the cell $i-1$ is $3$ with $3\,3$ on its left, hence part of a wall, and stays $3$. At $i$ with $i\in P_2(c)$: both equal $\pi(c_i)$. At any other position $j$, $\tau$ changes neither $c_j$ nor $F(c)_j$, and $F(\tau(c))_j=F(c)_j$ because the only cells where $\tau(c)$ differs from $c$ are the positions $i-1$, whose neighbours are $i$ (already treated) and $i-2$, a wall cell, whose image under $F$ is $3$ regardless of its right neighbour by \eqref{eq:wall}. So $\tau F=F\tau$.

\emph{The rest.} $\tau$ sends $\cdots3\,3\,3\,0\,3\,3\cdots$ to $\cdots3\,3\,1\,0\,3\,3\cdots$, so it does not preserve walls, and therefore lies outside the subgroup generated by $\sh$, $F$ and the wall relabellings, all of which preserve walls. For $c=\cdots3\,3\,3\,0\,3\,3\,3\cdots$ one has $\tau D(c)=\tau(\cdots3\,3\,3\,1\,3\,3\,3\cdots)=\cdots3\,3\,2\,1\,3\,3\,3\cdots$, whereas $\tau(c)=\cdots3\,3\,1\,0\,3\,3\,3\cdots$ contains no singleton, so $D\tau(c)=\tau(c)\ne\tau D(c)$.
\end{proof}


\subsection{Radius 3}\label{sec:54r3}

The enumeration of $\Cent^{\times}_3(F_{54})$ needs more than the balance filter, because the balanced commuting maps at radius $3$ number more than $2\times10^4$ and almost all of them are not injective. The computation adds to the CNF the injectivity of $g$ on all spatially periodic configurations of period at most $5$, and then learns further necessary conditions lazily: whenever the solver returns a non-injective map, a local witness is extracted, either a diamond (two words of length $13$ with the same prefix and suffix of length $6$ and equal images of length $7$, which by Hedlund's theorem no surjective map admits) or a pair of periodic configurations of period at most $10$ with equal images, and the clause forbidding that witness is added and shared among all worker processes. Every surviving map is decided by the pair-graph test. This enumeration is not complete, and we give no count of $\Cent^{\times}_3(F_{54})$ (Appendix~\ref{app:census}). Every automorphism of radius $3$ that it has produced differs from some $\sh^aF^b$, $|a|+|b|\le3$, only on windows of length $7$ that contain a wall, and most of them, like $\tau$, do not preserve walls. The wall relabellings of Theorem~\ref{thm:54} and the wall-trading involutions of Theorem~\ref{thm:tau} are two instances of one principle: between two walls, $F$ acts as a finite permutation, and a local, reversible re-encoding of the walls and segments that intertwines these permutations commutes with $F$.

\begin{remark}
Nothing in Lemma~\ref{lem:walls}(i),(ii) uses more of rule $54$ than the three values $f(0,1,1)=f(1,1,0)=f(1,1,1)=0$, so $33$ is a wall for the $32$ elementary rules with these values, among them rules $0$--$7$, $16$--$23$, $32$--$39$ and $48$--$55$, and each such second-order automaton admits the relabelling automorphisms $\Phi_L(\varphi)$ for $\varphi$ in the centralizer of its own $F_L$. Rule $110$ has $f(0,1,1)=1$, rule $30$ has $f(0,1,1)=1$, and rule $90$ has $f(0,1,1)=1$; none of them has $33$ as a wall. For rules $110$ and $30$ the automorphism groups are $\langle\sh,F\rangle\cong\Z^2$ \cite{PaperA,PaperB}.
\end{remark}


\section{Walls in general}\label{sec:walls}

In this section $F$ is a reversible cellular automaton of radius $1$ on a finite alphabet $\A$. A nonempty word $W$ is \emph{fixed in every context} if $F(c)$ contains $W$ at every position at which $c$ contains $W$.

\begin{lemma}\label{lem:exact}
Under the hypothesis of Theorem~\ref{thm:walls}, $F$ neither creates nor destroys occurrences of $W$: $F(c)$ contains $W$ at position $i$ if and only if $c$ does.
\end{lemma}

\begin{proof}
Fix $n$ and let $P_n$ be the finite set of configurations of spatial period $n$. $F$ maps $P_n$ bijectively to itself, being injective and commuting with $\sh^n$. Write $N(c)$ for the number of positions $i\in\{0,\dots,n-1\}$ at which $c$ contains $W$. By hypothesis $N(F(c))\ge N(c)$ for every $c\in P_n$, and $\sum_{c\in P_n}N(F(c))=\sum_{c\in P_n}N(c)$ because $F$ permutes $P_n$. Hence $N(F(c))=N(c)$, so on $P_n$ the occurrences of $W$ in $F(c)$ are exactly those of $c$. If some configuration $c$ had an occurrence of $W$ in $F(c)$ at a position $i$ where $c$ has none, the same would hold for the periodic configuration obtained by repeating a long enough window of $c$ around $i$, since both properties depend only on $c_{i-1},\dots,c_{i+|W|}$.
\end{proof}

\begin{proof}[Proof of Theorem~\ref{thm:walls}]
Let $k=|W|$. A \emph{car} of $c$ is a position $i$ with $c$ containing $W$ at $i-k$ and at $i+1$ and $c_i\in C$; by the definition of $C$ no occurrence of $W$ covers $i$. By Lemma~\ref{lem:exact} the occurrences of $W$ in $F(c)$ are those of $c$, and the cells of each occurrence keep their values, so the neighbours of a car keep their values and the car evolves by a map $\varphi\colon C\to\A$, $\varphi(s)=F(W\,s\,W)$ at the middle position, independent of the rest of $c$. The image $\varphi(s)$ lies in $C$, because an occurrence of $W$ in $W\,\varphi(s)\,W$ other than at its ends would be an occurrence created by $F$. And $\varphi$ is injective, because $F$ is injective on the periodic configurations $(W\,s)^{\infty}$, which differ only at their cars. So $\varphi$ is a permutation of $C$, and $F$ maps cars to cars. Conversely every car of $F(c)$ is a car of $c$: its position is flanked by the same occurrences of $W$, and if the symbol of $c$ there were outside $C$, then $c$ would contain an occurrence of $W$ covering it, which $F$ would preserve. So cars are invariant, and $F$ acts on every car by the same $\varphi$; this uses that the neighbours of a car lie in occurrences of $W$, whose images under $F$ do not depend on the context, which is where the hypothesis ``fixed in every context'' enters. A \emph{train with $m$ cars} is a chain of cars $i_1<\dots<i_m$ with $i_{j+1}=i_j+k+1$, such that $c$ contains $W$ at $i_1-2k$ and at $i_m+k+1$ (two copies of $W$ with no car between them at each end), and that is maximal with these properties. Trains are determined by the occurrences of $W$ and the cars, so they are invariant under $F$. As in the proof of Theorem~\ref{thm:finite}, permuting the cars of every $m$-car train by $\sigma\in S_m$ defines a block map $\Psi_m(\sigma)$ of radius at most $m(k+1)+2k+1$, which is an automorphism commuting with $F$ because $F$ applies the same $\varphi$ to every car. With two distinct symbols $a,b\in C$, the periodic configuration whose period is one $m$-car train with cars $(b,\dots,b,a,b,\dots,b)$ shows that $\sigma$ is recovered from $\Psi_m(\sigma)$. The map $\Psi_m(\sigma)$ preserves the occurrences of $W$, since any occurrence covering a car lies inside $W\,s\,W$ and the new symbol is again in $C$; hence it preserves cars and trains. Maps for different $m$ act on disjoint cells and commute, and the resulting homomorphism of $\bigoplus_mS_m$ is injective, as one sees by evaluating on periodic configurations $(W\,W\,s_1\,W\,s_2\cdots W\,s_m)^{\infty}$ whose only trains have $m$ cars.
\end{proof}

The train definition here, bounded by two adjacent copies of $W$, differs from the one used for rule $54$ in Theorem~\ref{thm:finite}, where outer walls of length at least $3$ suffice; both definitions work, and the general one is what is needed for words that overlap themselves.

\begin{table}[t]
\caption{Second-order elementary automata with a word $W$ fixed in every context, one rule per conjugacy class, with the shortest such $W$ (cells written $0$--$3$). The words are the fixed walls (i)--(viii) of Takesue \cite[Fig.~2, Table~III]{Takesue1989}, and the $31$ classes are those of the rules listed there for these walls. Classes are taken under reflection and the cellwise complement $\kappa(x,y)=(1-x,1-y)$, which conjugates $F_f$ to $F_{f'}$ with $f'(a,b,c)=f(1-a,1-b,1-c)$, as in \cite[Table~I]{Takesue1989}; Wolfram's complement $1\oplus f(1-a,1-b,1-c)$ is not a conjugacy of second-order automata \cite[p.~377]{Takesue1989}, and for instance $110$R and $137$R have different numbers of periodic points. The property is constant on each class. For each listed rule $|C|\ge2$. Of the $57$ classes in the last row, only $133=\{133,161\}$ has a fixed word of greater length ($W=0330$, $k=4$, the wall (ix) of \cite{Takesue1989}; Section~\ref{cl:sec:133}).}
\label{tab:walls}
\centering\small
\begin{tabular}{l}
\toprule
$W=0$: 0, 4, 12 \qquad $W=3$: 1, 2, 3, 18, 19, 33, 35, 50, 51\\
$W=30$: 9, 10, 11, 131, 139 \qquad $W=03$: 129 \qquad $W=030$: 73 \qquad $W=303$: 147 \qquad $W=00$: 36\\
$W=33$: 5, 6, 7, 22, 23, 37, 38, 39, 54, 55\\
\midrule
no fixed word of length $\le3$ ($57$ classes): 13--15, 24--31, 41, 43, 45--47, 57--63, 75, 77, 79, 90, 91,\\
\quad 94, 95, 105, 107, 109, 111, 123, 126, 127, 133, 135, 141, 143, 151, 153, 155, 157, 159, 165, 167,\\
\quad 175, 179, 183, 187, 189, 191, 219, 223, 255\\
\bottomrule
\end{tabular}
\end{table}

\begin{table}[t]
\caption{Spread of a single changed cell after $300$ steps, mean over $40$ random configurations of length $801$, and number of trials in which the difference disappeared.}
\label{tab:spread}
\centering
\begin{tabular}{lrrr}
\toprule
rule & leftward spread & rightward spread & healed\\
\midrule
54R & 15 & 21 & 0\\
110R & 195 (min 167) & 41 (min 22) & 0\\
30R & 98 (min 72) & 300 (min 299) & 0\\
\bottomrule
\end{tabular}
\end{table}


\begin{proof}[Proof of Corollary~\ref{cor:48}]
For each rule of Table~\ref{tab:walls} the listed word $W$ is fixed in every context, and the set $C$ of Theorem~\ref{thm:walls} has at least two elements; both are finite checks of the local rule over the $4^2$ contexts of $W$ and the four symbols $s$ in $W\,s\,W$. Theorem~\ref{thm:walls} applies, and by the conjugacy invariance of Section~\ref{sec:prelim} it applies to every rule of the $31$ classes, $81$ rules in all. The class of $133$ is Proposition~\ref{cl:prop:133}(1), and the last statement is Proposition~\ref{cl:prop:133}(3).
\end{proof}

The walls also confine the spreading of differences: in random configurations of $54$R a single changed cell produces a difference that spreads only a few cells, while in $110$R and $30$R it spreads in both directions at linear speed (Table~\ref{tab:spread}).

\subsection{The class of 133: a wall of length 4}\label{cl:sec:133}

Table~\ref{tab:walls} lists the fixed words of length at most $3$. One more class has a fixed word, of length $4$. It is the wall (ix) of Takesue, who notes that it is not detected with fewer than four cells \cite[Table~III and p.~391]{Takesue1989}. In the next proposition the fixed word of part~(1) is his; the automorphisms of parts~(1) and~(2) and the completeness statement~(3) are what is added. The examination in \cite[Section~4]{Takesue1989} covers the words of at most $6$ cells, and it is accompanied by the conjecture, stated there without proof, that $4$ cells suffice to detect a local conservation law of a second-order elementary automaton. Part~(3) concerns fixed words only, and for them it gives a proof for every length.

\begin{proposition}[the fixed word of $133$R: {\cite[p.~391]{Takesue1989}}]\label{cl:prop:133}
Let $f\in\{133,161\}$ and $F=F_f$.
\begin{enumerate}
\item The word $W=0330$ is fixed by $F_{133}$ in every context, and $W=3003$ is fixed by $F_{161}$ in every context. In both cases the set $C$ of Theorem~\ref{thm:walls} is all of $\A$, so $\bigoplus_mS_m$ embeds in $\Aut(X_F)$.
\item For $f=133$, let $D$ be the block map of radius $4$ that replaces the middle cell $s$ of every occurrence of $W\,s\,W$ by $F(W\,s\,W)$ at that cell and fixes every other cell. Then $D\in\Aut(X_F)$, $D$ acts on each such cell as the $4$-cycle $0\mapsto1\mapsto3\mapsto2\mapsto0$, $D^4=\mathrm{id}\ne D^2$, and $D\notin\langle\sh,F\rangle$.
\item Every word fixed by some $F_{f'}$ in every context lies in $\{0,3\}^k$. Among the $57$ classes with no such word of length at most $3$, the class of $133$ is the only one with such a word of any length.
\end{enumerate}
\end{proposition}

\begin{proof}
(1) A cell of a word fixed in every context keeps its value, so by \eqref{eq:F} its two bits agree and $f(y_{j-1},y_j,y_{j+1})=0$ at each of its positions $j$, the outer neighbour bits being arbitrary. For $133$, $f(a,b,c)=1$ exactly on $000$, $010$ and $111$. The $y$-row of $W=0330$ is $0110$, and the triples $(u,0,1)$, $(0,1,1)$, $(1,1,0)$ and $(1,0,v)$, $u,v\in\{0,1\}$, all avoid $000$, $010$ and $111$; so $W$ is fixed in all $16$ contexts. Since $161=\kappa(133)$ and $\kappa$ maps $0330$ to $3003$, the statement for $161$ follows. Finally $W\,s\,W$ contains $W$ only at its two ends for each $s\in\A$, so $C=\A$. Theorem~\ref{thm:walls} applies.

(2) The neighbours of the middle cell $s$ of $W\,s\,W$ have $y$-bit $0$, so it evolves by $(x,y)\mapsto(y,x\oplus f(0,y,0))=(y,x\oplus1)$, the $4$-cycle $0\to1\to3\to2\to0$. So $D$ is the map that applies the car permutation $\varphi$ of the proof of Theorem~\ref{thm:walls} to every car. As there, every length-$4$ window covering a car lies in $W\,s\,W$, and $W\,\varphi(s)\,W$ contains $W$ only at its ends, so $D$ preserves the occurrences of $W$ and the cars, $D^{-1}=D^3$, and $D$ commutes with $F$ because $F$ applies $\varphi$ to every car and the neighbours of a car lie in occurrences of $W$, whose images do not depend on the context. Since $\varphi$ is a $4$-cycle, $D^4=\mathrm{id}\ne D^2$. Suppose $D=\sh^aF^b$. The configuration $0^{\Z}$ contains no $W$, so $D$ fixes it; $F$ acts on constant configurations as the $4$-cycle $0\to1\to3\to2\to0$, and $\sh$ fixes them, so $4\mid b$. On $c=(0330\,0)^{\infty}$, every fifth cell is a car, so $D(c)=F(c)$, and therefore $\sh^aF^{b-1}c=c$. But $F^{b-1}c=(0330\,s)^{\infty}$ with $s=\varphi^{b-1}(0)\ne0$, as $b-1\not\equiv0\pmod4$, and no shift of it equals $c$. So $D\notin\langle\sh,F\rangle$.

(3) The first statement is the first sentence of the proof of (1). With $w$ the $y$-row, a fixed word of length $k\ge2$ is a word $w\in\{0,1\}^k$ with $f(u,w_0,w_1)=0$ and $f(w_{k-2},w_{k-1},v)=0$ for $u,v\in\{0,1\}$ and $f=0$ on its interior triples: a walk through $k-1$ states in a graph on the $4$ pairs $(w_j,w_{j+1})$, from a set of allowed initial pairs to a set of allowed final pairs. If such a walk exists, a shortest one visits each state at most once, so some fixed word has length at most $5$. A search of all lengths $k\le40$ for all $256$ rules therefore decides the question for every length; it finds fixed words for exactly the $81$ rules of Table~\ref{tab:walls} at length at most $3$, and, among the other rules, only for $133$ and $161$, at the lengths $4$, $7$, $8$ and every $k$ with $10\le k\le40$. A brute-force search over all of $\A^k$, $k\le4$, agrees with the reduction to $\{0,3\}^k$. The search is a finite computation, described in Appendix~\ref{app:proofs}.
\end{proof}



\section{Walls with hidden bits}\label{sec:hidden}

The walls of Sections~\ref{sec:54} and~\ref{sec:walls} are words that $F$ fixes in every context, and the automorphisms of Theorems~\ref{thm:finite} and~\ref{thm:walls} permute the cars that sit between them. A wall need not be fixed for this to work. Walls that are periodic in time and not fixed are in the list of Takesue \cite[Section~4]{Takesue1989}: for $27$R and $59$R the runs of cells in equal phase given by the local conservation law \cite[(4.1)--(4.3)]{Takesue1989}, of period $3$; for $179$R the pairs of cells (xii) and (xiii) of \cite[Fig.~2, Table~III]{Takesue1989}, of period $3$; for $219$R and $223$R the pairs (xv), and for $126$R and $127$R the pairs (xvi), of period $4$. The walls of this section are built from these, and we say in each case how. What is added here is the bit that such a wall can carry, the flips, and the embeddings that result. In this section the wall is a word that moves under $F$, but whose image never depends on the context, except perhaps at cells that do not matter, and which carries one bit that $F$ cannot see. Permuting these hidden bits along chains of walls gives every finite group again. The main example is rule $59$R, which has no word of any length fixed in every context (Lemma~\ref{hw:nofixed}) and whose walls are not fixed by any power of $F$ that we have checked; the general mechanism also covers walls that are fixed by a power of $F$, which we call time-periodic, and these occur for rules $127$R, $179$R and $223$R. A variant of the rule-$59$ walls with windows of length $6$ gives the same conclusion for rule $27$R (Theorem~\ref{hw:27}). Finally, for rules $219$R and $126$R the walls are runs of cells bound by a relation that $F$ preserves, and they rotate through the alphabet in every context; we call them rotating walls (Section~\ref{hw:sec:moving}). All seven classes discussed here, $27$, $59$, $126$, $127$, $179$, $219$ and $223$, lie in the last row of Table~\ref{tab:walls}: none of them has a word of length at most $3$ fixed in every context.

Throughout, $F$ is a reversible cellular automaton of radius $1$ on $\A^{\Z}$, and for a configuration $c$ and integers $i<j$ we write $c_{[i,j)}=c_i c_{i+1}\cdots c_{j-1}$.

\subsection{The general mechanism}\label{hw:sec:general}

Fix $k\ge3$, a set of words $W\subseteq\A^k$, a position $j$ with $1\le j\le k-2$, and an involution $\iota\colon W\to W$ such that $\iota(w)$ differs from $w$ exactly at position $j$, for every $w\in W$. For a configuration $c$ let
\[
\Occ(c)=\{i\in\Z:\ c_{[i,i+k)}\in W\}
\]
be its set of \emph{occurrences}; for $i\in\Occ(c)$ write $w_i(c)=c_{[i,i+k)}$, and let $\iota_ic$ be the configuration obtained from $c$ by replacing $w_i(c)$ by $\iota(w_i(c))$. Only the cell $i+j$, the \emph{flip cell} of the occurrence, changes. We consider the following hypotheses. Each is stated at position $0$; since $F$ commutes with $\sh$, it then holds at every position.
\begin{enumerate}
\item[(W1)] If $0\in\Occ(c)$ then $0\in\Occ(F(c))$, for every $c$.
\item[(H1)] If $0\in\Occ(c)$ then $F(\iota_0c)=\iota_0F(c)$.
\item[(H2)] If $0,i\in\Occ(c)$ with $i\ne0$, then $i\in\Occ(\iota_0c)$, and $\iota_0\iota_ic=\iota_i\iota_0c$.
\item[(H3)] There are $\beta\colon W\to\Z/2$ and $\varepsilon\in\Z/2$ such that $\beta\circ\iota=\beta+1$ and $\beta(w_0(F(c)))=\beta(w_0(c))+\varepsilon$ whenever $0\in\Occ(c)$.
\item[(H3$'$)] There are $\beta\colon W\to\Z/2$, a surjection $\kappa\colon W\to K$ onto a finite set, and maps $\theta\colon K\to K$, $\varepsilon\colon K\to\Z/2$ such that $\kappa\circ\iota=\kappa$, $\beta\circ\iota=\beta+1$, and whenever $0\in\Occ(c)$,
\[
\kappa(w_0(F(c)))=\theta(\kappa(w_0(c))),\qquad \beta(w_0(F(c)))=\beta(w_0(c))+\varepsilon(\kappa(w_0(c))).
\]
\item[(H4)] If $0,i\in\Occ(c)$ with $i\ne0$, then $\beta(w_i(\iota_0c))=\beta(w_i(c))$, and under (H3$'$) also $\kappa(w_i(\iota_0c))=\kappa(w_i(c))$.
\end{enumerate}
We call $\beta(w_i(c))$ the \emph{bit} and $\kappa(w_i(c))$ the \emph{class} of the occurrence $i$ of $c$. Hypothesis (H3) is the case $|K|=1$ of (H3$'$), and we work with (H3$'$) throughout. In (H2), $\iota_i\iota_0c$ is defined by the first half of (H2), and $\iota_0\iota_ic$ by the first half with $0$ and $i$ exchanged. The bit is ``hidden'' in the sense of (H1): flipping it commutes with $F$, so the rest of the configuration evolves as if it were not there.

\begin{lemma}\label{hw:exact}
Under \textup{(W1)}, $\Occ(F(c))=\Occ(c)$ and $\Occ(F^{-1}(c))=\Occ(c)$ for every configuration $c$.
\end{lemma}

\begin{proof}
This is the proof of Lemma~\ref{lem:exact}, with ``$c$ contains $W$ at $i$'' read as $i\in\Occ(c)$. On the finite set of configurations of spatial period $n$, which $F$ permutes, the number $N(c)$ of occurrences in a period satisfies $N(F(c))\ge N(c)$ by (W1), and summing over the set gives equality, so $\Occ(F(c))=\Occ(c)$ there. For general $c$, whether $i\in\Occ(c)$ and whether $i\in\Occ(F(c))$ depend only on $c_{[i-1,i+k+1)}$, so a counterexample would yield a periodic one.
\end{proof}

For $c$ and $I\subseteq\Occ(c)$ define $\iota_Ic$ cellwise by
\[
(\iota_Ic)_{i+j}=\iota(w_i(c))_j\quad(i\in I),\qquad (\iota_Ic)_n=c_n\quad(n\notin I+j).
\]
This is well defined because distinct occurrences have distinct flip cells. Note that $\iota_Ic$ is computed from the words of $c$, all at once.

\begin{lemma}\label{hw:flips}
Assume \textup{(W1)} and \textup{(H2)}, and let $I,J\subseteq\Occ(c)$.
\begin{enumerate}
\item[(a)] $\Occ(\iota_Ic)=\Occ(c)$; in particular flips neither destroy nor create occurrences.
\item[(b)] $\iota_J\iota_Ic=\iota_{I\triangle J}c$. In particular $\iota_I\iota_Ic=c$.
\item[(c)] If also \textup{(H3$'$)} and \textup{(H4)} hold, then for every $i\in\Occ(c)$ the class of $i$ in $\iota_Ic$ is its class in $c$, and the bit of $i$ in $\iota_Ic$ is its bit in $c$ plus $[i\in I]$.
\item[(d)] If also \textup{(H1)} holds, then $F(\iota_Ic)=\iota_IF(c)$.
\end{enumerate}
\end{lemma}

\begin{proof}
\emph{One flip.} Let $i\in\Occ(c)$. The word of $\iota_ic$ at $i$ is $\iota(w_i(c))$, since $\iota$ changes only position $j$; hence $\iota_i\iota_ic=c$. By (H2), every $i'\in\Occ(c)$ lies in $\Occ(\iota_ic)$. Conversely, if $i'\in\Occ(\iota_ic)$ with $i'\ne i$, then (H2) applied to the configuration $\iota_ic$, which has occurrences at $i$ and $i'$, shows that $i'\in\Occ(\iota_i\iota_ic)=\Occ(c)$. So $\Occ(\iota_ic)=\Occ(c)$: flips do not create occurrences, because $\iota$ is an involution.

\emph{Finitely many flips.} Let $I=\{i_1,\dots,i_r\}$ be finite. By the first step every intermediate configuration has occurrence set $\Occ(c)$, so the composition $\iota_{i_1}\circ\dots\circ\iota_{i_r}(c)$ is defined, and by (H2), applied to intermediate configurations, it does not depend on the order. It equals $\iota_Ic$: the cell $i_s+j$ is changed only by $\iota_{i_s}$, and ordering the composition so that $\iota_{i_s}$ acts first shows that this cell becomes $\iota(w_{i_s}(c))_j$; cells outside $I+j$ are never changed. This proves (a) for finite $I$, and (b) follows, since $\iota_J\iota_Ic$ is then a composition of single flips in which those at $I\cap J$ occur twice and cancel. For (c), a single flip at $i_s$ changes the word at $i_s$ to its $\iota$-partner, which has the same class and the other bit, and by (H4), applied to the intermediate configuration, it leaves the bit and class of every other occurrence unchanged. For (d), if $i\in\Occ(y)$ then $F(\iota_iy)=\iota_iF(y)$ by (H1), and $\Occ(F(y))=\Occ(y)$ by Lemma~\ref{hw:exact}; induction on $r$ gives $F(\iota_Ic)=\iota_IF(c)$.

\emph{Infinite sets.} For a finite set $U$ of cells, $\iota_Ic$ agrees on $U$ with $\iota_{I_U}c$, where $I_U=\{i\in I:\ i+j\in U\}$ is finite. Each of (a)--(d) is a statement about finitely many cells at a time (an occurrence reads $k$ cells, a flip reads one window, $F$ reads three cells), so each follows from the finite case applied to suitable finite subsets of $I$ and $J$.
\end{proof}

\begin{corollary}\label{hw:flipall}
Under \textup{(W1)}, \textup{(H1)} and \textup{(H2)}, the map $\tau_W(c)=\iota_{\Occ(c)}c$, which flips every occurrence, is an involutive automorphism of $\A^{\Z}$ of radius at most $\max(j,k-1-j)$ commuting with $F$.
\end{corollary}

\begin{proof}
Cell $n$ of $\tau_W(c)$ depends only on $c_{[n-j,n-j+k)}$. By Lemma~\ref{hw:flips}(a),(b), $\tau_W(\tau_W(c))=\iota_{\Occ(c)}\iota_{\Occ(c)}c=c$, and by Lemma~\ref{hw:flips}(d) and Lemma~\ref{hw:exact},
\[
F(\tau_W(c))=\iota_{\Occ(c)}F(c)=\iota_{\Occ(F(c))}F(c)=\tau_W(F(c)).
\]
\end{proof}

\begin{definition}\label{hw:train}
Fix a gap $d\ge1$ and assume (H3$'$). A \emph{train of length $m$} of $c$ is a set $\{i,i+d,\dots,i+(m-1)d\}\subseteq\Occ(c)$ of occurrences that all have the same class $x$, such that neither $i-d$ nor $i+md$ is an occurrence of class $x$. Its \emph{bit vector} is $(b_1,\dots,b_m)$, where $b_r$ is the bit of $i+(r-1)d$. For $\sigma\in S_m$ let $I_\sigma(c)$ be the set of occurrences $i+(r-1)d$, taken over all trains of length $m$, with $b_{\sigma^{-1}(r)}\ne b_r$, and define
\[
\Psi_m(\sigma)(c)=\iota_{I_\sigma(c)}c .
\]
\end{definition}

Every occurrence lies in at most one train: the occurrences of a given class $x$ whose positions are congruent modulo $d$ split into maximal arithmetic progressions of step $d$, and the trains are the finite ones among them. An occurrence that lies in an infinite progression, one-sided or two-sided, belongs to no train; nothing below needs these, since $\Psi_m$ touches only trains of length exactly $m$. By Lemma~\ref{hw:flips}(c), $\Psi_m(\sigma)$ replaces the bit vector $b$ of every train of length $m$ by $b\circ\sigma^{-1}=(b_{\sigma^{-1}(1)},\dots,b_{\sigma^{-1}(m)})$ and changes no other bit and no class.

\begin{lemma}\label{hw:trainlemma}
Assume \textup{(W1)}, \textup{(H1)}, \textup{(H2)}, \textup{(H3$'$)} and \textup{(H4)}, and let $d\ge1$. Then for every $m\ge1$:
\begin{enumerate}
\item[(i)] $\Psi_m(\sigma)$ is a block map of radius at most $md+\max(j,k-1-j)$;
\item[(ii)] $\Psi_m(\sigma)\Psi_m(\rho)=\Psi_m(\sigma\rho)$ and $\Psi_m(\mathrm{id})=\mathrm{id}$, so each $\Psi_m(\sigma)$ is an automorphism of $\A^{\Z}$;
\item[(iii)] $F\Psi_m(\sigma)=\Psi_m(\sigma)F$, so $\Psi_m(\sigma)\in\Aut(X_F)$ by Lemma~\ref{lem:spacetime};
\item[(iv)] $\Psi_m(\sigma)$ and $\Psi_{m'}(\rho)$ commute for $m\ne m'$, and $\Psi_{m'}(\rho)$ does not change the bits of the trains of length $m$.
\end{enumerate}
\end{lemma}

\begin{proof}
(i) Cell $n$ can change only if $n-j=:i$ is an occurrence. Whether $i$ is the $r$-th member of a train of length exactly $m$, and the bits of that train, are decided by the words at $i+td$, $-m\le t\le m$, and these lie in $c_{[n-j-md,\,n-j+md+k)}$. The construction is the same at every position, so $\Psi_m(\sigma)$ commutes with $\sh$.

We show that trains are invariant. By Lemma~\ref{hw:flips}(a),(c), $\iota_Ic$ has the same occurrences, with the same classes, as $c$; so it has the same trains. For $F$: the map $\theta$ is a permutation of $K$. Indeed, given $x\in K$ choose $c$ with $w_0(c)$ of class $x$; then $0\in\Occ(F^{-1}(c))$ by Lemma~\ref{hw:exact}, and $\theta$ maps the class of $0$ in $F^{-1}(c)$ to $x$, so $\theta$ is onto, hence bijective. Since $\Occ(F(c))=\Occ(c)$ and every class is replaced by its image under the injective map $\theta$, two occurrences have the same class in $F(c)$ if and only if they do in $c$. Hence $F(c)$ has the same trains as $c$, and by (H3$'$) the bit vector of a train of class $x$ changes from $b$ to $b+\varepsilon(x)(1,\dots,1)$. The same constant is added to every bit of a train, so the set of $r$ with $b_{\sigma^{-1}(r)}\ne b_r$ is unchanged: $I_\sigma(F(c))=I_\sigma(c)$.

(iii) By Lemma~\ref{hw:flips}(d), $F\Psi_m(\sigma)(c)=F(\iota_{I_\sigma(c)}c)=\iota_{I_\sigma(c)}F(c)=\iota_{I_\sigma(F(c))}F(c)=\Psi_m(\sigma)F(c)$.

(ii) Let $c'=\Psi_m(\rho)(c)=\iota_{I_\rho(c)}c$. It has the trains of $c$, and the bit vector $b$ of each train of length $m$ has become $b\circ\rho^{-1}$. Then $\Psi_m(\sigma)(c')=\iota_{I'}c'$, where $I'$ consists of the members at which $b\circ\rho^{-1}\circ\sigma^{-1}$ differs from $b\circ\rho^{-1}$. By Lemma~\ref{hw:flips}(b), $\Psi_m(\sigma)\Psi_m(\rho)(c)=\iota_{I_\rho(c)\triangle I'}c$, and $I_\rho(c)\triangle I'$ is the set of members at which $b\circ(\sigma\rho)^{-1}$ differs from $b$, which is $I_{\sigma\rho}(c)$. Clearly $I_{\mathrm{id}}(c)=\emptyset$. So $\Psi_m(\sigma^{-1})$ inverts $\Psi_m(\sigma)$, and an invertible block map is an automorphism of $\A^{\Z}$.

(iv) $\Psi_{m'}(\rho)$ flips only members of trains of length $m'$, so by Lemma~\ref{hw:flips}(c) it leaves the trains and the bits of the trains of length $m$ unchanged; hence $I_\sigma(\Psi_{m'}(\rho)(c))=I_\sigma(c)$, and symmetrically. By Lemma~\ref{hw:flips}(b) both composites equal $\iota_{I_\sigma(c)\cup I_\rho(c)}c$.
\end{proof}

\begin{theorem}\label{hw:general}
Let $F$ be a reversible cellular automaton of radius $1$ on $\A^{\Z}$, and let $W\subsetneq\A^k$, $j$ and $\iota$ be as above. Assume \textup{(W1)}, \textup{(H1)}, \textup{(H2)}, \textup{(H4)}, and either \textup{(H3)} or \textup{(H3$'$)}. Then for every $d\ge k$ and $m\ge1$ the maps $\Psi_m$ of Definition~\ref{hw:train} form an injective homomorphism $S_m\to\Aut(X_F)$ whose image consists of block maps of radius at most $md+\max(j,k-1-j)$; the images for different $m$ commute, and together they give an embedding of $\bigoplus_{m\ge1}S_m$ in $\Aut(X_F)$. In particular $\Aut(X_F)$ contains every finite group and is not virtually $\Z^n$ for any $n$.
\end{theorem}

\begin{proof}
Everything except injectivity is Lemma~\ref{hw:trainlemma}, and (iv) there shows that the $\Psi_m$ assemble to a homomorphism $\Psi$ of the restricted direct sum. Let $g=\Psi((\sigma_{m'})_{m'})$ with $\sigma_m\ne\mathrm{id}$ for some $m$, and choose $r$ with $\sigma_m(r)\ne r$. Pick $w\in W$, a word $z\in\A^k\setminus W$ and a letter $e\in\A$, and let $u_t=\iota(w)$ for $t=r$ and $u_t=w$ for $t\ne r$, $1\le t\le m$. Let $c$ be the configuration of spatial period $(m+1)d$ whose period, starting at position $0$, is
\[
u_1\,e^{d-k}\,u_2\,e^{d-k}\cdots u_m\,e^{d-k}\,z\,e^{d-k};
\]
this uses $d\ge k$. The positions $0,d,\dots,(m-1)d$ are occurrences, all of the class $\kappa(w)$ because $\kappa\circ\iota=\kappa$, and $md$ and $-d$ are not occurrences because $c_{[md,md+k)}=c_{[-d,-d+k)}=z$. So these positions form a train of length $m$ whose bit vector $b$ has $b_t$ equal to $\beta(w)+[t=r]$. Other occurrences, created by the fillers, may form further trains; this does not matter. By Lemma~\ref{hw:trainlemma}(iv) and the remark before it, the bit vector of this train in $g(c)$ is $b\circ\sigma_m^{-1}$, which differs from $b$ at $\sigma_m(r)$. Since bits are functions of the configuration, $g(c)\ne c$. So $\Psi$ is injective. Every finite group embeds in some $S_m$ by Cayley's theorem. Finally $\bigoplus_mS_m$ is not finitely generated, and every subgroup of a virtually $\Z^n$ group is finitely generated, as in the proof of Theorem~\ref{thm:54}.
\end{proof}

\begin{remark}\label{hw:H4}
Hypothesis (H4) cannot be dropped from the argument. (H2) guarantees that a flip at $0$ preserves an overlapping occurrence at $i$ and commutes with the flip there, but the flip cell $j$ may be another letter of the word at $i$, so the word at $i$ can move to a different $\iota$-pair, with a different bit or class. Then the bit vectors after $\Psi_m(\rho)$ need not be $b\circ\rho^{-1}$, and both the homomorphism law and the train structure under (H3$'$) can fail; commutation with $F$ is unaffected, since it uses only $I_\sigma(F(c))=I_\sigma(c)$. Rule $59$ below shows that such overlaps occur: a flip changes the protector of the next wall. There (H4) holds because the bit does not read the protector; it would fail if the class were taken to be the $\iota$-pair of the whole word. Alternatively, (H4) can be avoided by building trains only from occurrences that overlap no other occurrence. That set is preserved by $F$ and by flips, since both preserve $\Occ$, and a flip at such an occurrence changes no other word; the injectivity argument then needs a configuration in which the members of the train are of this kind.
\end{remark}

\begin{remark}\label{hw:cent}
The proof uses about $\sigma\in S_m$ only that $b\mapsto b\circ\sigma^{-1}$ is a permutation of $(\Z/2)^m$ commuting with $b\mapsto b+(1,\dots,1)$, which is how $F$ acts on the bit vector of a train when $\varepsilon=1$. For every permutation $\pi$ of $(\Z/2)^m$ commuting with $b\mapsto b+(1,\dots,1)$, replacing the bit vector $b$ of every train of length $m$ by $\pi(b)$ defines an automorphism $\Psi_m(\pi)$ commuting with $F$, by the same proof with $I_\pi(c)$ the set of members of trains of length $m$ at which $\pi(b)$ and $b$ differ, and $\pi\mapsto\Psi_m(\pi)$ is a homomorphism. For injectivity the witness must be a train with an arbitrary bit vector, rather than the vector with one odd bit used for $S_m$. The construction in the proof of Theorem~\ref{hw:general} gives one: choose each $u_t\in\{w,\iota(w)\}$ freely, so that the train at $0,d,\dots,(m-1)d$ has bit vector $b$ with $b_t=\beta(u_t)$, which is arbitrary. Given $\pi\ne\mathrm{id}$, choose $b$ with $\pi(b)\ne b$; then the bit vector of this train in $\Psi_m(\pi)(c)$ is $\pi(b)\ne b$, so $\Psi_m(\pi)(c)\ne c$. For rule $59$ the witness of Theorem~\ref{hw:59} serves in the same way, with the centres $C_t$ chosen freely. So the centralizer of $b\mapsto b+(1,\dots,1)$ in $\Sym((\Z/2)^m)$, which contains $(\Z/2)^m\rtimes S_m$, embeds in $\Aut(X_F)$.
\end{remark}

Theorem~\ref{hw:general} generalizes the mechanism of Theorem~\ref{thm:walls}, though not its statement, which it does not contain. In Theorem~\ref{thm:walls} the walls are fixed and the permuted information sits in cars between them; here the information is a letter inside the wall, and the rest of the wall moves. The common point is that every cell next to the permuted information has a trajectory that does not depend on the context.

The simplest case is a wall that is a single periodic orbit of $F$ in every context.

\begin{corollary}\label{hw:periodic}
Let $P=2h$ be even and let $w_0,\dots,w_{P-1}\in\A^k$ be distinct words, indices modulo $P$, such that:
\begin{enumerate}
\item[(a)] if $c_{[0,k)}=w_s$ then $F(c)_{[0,k)}=w_{s+1}$, for every $c$ and $s$;
\item[(b)] $w_{s+h}$ differs from $w_s$ exactly at a position $j$ with $1\le j\le k-2$, the same for all $s$;
\item[(c)] no configuration has occurrences of $W=\{w_0,\dots,w_{P-1}\}$ at two positions at distance $\delta$ with $1\le\delta\le\max(j,k-1-j)$.
\end{enumerate}
Then $\bigoplus_{m\ge1}S_m$ embeds in $\Aut(X_F)$. In particular each $w_s$ is fixed by $F^P$ in every context, and the flip-all map $\tau_W$ of Corollary~\ref{hw:flipall} is an involutive automorphism of radius $\max(j,k-1-j)$.
\end{corollary}

\begin{proof}
Let $\iota(w_s)=w_{s+h}=F^h$ on $W$. (W1) is (a). For (H1), let $c_{[0,k)}=w_s$. The configurations $\iota_0c$ and $c$ differ only at cell $j$, so $F(\iota_0c)$ and $F(c)$ differ at most at cells $j-1,j,j+1$, which lie in $[0,k)$ because $j$ is interior. On $[0,k)$, $F(\iota_0c)$ carries $w_{s+h+1}$ by (a), which is what $\iota_0F(c)$ carries there; outside $[0,k)$ both agree with $F(c)$. So $F(\iota_0c)=\iota_0F(c)$. For (H3$'$) take $K=\Z/h$, $\kappa(w_s)=s\bmod h$, $\beta(w_s)=\lfloor s/h\rfloor$ for $0\le s<P$, $\theta(x)=x+1$ and $\varepsilon(x)=[x=h-1]$: then $\kappa\circ\iota=\kappa$, $\beta\circ\iota=\beta+1$, and passing from $w_s$ to $w_{s+1}$ changes $\lfloor s/h\rfloor$ exactly when $s\equiv h-1\pmod h$. For (H2) and (H4): if $0,i\in\Occ(c)$ with $i\ne0$, then $|i|>\max(j,k-1-j)$ by (c), so the flip cell of each of the two occurrences lies outside the window of the other. Hence $\iota_0$ leaves $w_i(c)$ unchanged, the new value it writes at cell $j$ depends only on $w_0(c)$, which $\iota_i$ leaves unchanged, and so (H2) and (H4) hold. Finally $W\ne\A^k$ by (c). Theorem~\ref{hw:general} and Corollary~\ref{hw:flipall} apply.
\end{proof}

Condition (b) contains the requirement that $w_s$ and $w_{s+h}$ agree at their end positions $0$ and $k-1$; this is sufficient for (H1) but not necessary, since a disagreement at an end letter that the outside neighbour's update ignores would not break (H1). Conditions (a)--(c) are finite checks, but (H2)--(H4) do not follow from (a) and (b) alone, which is why (c) is assumed.

\subsection{Rule 59R: clocks}\label{hw:sec:59}

Rule $59$ is $f(a,b,c)=\lnot b\lor(\lnot a\land c)$. Call a cell $i$ \emph{triggered} in $c$ if $y_{i-1}=0$ and $y_{i+1}=1$, that is $c_{i-1}\in\{0,2\}$ and $c_{i+1}\in\{1,3\}$. Since $f(a,0,c)=1$ and $f(a,1,c)=[a=0,\,c=1]$, \eqref{eq:F} gives, for $F=F_{59}$:
\begin{equation}\label{hw:clock}
0\mapsto1,\qquad 2\mapsto0,\qquad 1\mapsto\begin{cases}3&\text{triggered}\\2&\text{otherwise,}\end{cases}\qquad 3\mapsto\begin{cases}2&\text{triggered}\\3&\text{otherwise.}\end{cases}
\end{equation}
So every cell is a clock $0\to1\to2\to0$ of period $3$, whose $y$-bit is $1$ exactly at phase $1$; the state $3$ is phase $1$ held, a trigger at phase $1$ holds the clock, and a later trigger releases it. A cell can only be triggered while its $y$-bit is $1$, and only by the pattern ``left neighbour silent, right neighbour firing''. In particular two adjacent cells in the same phase never disturb each other.

Say that a word $u$ of length $n$ is \emph{fixed in every context} if $F(c)_{[0,n)}=u$ for every $c$ with $c_{[0,n)}=u$.

\begin{lemma}\label{hw:nofixed}
For $F=F_{59}$, no word of any length is fixed in every context.
\end{lemma}

\begin{proof}
By \eqref{hw:clock}, a cell with value $0$, $1$ or $2$ changes its value in every context, so a word fixed in every context would be $3^n$. Let $c_{[0,n)}=3^n$, $c_{-1}=0$ and $c_n=1$. The left neighbour of cell $0$ has $y=0$, and its right neighbour (the next $3$, or the outside cell $c_1=1$ when $n=1$) has $y=1$. So cell $0$ is triggered, and $F(c)_0=2$.
\end{proof}

For $p\in\{0,1,2\}$, expressions such as $p+1$ and $p+2$ are taken modulo $3$ and denote cells in $\{0,1,2\}$. Let
\begin{align*}
X_A&=\{\,a\,p\,C\,(p+1)\ :\ p\in\{0,1,2\},\ a\in\{p,3\},\ C\in\{p+1,3\}\,\},\\
X_B&=\{\,a\,p\,C\,(p+2)\ :\ p\in\{0,1,2\},\ a\in\{p,3\},\ C\in\{p+2,3\}\,\},
\end{align*}
two disjoint sets of $12$ words of length $4$. In an \emph{A-wall} the right side is one phase ahead of the left, in a \emph{B-wall} one phase behind. We say $c$ has an A-wall (B-wall) at $i$ if $c_{[i,i+4)}\in X_A$ ($\in X_B$), and call $c_i$, $c_{i+1}$, $c_{i+2}$, $c_{i+3}$ its \emph{protector} $a$, its \emph{left cell} $L=p$, its \emph{centre} $C$ and its \emph{right cell} $R$; the pair $(L,C)$ is its \emph{core}. Write $\delta_A=1$ and $\delta_B=2$, so that $R=L+\delta$ and $C\in\{L+\delta,3\}$.

In this notation the local conservation law of $27$R and $59$R found by Takesue \cite[(4.1), (4.2)]{Takesue1989} says that the set of pairs $a\,p$ with $p\in\{0,1,2\}$ and $a\in\{p,3\}$ is invariant: such a pair remains one in every context. His wall is a run $a\,p\,p\cdots p$ of at least three cells with $a\in\{p,3\}$, in which every cell but the first passes from $p$ to $p+1$ in every context \cite[(4.3)]{Takesue1989}. A word of $X_A\cup X_B$ consists of two such pairs, the protector with the left cell and the centre with the right cell, whose phases differ by $\delta$. The cycles of the cores, the bit carried by the centre and the flips (Lemmas~\ref{hw:59walls}--\ref{hw:59overlap}) are not in \cite{Takesue1989}. For $59$R the invariance of the pairs is contained in the proof of Lemma~\ref{hw:59walls}, in the images of the protector and of the left cell, where it is needed together with the images of the other cells.

\begin{lemma}\label{hw:59walls}
Let $F=F_{59}$ and $X\in\{A,B\}$. If $c$ has an $X$-wall at $i$ with left cell $p$, then $F(c)$ has an $X$-wall at $i$ with left cell $p+1$, and $F(c)_{i+1},F(c)_{i+2},F(c)_{i+3}$ depend only on the core of $c$ at $i$, not on the rest of $c$. The cores run through the $6$-cycles
\begin{align*}
A:&\quad(0,1)\to(1,3)\to(2,3)\to(0,3)\to(1,2)\to(2,0)\to(0,1),\\
B:&\quad(0,2)\to(1,0)\to(2,1)\to(0,3)\to(1,3)\to(2,3)\to(0,2).
\end{align*}
Moreover $F(c)$ has an $X$-wall at $i$ if and only if $c$ does, and the same holds for $F^{-1}$.
\end{lemma}

\begin{proof}
By \eqref{hw:clock}, cells with value $0$ or $2$ advance regardless of their neighbours, so only cells with value $1$ or $3$ need attention. The six cases are:
\[
\begin{array}{ccccl}
 & (a,L,C,R) & (a',L',C',R') & C\mapsto C' & \text{triggered}\\ \hline
A,\,p=0 & (\{0,3\},0,\{1,3\},1) & (\{1,3\},1,\{3,2\},2) & 1\mapsto3,\ 3\mapsto2 & C\\
A,\,p=1 & (\{1,3\},1,\{2,3\},2) & (\{2,3\},2,\{0,3\},0) & 2\mapsto0,\ 3\mapsto3 & \text{possibly }a\\
A,\,p=2 & (\{2,3\},2,\{0,3\},0) & (\{0,3\},0,\{1,3\},1) & 0\mapsto1,\ 3\mapsto3 & \text{none}\\
B,\,p=0 & (\{0,3\},0,\{2,3\},2) & (\{1,3\},1,\{0,3\},0) & 2\mapsto0,\ 3\mapsto3 & \text{none}\\
B,\,p=1 & (\{1,3\},1,\{0,3\},0) & (\{2,3\},2,\{1,3\},1) & 0\mapsto1,\ 3\mapsto3 & \text{possibly }a\\
B,\,p=2 & (\{2,3\},2,\{1,3\},1) & (\{0,3\},0,\{2,3\},2) & 1\mapsto3,\ 3\mapsto2 & C
\end{array}
\]
We justify the entries. If $L=1$ (case $p=1$), its left neighbour $a\in\{1,3\}$ has $y=1$, so $L$ is not triggered and $L'=2$. If $R=1$ (cases A, $p=0$ and B, $p=2$), its left neighbour $C\in\{1,3\}$ has $y=1$, so $R'=2$. The centre, when it equals $1$ or $3$, is triggered exactly when $y_L=0$ and $y_R=1$, that is when $p\ne1$ and $R=1$: this happens in the cases A, $p=0$ and B, $p=2$ and in no other case (in the other cases $p=1$, or $R\in\{0,2\}$). The protector: $a=0\mapsto1$, $a=2\mapsto0$, $a=1\mapsto\{2,3\}$, and $a=3$ is triggered only if $y_L=1$, so $a=3\mapsto3$ when $p\in\{0,2\}$ and $a=3\mapsto\{2,3\}$ when $p=1$. In every case $a'\in\{p+1,3\}$. Thus each row maps the family into itself with $p$ replaced by $p+1$, and the images of $L$, $C$, $R$ are determined by $(L,C)$, since $R=L+\delta$; only $a'$ depends on the context. Reading off the core images gives the two $6$-cycles. The last statement follows from Lemma~\ref{hw:exact} applied to $W=X_A$ and to $W=X_B$.
\end{proof}

Only the protector depends on the context, through its own left neighbour. The left cell $L$ of a wall is therefore a column whose whole trajectory, forwards and backwards, is determined by the wall (backwards because the core map is a bijection), and the same holds for $C$ and $R$. Cells to the left of the wall feel the right side only through the $y$-sequence of $L$, and cells to the right feel the left side only through that of $R$.

At a wall the centre takes one of two values, $C\in\{L+\delta,3\}$. Let $\iota$ exchange them; on each family it is a fixed-point-free involution changing only position $j=2$. On the core, $\iota$ is the shift by three steps of the $6$-cycle, so $\iota$ acts on the wall as $F^3$ acts on its core, while leaving every other cell alone. Define
\[
\beta_A(a\,L\,C\,R)=[C=3]+[L=2],\qquad \beta_B(a\,L\,C\,R)=[C=3]+[L=1]\qquad(\bmod 2).
\]

\begin{lemma}\label{hw:59flip}
Let $F=F_{59}$, $X\in\{A,B\}$, and let $c$ have an $X$-wall at $i$. Then $F(\iota_ic)=\iota_iF(c)$ and $F^{-1}(\iota_ic)=\iota_iF^{-1}(c)$. Moreover $\beta_X\circ\iota=\beta_X+1$, and $\beta_X(F(c)_{[i,i+4)})=\beta_X(c_{[i,i+4)})+1$ in every context.
\end{lemma}

\begin{proof}
The configurations $\iota_ic$ and $c$ differ only at the centre $i+2$, so $F(\iota_ic)$ and $F(c)$ can differ only at $i+1$, $i+2$, $i+3$. The images of $L$ and $R$ computed in the proof of Lemma~\ref{hw:59walls} do not depend on which of the two values $C$ has: $L$ and $R$ read $C$ only through its $y$-bit, and either $y_C=1$ for both values (A, $p=0$; B, $p=2$), or the cell reading $C$ is $L=1$ with $y_a=1$ and is never triggered ($p=1$), or $L$ and $R$ are in $\{0,2\}$ and ignore their neighbours. The centre's update is, by the table, a bijection from $\{L+\delta,3\}$ onto $\{L'+\delta,3\}$, the centre pair of the image wall. Hence $F(\iota_ic)_{i+2}$ is the element of that pair other than $F(c)_{i+2}$, that is, $(\iota_iF(c))_{i+2}$, and $F(\iota_ic)=\iota_iF(c)$. Applying this to $F^{-1}(c)$, which has an $X$-wall at $i$ by Lemma~\ref{hw:59walls}, gives the statement for $F^{-1}$. The partners under $\iota$ are $(0,1)\leftrightarrow(0,3)$, $(1,3)\leftrightarrow(1,2)$, $(2,3)\leftrightarrow(2,0)$ for A and $(0,2)\leftrightarrow(0,3)$, $(1,0)\leftrightarrow(1,3)$, $(2,1)\leftrightarrow(2,3)$ for B, and along the A-cycle $\beta_A$ takes the values $0,1,0,1,0,1$, along the B-cycle $\beta_B$ takes the values $0,1,0,1,0,1$; both statements about $\beta$ follow, since $\beta_X$ depends only on the core.
\end{proof}

In spacetime the centre column of an A-wall has $y$-history $1,1,1,1,0,0$ repeated, and $\iota$ replaces it by its shift by half a period; the neighbours read that column only at times when the two histories agree. By Lemma~\ref{hw:59flip}, $F^n(\iota_ic)=\iota_iF^n(c)$ for all $n\in\Z$, so flipping the bit of one wall changes exactly one column of the spacetime diagram, at all times.

\begin{lemma}\label{hw:59overlap}
Let $F=F_{59}$ and $W=X_A$, with $k=4$ and $j=2$.
\begin{enumerate}
\item[(i)] No configuration has A-walls at $i$ and $i+1$.
\item[(ii)] If $c$ has A-walls at $0$ and $i\ne0$, then $\iota_0c$ has an A-wall at $i$, $\iota_0\iota_ic=\iota_i\iota_0c$, and the bit $\beta_A$ of the wall at $i$ is the same in $c$ and in $\iota_0c$.
\end{enumerate}
Hence \textup{(W1)}, \textup{(H1)}, \textup{(H2)}, \textup{(H3)} with $\varepsilon=1$, and \textup{(H4)} hold for $X_A$, $\iota$ and $\beta_A$.
\end{lemma}

\begin{proof}
(i) Suppose $c$ has A-walls at $i$ and $i+1$, and let $p=c_{i+1}$. The left cell of the second wall is $c_{i+2}$, the centre of the first, which must then lie in $\{0,1,2\}$; so $c_{i+2}=p+1$. The protector of the second wall is $c_{i+1}=p$, which must lie in $\{p+1,3\}$; impossible.

(ii) The flip $\iota_0$ changes only cell $2$, writing the element of $\{c_1+1,3\}$ other than $c_2$; it reads only $c_1$ and $c_2$. The window $[i,i+4)$ contains cell $2$ only if $-1\le i\le2$, so by (i) only if $i=2$. In that case cell $2$ is the protector of the wall at $2$, whose left cell is $c_3=c_1+1$, so its allowed protectors are $\{c_1+1,3\}$, the same pair; the wall at $2$ survives. In all cases the wall at $i$ survives, and its core $(c_{i+1},c_{i+2})$ is untouched, since $i+1\ne2$ by (i) and $i+2\ne2$. So its bit is unchanged. The same argument with $0$ and $i$ exchanged shows that $\iota_i$ does not touch $c_1$ or $c_2$; as each flip reads only the core of its own wall, the two flips commute. The hypotheses: (W1) is Lemma~\ref{hw:59walls}, (H1) and (H3) are Lemma~\ref{hw:59flip}, (H2) and (H4) are (ii).
\end{proof}

The same holds for walls of both families together; this is Lemma~\ref{hw:59classes}(i) below. Overlapping walls then occur only at distances $2$ and $3$, and by computation on all words of length at most $9$ both distances occur for all four type pairs AA, AB, BA, BB.

\begin{theorem}\label{hw:59}
For every $m\ge1$ there is an injective homomorphism $\Psi_m\colon S_m\to\Aut(X_{F_{59}})$ whose image consists of block maps of radius at most $2m+2$, and the images for different $m$ commute. So $\bigoplus_{m\ge1}S_m$ embeds in $\Aut(X_{F_{59}})$, which therefore contains every finite group and is not virtually $\Z^n$ for any $n$. The same holds for the rules $115$, $206$ and $220$, which are conjugate to $59$ by reflection and complement.
\end{theorem}

Here a train of length $m$ is a maximal chain of A-walls at gap $2$: A-walls at $i,i+2,\dots,i+2(m-1)$ and none at $i-2$ or $i+2m$. Consecutive walls overlap, the centre of each being the protector of the next, so a train looks like a staircase $a\;p\;C_1\,(p+1)\,C_2\,(p+2)\cdots C_m\,(p+m)$. This is not literally an instance of Theorem~\ref{hw:general}, which requires the gap $d$ to be at least the word length $k=4$, so that the members of a train can be written down independently; at gap $2$ they overlap, and the injectivity witness has to be built by hand. (With $d=4$, Theorem~\ref{hw:general} does apply to $X_A$, by Lemma~\ref{hw:59overlap}, with the weaker radius bound $4m+2$.)

\begin{proof}
By Lemma~\ref{hw:59overlap}, the hypotheses of Lemma~\ref{hw:trainlemma} hold for $W=X_A$ with $|K|=1$ and $\varepsilon=1$. Applying that lemma with $d=2$, the maps $\Psi_m(\sigma)$ are automorphisms of radius at most $2m+\max(2,1)=2m+2$, commuting with $F$, forming homomorphisms of the $S_m$ whose images commute for different $m$. It remains to prove injectivity, as in the proof of Theorem~\ref{hw:general} but with a different witness. Fix $m$, $p\in\{0,1,2\}$ and choices $C_t\in\{p+t,3\}$ for $1\le t\le m$, and let
\[
c=\cdots3\,3\,3\;p\;C_1\,(p+1)\,C_2\,(p+2)\cdots C_m\,(p+m)\;3\,3\,3\cdots,
\]
with $c_{-1}=p$, $c_{2t-2}=C_t$ and $c_{2t-1}=p+t$, and all other cells equal to $3$. For $1\le t\le m$ the window $c_{[2t-4,2t)}$ is $(c_{2t-4},\,p+t-1,\,C_t,\,p+t)$, with protector $c_{-2}=3$ for $t=1$ and $C_{t-1}\in\{p+t-1,3\}$ for $t\ge2$; so it is an A-wall. There is no A-wall at $-4$, since $c_{-3}=3$ is not a possible left cell, and none at $2m-2$, since $c_{2m+1}=3$ is not a possible right cell. So the walls at $-2,0,\dots,2m-4$ form a train of length $m$, and its bits $\beta_A=[C_t=3]+[p+t-1=2]$ can be prescribed arbitrarily by the choice of the $C_t$. (By direct check, for $m\le7$ these are the only walls of $c$ of either family, but this is not needed.) Given $(\sigma_{m'})_{m'}$ with $\sigma_m(r)\ne r$, choose the bits to be $1$ at $r$ and $0$ elsewhere; as in the proof of Theorem~\ref{hw:general}, the bit vector of this train in $\Psi((\sigma_{m'}))(c)$ is $b\circ\sigma_m^{-1}\ne b$. The statement for the conjugate rules follows because conjugation by the reflection or by the cellwise complement induces an isomorphism of the automorphism groups of the spacetime subshifts.
\end{proof}

By Remark~\ref{hw:cent}, $\Aut(X_{F_{59}})$ also contains the centralizer of $b\mapsto b+(1,\dots,1)$ in $\Sym((\Z/2)^m)$ for every $m$.

\subsection{The radius-2 census of rule 59R}\label{hw:sec:census}

For $X\in\{A,B\}$, call an $X$-wall at $i$ \emph{level} if $c_{i+4}=c_{i+3}$ (the cell after $R$ is in phase with $R$), and \emph{doubly protected} if $c_i=c_{i+1}$ and $c_{i-1}\in\{c_{i+1},3\}$ (the protector is itself protected). For a property $Q$ of $X$-walls, flipping the centre of every $X$-wall with property $Q$, all at once, defines a block map; we write $\tau_{X,=}$, $\tau_{X,\ne}$ for the maps that flip the level and the non-level $X$-walls, $\lambda_{X,\mathrm{pp}}$, $\lambda_{X,\lnot\mathrm{pp}}$ for those that flip the doubly protected and the other $X$-walls, and $\tau_X=\tau_{X,=}\tau_{X,\ne}=\lambda_{X,\mathrm{pp}}\lambda_{X,\lnot\mathrm{pp}}$ for the flip of all $X$-walls, which is the map $\tau_W$ of Corollary~\ref{hw:flipall} for $W=X_A$ or $W=X_B$. The maps $\tau_{X,=}$ and $\tau_{X,\ne}$ have radius $2$ and differ from the identity on $12$ and $36$ entries of the table for $X=A$; the maps $\lambda$ have radius $3$. Let
\[
T_R=\langle\tau_{A,=},\tau_{A,\ne},\tau_{B,=},\tau_{B,\ne}\rangle,\qquad T_L=\langle\lambda_{A,\mathrm{pp}},\lambda_{A,\lnot\mathrm{pp}},\lambda_{B,\mathrm{pp}},\lambda_{B,\lnot\mathrm{pp}}\rangle,\qquad T=T_RT_L .
\]
The \emph{class} of an $X$-wall at $i$ is the triple $(X,\text{level or not},\text{doubly protected or not})$; there are $8$ possible classes. The class reads the cells $i-1$, $i$, $i+1$, $i+3$, $i+4$ (the family is determined by $L$ and $R$), but not the centre $i+2$. Let $W=X_A\cup X_B$, with $k=4$, $j=2$, and $\iota$ exchanging the two values of the centre in each family. The flip $\iota_i$ of a wall at $i$ writes only cell $i+2$, and reads only cells $i+1$, $i+2$, $i+3$, since the new centre is the other element of $\{L+\delta,3\}$ and $\delta=R-L$. Let $Y=\{a\,p\,p:\ p\in\{0,1,2\},\ a\in\{p,3\}\}$. An $X$-wall at $i$ is doubly protected if and only if $c_{[i-1,i+2)}\in Y$, since $L=c_{i+1}\in\{0,1,2\}$. It is level if and only if $c_{[i+2,i+5)}\in Y$, since the centre always lies in $\{R,3\}$.

\begin{lemma}\label{hw:59classes}
Let $F=F_{59}$ and $W=X_A\cup X_B$.
\begin{enumerate}
\item[(i)] No configuration has walls, of either family, at two positions at distance $1$. The pair $(W,\iota)$ satisfies \textup{(W1)}, \textup{(H1)} and \textup{(H2)}, so Lemma~\ref{hw:flips}(a), (b), (d) applies to it.
\item[(ii)] If $c$ has walls at $i$ and $i'$, possibly $i=i'$, then the wall at $i$ has the same class in $\iota_{i'}c$ as in $c$. Hence for every set $I$ of walls of $c$, the configuration $\iota_Ic$ has the same walls as $c$, with the same classes.
\item[(iii)] $F(c)$ has a wall at $i$ of class $x$ if and only if $c$ does.
\item[(iv)] All $8$ classes occur.
\end{enumerate}
\end{lemma}

\begin{proof}
(i) Suppose $c$ has walls at $i$ and $i+1$, and let $p=c_{i+1}$ and $\delta$ belong to the first wall. The left cell of the second wall is $c_{i+2}$, the centre of the first, so it lies in $\{0,1,2\}$, and hence $c_{i+2}=p+\delta$. The protector of the second wall is $c_{i+1}=p$, which must lie in $\{p+\delta,3\}$. This is impossible. (W1) is Lemma~\ref{hw:59walls} and (H1) is Lemma~\ref{hw:59flip}, applied to each family. For (H2), let $c$ have walls at $0$ and $i\ne0$, so $|i|\ge2$. The cell $2$ written by $\iota_0$ lies in the window $[i,i+4)$ only if $-1\le i\le2$, hence only if $i=2$. In that case cell $2$ is the protector of the wall at $2$ and the centre of the wall at $0$. The centre pair of the wall at $0$ is $\{R_0,3\}=\{c_3,3\}$, and $c_3$ is the left cell of the wall at $2$, so this pair is exactly the set of allowed protectors of that wall. In every case the wall at $i$ survives, with its family. The flip $\iota_i$ writes cell $i+2$, which is not among the cells $1$, $2$, $3$ read by $\iota_0$, since $|i|\ne1$ and $i\ne0$. Symmetrically, $\iota_0$ writes no cell that $\iota_i$ reads, so the two flips commute.

(ii) If $i'=i$, the flip changes only the centre, which the class does not read. Otherwise $|i'-i|\ge2$. The cell $i'+2$ written by $\iota_{i'}$ is one of the cells $i-1$, $i$, $i+4$ read by the class only if $i'\in\{i-3,i-2,i+2\}$, because $i'+2\in\{i+1,i+3\}$ would give $|i'-i|=1$. The wall at $i$ survives by (i), and we check these three cases.
\begin{itemize}
\item $i'=i+2$: cell $i+4$ is the centre of the wall at $i+2$, whose left cell is $c_{i+3}$. Both possible values of that centre differ from its left cell. So the wall at $i$ is not level, before or after the flip.
\item $i'=i-2$: the wall at $i-2$ has left cell $c_{i-1}\in\{0,1,2\}$ and right cell $c_{i+1}\ne c_{i-1}$. So $c_{i-1}\notin\{c_{i+1},3\}$, and the wall at $i$ is not doubly protected, before or after the flip. The flip writes $c_i$ but not $c_{i-1}$.
\item $i'=i-3$: the flip exchanges $c_{i-1}$, the centre of the wall at $i-3$, between $c_i$ (its right cell) and $3$. If $c_i=c_{i+1}$, both values lie in $\{c_{i+1},3\}$, so the wall at $i$ is doubly protected before and after the flip. If $c_i\ne c_{i+1}$, it is doubly protected in neither.
\end{itemize}
For a set $I$, the walls and the class of a given wall read finitely many cells, so the second statement follows from the first, Lemma~\ref{hw:flips}(a), and induction, as in the proof of Lemma~\ref{hw:flips}.

(iii) The family is preserved by Lemma~\ref{hw:59walls}. By the remark before the lemma, it remains to show that $F$ and $F^{-1}$ preserve the occurrences of $Y$. By Lemma~\ref{hw:exact} it suffices to check (W1) for $Y$ and $F$. Let $c_{[0,3)}=a\,p\,p$ with $a\in\{p,3\}$. Cells $1$ and $2$ have left neighbours in $\{p,3\}$. For $p\in\{0,2\}$ they advance to $p+1$ by \eqref{hw:clock}. For $p=1$ their left neighbours have $y=1$, so they are not triggered and also become $2=p+1$. The protector $a=p$ becomes $p+1$, or $3$ when $p=1$ and it is triggered. The protector $a=3$ can be triggered only if $y_1=1$, that is only if $p=1$, and then it becomes $2=p+1$. So $F(c)_{[0,3)}\in Y$.

(iv) Let $X$ have $\delta$, and take a window $c_{[i-1,i+5)}=u\,a\,p\,C\,(p+\delta)\,v$ with $C\in\{p+\delta,3\}$. With $a=u=p$ the wall at $i$ is doubly protected; with $a=3$ it is not, since then $c_i\ne c_{i+1}$. With $v=p+\delta$ it is level; with $v=3$ it is not. For instance $000111$, $000113$, $030111$, $030113$ give the four A-classes, and $000222$, $000223$, $030222$, $030223$ give the four B-classes, each for the wall at position $1$.
\end{proof}

\begin{proposition}\label{hw:census}
Let $F=F_{59}$. The $64$ elements of $T$ are distinct automorphisms commuting with $F$ and with each other, and each is an involution or the identity; $T_R\cong T_L\cong(\Z/2)^4$, $T_R\cap T_L=\{1,\tau_A,\tau_B,\tau_A\tau_B\}$, and $T\cong(\Z/2)^6$. The group $\Cent^{\times}_2(F)$ has exactly $133$ elements, and each of them has exactly one expression of one of the following forms:
\begin{enumerate}
\item[(i)] $\sh^aF^b$ with $|a|+|b|\le2$ ($13$ maps);
\item[(ii)] $F^bt$ with $|b|\le2$ and $1\ne t\in T_R$ ($75$ maps);
\item[(iii)] $\sh F^bt$ with $|b|\le1$ and $1\ne t\in T_L$ ($45$ maps).
\end{enumerate}
\end{proposition}

\begin{proof}
\emph{The group $T$.} This part is exact. For a set $U$ of classes, let $I_U(c)$ be the set of walls of $c$ whose class lies in $U$, and $\varphi_U(c)=\iota_{I_U(c)}c$. This is a block map of radius $3$, since cell $n$ changes only if there is a wall at $n-2$, and its class is read from $c_{[n-3,n+3)}$. By Lemma~\ref{hw:59classes}(ii), $I_V(\iota_Jc)=I_V(c)$ for every set $J$ of walls of $c$. So Lemma~\ref{hw:flips}(b) gives
\[
\varphi_V\varphi_U(c)=\iota_{I_V(c)}\iota_{I_U(c)}c=\iota_{I_U(c)\triangle I_V(c)}c=\varphi_{U\triangle V}(c),
\]
and $U\mapsto\varphi_U$ is a homomorphism from the group of subsets of the $8$ classes under $\triangle$, which is $(\Z/2)^8$. In particular each $\varphi_U$ is an involution or the identity, hence an automorphism, and any two commute. By Lemma~\ref{hw:59classes}(iii), $I_U(F(c))=I_U(c)$, so $F\varphi_U=\varphi_UF$ by Lemma~\ref{hw:flips}(d). The homomorphism is injective. If $U\ne\emptyset$, pick $x\in U$ and, by Lemma~\ref{hw:59classes}(iv), a configuration $c$ with a wall of class $x$; then $\varphi_U$ changes the centre of that wall, so $\varphi_U\ne\mathrm{id}$.

Write $R_{X,=}$, $R_{X,\ne}$ for the sets of the two level and the two non-level classes of family $X$, and $L_{X,\mathrm{pp}}$, $L_{X,\lnot\mathrm{pp}}$ for the sets of the two doubly protected and the two other classes. Then $\tau_{X,=}=\varphi_{R_{X,=}}$, $\tau_{X,\ne}=\varphi_{R_{X,\ne}}$, $\lambda_{X,\mathrm{pp}}=\varphi_{L_{X,\mathrm{pp}}}$ and $\lambda_{X,\lnot\mathrm{pp}}=\varphi_{L_{X,\lnot\mathrm{pp}}}$. Hence the map from $(\Z/2)^8$, indexed by the eight generators, onto $T$ that sends $s$ to the product of the generators with exponents $s$ sends $s$ to $\varphi_{U(s)}$. Here the class $(X,l,q)$ lies in $U(s)$ if and only if the exponents of $\tau_{X,l}$ and $\lambda_{X,q}$ have sum $1$. So $s$ lies in the kernel if and only if, for each $X$, the four exponents of $\tau_{X,=}$, $\tau_{X,\ne}$, $\lambda_{X,\mathrm{pp}}$, $\lambda_{X,\lnot\mathrm{pp}}$ are equal. The kernel has dimension $2$, so $|T|=2^6=64$ and $T\cong(\Z/2)^6$. The four sets $R_{X,l}$ are disjoint and nonempty, so $T_R\cong(\Z/2)^4$, and likewise $T_L\cong(\Z/2)^4$. A set of classes that is a union of sets $R_{X,l}$ and also a union of sets $L_{X,q}$ contains, for each $X$, either none or all four classes of $X$. So $T_R\cap T_L=\{\varphi_\emptyset,\varphi_A,\varphi_B,\varphi_{A\cup B}\}=\{1,\tau_A,\tau_B,\tau_A\tau_B\}$, where $\varphi_X$ flips all $X$-walls and so equals $\tau_X$. This also gives $\tau_X=\tau_{X,=}\tau_{X,\ne}=\lambda_{X,\mathrm{pp}}\lambda_{X,\lnot\mathrm{pp}}$.

\emph{The census.} This part is a computation. $\Cent^{\times}_2(F)$ was enumerated as in Section~\ref{sec:method}, and every solution passed the pair-graph test. The tables of the maps in (i)--(iii) were computed by exact composition, and each of the $133$ tables was matched against them; each matched exactly one, so the three lists are disjoint. In (iii), $\sh\lambda$ has radius $2$ although $\lambda$ has radius $3$, because $\lambda$ reads three cells to the left of the flip cell and one to the right.
\end{proof}

The census is consistent with the picture of Theorem~\ref{hw:59}: radius $2$ is too small for any train of length $m\ge2$, and what appears instead are flips of whole classes of walls, defined by properties of the walls and their immediate surroundings that $F$ and the flips preserve.

\begin{proposition}\label{hw:group}
Let $F=F_{59}$ and $G=\langle\sh,F,\Cent^{\times}_2(F)\rangle$. Then:
\begin{enumerate}
\item[(i)] $G=\langle\sh,F\rangle\,T$, and $T\cong(\Z/2)^6$ is central in $G$;
\item[(ii)] $\langle\sh,F\rangle\cap T=\{1\}$;
\item[(iii)] $G\cong\Z^2\times(\Z/2)^6$.
\end{enumerate}
Every $t\ne1$ in $T$ has order $2$, and every $\sh^aF^bt$ with $(a,b)\ne(0,0)$ has infinite order.
\end{proposition}

\begin{proof}
(i) By Proposition~\ref{hw:census} the generators of $G$ are of the form $\sh^aF^bt$ with $t\in T$. The elements of $T$ are automorphisms commuting with $F$, hence with $\sh$ and $F$, and with each other; so $T$ is central and $G=\langle\sh,F\rangle T$.

(ii) If $\sh^aF^b=t\in T$, then $\sh^{2a}F^{2b}=t^2=1$, so $a=b=0$ by Theorem~\ref{cl:thm:A}(3); rule $59$ depends on its first and third arguments.

(iii) By Theorem~\ref{cl:thm:A}(3), $\langle\sh,F\rangle\cong\Z^2$. By (ii) and since $T$ is central, $G\cong\langle\sh,F\rangle\times T\cong\Z^2\times(\Z/2)^6$.

For the orders: $(\sh^aF^bt)^2=\sh^{2a}F^{2b}$, which has infinite order if $(a,b)\ne(0,0)$.
\end{proof}

\begin{corollary}\label{hw:order}
$F_{59}$ has infinite order, and $\langle\sh,F_{59},\Cent^{\times}_2(F_{59})\rangle\cong\Z^2\times(\Z/2)^6$.
\end{corollary}

\begin{proof}
Theorem~\ref{cl:thm:A} and Proposition~\ref{hw:group}(iii).
\end{proof}

Some orbits of $F_{59}$ move slowly. Let $c_0$ be the configuration of spatial period $4000$ with period $(021)^{666}\,02\,0^{2000}$; it has no wall of either family, and a direct computation shows $F^n(c_0)\ne c_0$ for $1\le n\le400{,}000$. At the interface between the staircase $(021)^{666}02$ and the uniform block $0^{2000}$ of $c_0$, $F$ behaves like a binary odometer: the disturbance first reaches depth $k$ at times close to $2^{k+2}$ (at $t=501$, $1014$, $2037$, \dots, $131{,}055$ for $k=7,\dots,15$), so the period of $c_0$ under $F$, which is finite because $c_0$ is spatially periodic, is very large. Apart from walls, the dynamics of rule $59$R is rather tame. A single changed cell on random rings stayed, in all trials, between the cores of the nearest walls on either side, and after transients most columns are periodic in time.

The walls of rule $59$R are not time-periodic in the sense of Corollary~\ref{hw:periodic}, because the protector depends on the context. For $p\not\equiv0\pmod6$ no word of $X_A\cup X_B$ is fixed by $F^p$ in every context, since $F^p$ moves the core along its $6$-cycle; for $p\in\{6,12,30,60,120\}$ a computer search found, for every word of $X_A\cup X_B$, a context in which $F^p$ changes the protector. We have not proved this for all multiples of $6$; nothing depends on it, since Lemma~\ref{hw:trainlemma} applies to the family $X_A$ as it stands.

\subsection{Time-periodic walls: rules 127R, 179R and 223R}\label{hw:sec:tp}

For these three rules the walls of Takesue are pairs of adjacent cells whose images do not depend on the context: for $127$R the pairs $b\,\kappa(b)$, $b\in\A$, and for $223$R the pairs $b\,b$, on which $F$ acts by the $4$-cycle $0\mapsto1\mapsto3\mapsto2\mapsto0$ in both cells, so that they have period $4$; and for $179$R the pairs of the two cycles $01\to12\to20\to01$ and $10\to21\to02\to10$, of period $3$ \cite[Fig.~2 and Table~III, walls (xii), (xiii), (xv), (xvi)]{Takesue1989}. Each word of the next proposition consists of two such pairs with one cell between them. The cell in the middle is the flip cell, and with it the words have period $8$, $8$ and $6$, twice the period of the pairs.

\begin{proposition}\label{hw:tp}
For $f\in\{127,179,223\}$ the following cycles of words of length $5$ satisfy the hypotheses of Corollary~\ref{hw:periodic} for $F=F_f$, with $k=5$, $j=2$ and $\iota=F^{P/2}$ on $W$:
\begin{align*}
127:&\quad 03112\to12230\to30021\to21103\to03312\to12330\to30221\to21003\quad(P=8),\\
179:&\quad 01110\to12321\to20302\to01310\to12221\to20002\quad(P=6),\\
223:&\quad 00122\to11300\to33211\to22133\to00322\to11200\to33011\to22033\quad(P=8).
\end{align*}
Hence $\bigoplus_{m\ge1}S_m$ embeds in $\Aut(X_{F_f})$, and so in $\Aut(X_{F_{f'}})$ for the conjugate rules $f'=254$, $205$ and $251$. Each word is fixed by $F^P$ in every context.
\end{proposition}

\begin{proof}
(a) The images of the three interior cells of a word of length $5$ depend only on the word, and the arrows above are the result of evaluating \eqref{eq:F} on them and on the end cells. It remains to see that the images of the end cells do not depend on the outside neighbours, which enter only through their $y$-bits. For $f=127$, $f(a,b,c)=\lnot(a\land b\land c)$ depends on $a$ only when $b=c=1$ and on $c$ only when $a=b=1$; for $f=179$ the same is true, since $f(0,b,c)=f(1,b,c)$ unless $b=c=1$ and $f(a,b,0)=f(a,b,1)$ unless $a=b=1$; so for these two rules it suffices that the $y$-bits of the first two and of the last two letters of each word are not both $1$, which one reads off the words ($y$ is the parity of the letter: for $127$ the $y$-words are $01110$, $10010$, $10001$, $01101$, $01110$, $10110$, $10001$, $01001$, and for $179$ they are $01110$, $10101$, $00100$, $01110$, $10001$, $00000$). For $f=223$, $f(a,b,c)=\lnot(a\land\lnot b\land c)$ depends on $a$ only when $(b,c)=(0,1)$ and on $c$ only when $(a,b)=(1,0)$, and the $y$-words $00100$, $11100$, $11011$, $00111$, $00100$, $11000$, $11011$, $00011$ never begin with $01$ or end with $10$.
(b) In each list the words $P/2$ apart differ exactly at position $2$, as one reads off: for instance $03112$ and $03312$, $01110$ and $01310$, $00122$ and $00322$.
(c) Comparing suffixes with prefixes shows that two words of $W$ can overlap only at offsets $3$ and $4$ for $127$ and $223$, and only at offset $4$ for $179$; in particular never at offsets $1$ or $2=\max(j,k-1-j)$.
Corollary~\ref{hw:periodic} applies. The rules $254$, $205$, $251$ are conjugate to $127$, $179$, $223$ by reflection and complement.
\end{proof}

For $179$, where $h=P/2=3$ is odd, one can also take $|K|=1$: $\beta(w_s)=s\bmod2$ satisfies (H3) with $\varepsilon=1$. For $127$ and $223$, where $h=4$, (H3) with a constant $\varepsilon$ is impossible, since it would give $\beta(\iota w)=\beta(F^4w)=\beta(w)+4\varepsilon=\beta(w)$; so the classes of (H3$'$) are genuinely needed there, and trains consist of walls in the same phase modulo $h$. By Corollary~\ref{hw:periodic} the flip of all walls is in each case an involutive automorphism of radius $2$ commuting with $F$; this was also confirmed by composing tables exactly at radius $4$ and by the pair-graph test. Rings of length $1500$ seeded with trains at gaps $5$, $6$ and $8$, with walls of other classes mixed in, confirmed commutation with $F$ and $F^{-1}$ and the homomorphism law, with no failures. We have not computed $\Cent^{\times}_2(F_f)$ for these three rules.

\subsection{Rule 27R}\label{hw:sec:27}

The class $\{27,83,202,216\}$ carries a version of the rule-$59$ mechanism with windows of length $6$. Rule $27$ has $f(a,0,c)=\lnot(a\land c)$ and $f(a,1,c)=\lnot a\land c$, so it differs from rule $59$ only at $(a,b,c)=(1,0,1)$. Call a cell $i$ \emph{pinched} if $y_{i-1}=y_{i+1}=1$. By \eqref{eq:F}, $F_{27}$ acts on cells with values $1$ and $3$ as in \eqref{hw:clock}, and
\begin{equation}\label{hw:clock27}
0\mapsto\begin{cases}0&\text{pinched}\\1&\text{otherwise,}\end{cases}\qquad 2\mapsto\begin{cases}1&\text{pinched}\\0&\text{otherwise.}\end{cases}
\end{equation}
Let
\[
W_{27}=\{\,a\,p\,p\,C\,(p+1)\,(p+1)\ :\ p\in\{0,1,2\},\ a\in\{p,3\},\ C\in\{p+1,3\}\,\},
\]
twelve words. As for rule $59$, call $a$ the \emph{protector}, $C$ the \emph{centre} and $(p,C)$ the \emph{core}. Let $\iota$ exchange the two values of $C$, at position $j=3$, and let
\[
\beta(a\,p\,p\,C\,(p+1)\,(p+1))=[C=3]+[p=2]\pmod 2.
\]

\begin{theorem}\label{hw:27}
Let $F=F_{27}$. The family $W_{27}$, with $k=6$, $j=3$, $\iota$ and $\beta$ as above, satisfies \textup{(W1)}, \textup{(H1)}, \textup{(H2)}, \textup{(H3)} with $\varepsilon=1$, and \textup{(H4)}. Consequently, for every $d\ge6$ and $m\ge1$ the maps $\Psi_m$ of Definition~\ref{hw:train} form an injective homomorphism $S_m\to\Aut(X_{F_{27}})$ whose image consists of block maps of radius at most $md+3$, and the images for different $m$ commute. So $\bigoplus_{m\ge1}S_m$ embeds in $\Aut(X_{F_f})$ for $f\in\{27,83,202,216\}$; each of these groups contains every finite group and is not virtually $\Z^n$ for any $n$. The map that flips every occurrence of $W_{27}$ is an involutive automorphism of radius $3$ commuting with $F_{27}$.
\end{theorem}

\begin{proof}
\emph{(W1) and the core.} Let $c_{[0,6)}=a\,p\,p\,C\,(p+1)\,(p+1)$. The cells $1,\dots,4$ read only cells of the word. Cells $0$ and $5$ read the outside neighbours, and only through their $y$-bits. By \eqref{hw:clock} and \eqref{hw:clock27}:
\[
\begin{array}{ccccl}
 & (a,L,L,C,R,R) & (a',L',L',C',R',R') & C\mapsto C' & \text{context enters}\\ \hline
p=0 & (\{0,3\},0,0,\{1,3\},1,1) & (\{1,3\},1,1,\{3,2\},2,2) & 1\mapsto3,\ 3\mapsto2 & \text{nowhere}\\
p=1 & (\{1,3\},1,1,\{2,3\},2,2) & (\{2,3\},2,2,\{0,3\},0,0) & 2\mapsto0,\ 3\mapsto3 & \text{at }a\\
p=2 & (\{2,3\},2,2,\{0,3\},0,0) & (\{0,3\},0,0,\{1,3\},1,1) & 0\mapsto1,\ 3\mapsto3 & \text{nowhere}
\end{array}
\]
We justify the entries.
\begin{itemize}
\item $p=0$. Cells $1$ and $2$ have value $0$ and are neighbours of each other with $y=0$, so they are not pinched and become $1$. The centre has $y_2=0$ and $y_4=1$, so it is triggered: $1\mapsto3$, $3\mapsto2$. Cells $4$ and $5$ have value $1$ and left neighbours with $y=1$, so they are not triggered and become $2$. The protector has $y_1=0$, so it is neither pinched nor triggered: $0\mapsto1$, $3\mapsto3$.
\item $p=1$. Cells $1$ and $2$ have value $1$ and left neighbours with $y=1$, so they become $2$. The centre $C=2$ has $y_4=0$, so it is not pinched and becomes $0$; the centre $C=3$ has $y_2=1$, so it is not triggered and stays $3$. Cells $4$ and $5$ have value $2$ and are neighbours of each other with $y=0$, so they become $0$. The protector has $y_1=1$ and is triggered exactly when $y_{-1}=0$, so $a'\in\{2,3\}$ depends on the context.
\item $p=2$. Cells $1$ and $2$ have value $2$ and are neighbours of each other with $y=0$, so they become $0$. The centre $C=0$ has $y_2=0$, so it becomes $1$; the centre $C=3$ has $y_4=0$, so it stays $3$. Cells $4$ and $5$ have value $0$ and are neighbours of each other with $y=0$, so they become $1$. The protector has $y_1=0$: $2\mapsto0$, $3\mapsto3$.
\end{itemize}
In each case $F(c)_{[0,6)}\in W_{27}$ with $p$ replaced by $p+1$. This is (W1), and by Lemma~\ref{hw:exact} $F$ and $F^{-1}$ preserve the occurrences. The cells $1,\dots,5$ of the image depend only on the core, which runs through the $6$-cycle
\[
(0,1)\to(1,3)\to(2,3)\to(0,3)\to(1,2)\to(2,0)\to(0,1),
\]
the A-cycle of Lemma~\ref{hw:59walls}.

\emph{(H3).} The map $\iota$ changes $[C=3]$ and not $p$, so $\beta\circ\iota=\beta+1$. Along the cycle, $\beta$ takes the values $0,1,0,1,0,1$. Since $\beta$ depends only on the core, $\beta(w_0(F(c)))=\beta(w_0(c))+1$ in every context, which is (H3) with $\varepsilon=1$.

\emph{(H1).} The configurations $\iota_0c$ and $c$ differ only at cell $3$, so $F(\iota_0c)$ and $F(c)$ can differ only at cells $2$, $3$, $4$. In the justifications above, the images of cells $2$ and $4$ do not depend on which of the two values $C$ has. For $p=0$, cell $2$ is decided by cell $1$, and cell $4$ reads $C$ only through $y_C=1$. For $p=1$ and $p=2$, cell $2$ is decided by cell $1$ and cell $4$ by cell $5$. The centre map $C\mapsto C'$ is a bijection from $\{p+1,3\}$ onto $\{p+2,3\}$, the centre pair of the image word. So $F(\iota_0c)_3$ is the element of that pair other than $F(c)_3$, and $F(\iota_0c)=\iota_0F(c)$.

\emph{Overlaps.} Suppose $c$ has occurrences at $0$ and $\delta\in\{1,2\}$, with $c_1=p$. If $\delta=1$, then $c_2=p$ and $c_3=C$ are both the left cell of the second word, so $C=p$, which contradicts $C\in\{p+1,3\}$. If $\delta=2$, then $c_3=C$ and $c_4=p+1$ are both the left cell of the second word, so its protector $c_2=p$ must lie in $\{p+1,3\}$, which is impossible. So distinct occurrences are at distance at least $3$. Distances $3$, $4$ and $5$ do occur: at distance $3$ the centre of the first word is the protector of the second.

\emph{(H2) and (H4).} The flip $\iota_0$ writes only cell $3$. It reads only cells $1$ and $3$, since the new value is the element of $\{c_1+1,3\}$ other than $c_3$. The bit $\beta$ reads positions $1$ and $3$ of its word. Let $i\ne0$ be another occurrence, so $|i|\ge3$. The window $[i,i+6)$ contains cell $3$ only if $-2\le i\le3$, hence only if $i=3$. Then cell $3$ is the protector of the word at $3$, whose left cell is $c_4=c_1+1$. So its allowed protectors are $\{c_1+1,3\}$, the pair between which $\iota_0$ exchanges $c_3$. Hence the occurrence at $i$ survives in $\iota_0c$ in every case. Its positions $1$ and $3$, the cells $i+1$ and $i+3$, are not cell $3$, so its bit is unchanged; this is (H4). The flip $\iota_i$ writes cell $i+3$, which is neither $1$ nor $3$ since $i\notin\{-2,0\}$. Symmetrically, $\iota_0$ writes no cell read by $\iota_i$. So the two flips commute, which is (H2).

\emph{Conclusion.} $W_{27}\ne\A^6$. Theorem~\ref{hw:general} applies with $d\ge k=6$ and gives the radius bound $md+\max(3,2)=md+3$. Corollary~\ref{hw:flipall} gives the statement about the flip-all map. The conjugate rules $83$, $202$, $216$ follow as in the proof of Theorem~\ref{hw:59}.
\end{proof}

Machine checks of each step of the proof are listed in Appendix~\ref{app:confirm}. At gap $3$ consecutive walls overlap, as at gap $2$ for rule $59$; Theorem~\ref{hw:general} does not cover this case, and we do not need it.

As for $59$, the protector depends on the context; for instance the image of $1\,1\,1\,2\,2\,2$ is $2\,2\,2\,0\,0\,0$ or $3\,2\,2\,0\,0\,0$ depending on the left neighbour. The flip-all map has radius $3$, and the maps $\Psi_m$ have larger radius. This is consistent with the census of \cite{PaperD}, in which the block maps of radius at most $2$ that commute with $F_{27}$ are the $13$ maps $\sh^aF^b$ and one constant map. A computer search (Appendix~\ref{app:stated}) looked for invariant families with a hidden involution satisfying (W1) and (H1), with windows of length at most $6$, among the $57$ classes with no fixed word of length at most $3$. It considered the closure of each word under $F$ and $F^{-1}$ in all contexts, if that closure has at most $300$ words. It found such families only in the classes of $27$, $59$, $127$, $179$ and $223$, all of which have periodic walls in \cite[Section~4]{Takesue1989}. It says nothing about longer windows or larger families.

\subsection{Rotating walls: rules 219R and 126R}\label{hw:sec:moving}

In the walls of Sections~\ref{hw:sec:general}--\ref{hw:sec:27}, the cells next to the permuted information have context-free trajectories because they belong to a word with that property. For rules $219$ and $126$ the same effect comes from a relation between adjacent cells that $F$ preserves. The cells that satisfy it all evolve in every context by one permutation of the alphabet, of order $4$, so the walls they form stay in place but rotate through the symbols. We call them \emph{rotating walls}. Throughout this subsection $\pi$ is the cell permutation $\pi(x,y)=(y,1-x)$, the $4$-cycle $0\mapsto1\mapsto3\mapsto2\mapsto0$, so that $\pi^2=\kappa$ and $\pi\kappa=\kappa\pi$. The pairs of equal cells of Lemma~\ref{hw:bonds}, which rotate by $\pi$, are the wall (xv) of Takesue for $219$R, and the pairs $b\,\kappa(b)$ used for $126$R below are his wall (xvi), both of period $4$ \cite[Fig.~2, Table~III]{Takesue1989}. The source lists these walls without a proof, and it does not state that $F$ creates no such pair; so Lemma~\ref{hw:bonds} is proved here. The cars, the trains and the automorphisms are what is added.

\begin{lemma}[bonds; cf.\ {\cite[Table~III, wall (xv)]{Takesue1989}}]\label{hw:bonds}
Let $F=F_{219}$, and say that $c$ has a \emph{bond} at $(i,i+1)$ if $c_i=c_{i+1}$.
\begin{enumerate}
\item[(a)] $F(c)$ has a bond at $(i,i+1)$ if and only if $c$ does, and the same holds for $F^{-1}$.
\item[(b)] If $c$ has a bond at $(i-1,i)$ or at $(i,i+1)$, then $F(c)_i=\pi(c_i)$, whatever the other cells of $c$ are.
\end{enumerate}
\end{lemma}

\begin{proof}
The rule $219$ vanishes exactly on $010$ and $101$, so $f(a,b,c)=1$ whenever $a=b$ or $b=c$. If $c_i=c_{i+1}$, then $y_i=y_{i+1}$, so $f=1$ on the triples $(y_{i-1},y_i,y_{i+1})$ and $(y_i,y_{i+1},y_{i+2})$, and by \eqref{eq:F} $F(c)_i=(y_i,1-x_i)=\pi(c_i)$ and $F(c)_{i+1}=\pi(c_{i+1})$. This proves (b) and one half of (a). Conversely, if $F(c)_i=F(c)_{i+1}$, the first coordinates give $y_i=y_{i+1}$, hence $f=1$ at $i$ and at $i+1$ as before, and then the second coordinates give $1-x_i=1-x_{i+1}$; so $c_i=c_{i+1}$. The statement for $F^{-1}$ follows by applying this to $F^{-1}(c)$.
\end{proof}

So every maximal run of at least two equal cells keeps its position and length, and its symbol runs through the $4$-cycle $\pi$ in every context. A \emph{car} of $c$ is a position $p$ with $c_{p-2}=c_{p-1}=c_{p+1}=c_{p+2}\ne c_p$; the common value $s$ of its four neighbours is its \emph{wall symbol}. For $s\in\A$ and $a\ne s$ let $\varphi_s(a)$ be the middle letter of $F(s\,s\,a\,s\,s)$, which by \eqref{eq:F} depends only on $s$ and $a$.

\begin{lemma}\label{hw:cars219}
Let $F=F_{219}$.
\begin{enumerate}
\item[(a)] The cars of $F(c)$ are exactly the cars of $c$. A car $p$ of $c$ with wall symbol $s$ is a car of $F(c)$ with wall symbol $\pi(s)$, and $F(c)_p=\varphi_s(c_p)$; and $\varphi_s$ is a bijection $\A\setminus\{s\}\to\A\setminus\{\pi(s)\}$.
\item[(b)] If $c'$ is obtained from $c$ by changing the values of some cars, each to a value different from its wall symbol, then $c'$ has the same bonds, the same cars and the same wall symbols as $c$.
\end{enumerate}
\end{lemma}

\begin{proof}
(a) Let $p$ be a car of $c$ with wall symbol $s$. The cells $p\pm1$, $p\pm2$ carry bonds, so they become $\pi(s)$ by Lemma~\ref{hw:bonds}(b), and $F(c)_p=\varphi_s(c_p)$ because $F(c)_p$ reads only $c_{p-1},c_p,c_{p+1}$. By Lemma~\ref{hw:bonds}(a) there is no bond at $(p-1,p)$ or $(p,p+1)$ in $F(c)$, so $\varphi_s(c_p)\ne\pi(s)$ and $p$ is a car of $F(c)$. The same argument for $F^{-1}$, under which bonded cells evolve by $\pi^{-1}$, shows that every car of $F(c)$ is a car of $c$, and that the corresponding map $\A\setminus\{\pi(s)\}\to\A\setminus\{s\}$ is inverse to $\varphi_s$. (Explicitly, $\varphi_0=(1\mapsto2,\,2\mapsto0,\,3\mapsto3)$, $\varphi_1=(0\mapsto0,\,2\mapsto1,\,3\mapsto2)$, $\varphi_2=(0\mapsto1,\,1\mapsto2,\,3\mapsto3)$ and $\varphi_3=(0\mapsto0,\,1\mapsto3,\,2\mapsto1)$.)

(b) Let $p$ be a changed car with wall symbol $s$. The bonds at $(p-2,p-1)$ and $(p+1,p+2)$ do not involve $p$, and there is no bond at $(p-1,p)$ or $(p,p+1)$ before or after, since the new value differs from $s$. Whether a position $q$ is a car depends on $c_{[q-2,q+2]}$. For $q\in\{p\pm1,p\pm2\}$ the cell $q$ carries a bond in $c$ and in $c'$, so it is not a car of either; for $|q-p|\ge3$ the window of $q$ does not contain $p$. Wall symbols are cells next to cars, which are not changed.
\end{proof}

Two cars at distance $3$ share the pair of wall cells between them, so they have the same wall symbol. A \emph{train of length $m$} of $c$ is a set $\{p,p+3,\dots,p+3(m-1)\}$ of cars of $c$ such that neither $p-3$ nor $p+3m$ is a car; its \emph{car vector} is $(c_p,c_{p+3},\dots,c_{p+3(m-1)})$. Every car lies in at most one train (it may also lie in an infinite progression of cars, which is not a train). By Lemma~\ref{hw:cars219}(a), $F(c)$ has the same trains as $c$, and $F$ applies the same bijection $\varphi_s$ to every entry of the car vector of a train with wall symbol $s$. By Lemma~\ref{hw:cars219}(b), changing car values within $\A\setminus\{s\}$ does not change the trains.

\begin{theorem}\label{hw:219}
Let $F=F_{219}$. For $\sigma\in S_m$ let $\Psi_m(\sigma)$ replace the car vector $(a_1,\dots,a_m)$ of every train of length $m$ by $(a_{\sigma^{-1}(1)},\dots,a_{\sigma^{-1}(m)})$ and leave every other cell unchanged. Then $\Psi_m(\sigma)$ is a block map of radius at most $3m+2$ that commutes with $F$, and the $\Psi_m$ assemble to an injective homomorphism $\bigoplus_{m\ge1}S_m\to\Aut(X_F)$. In particular $\Aut(X_{F_{219}})$ contains every finite group and is not virtually $\Z^n$ for any $n$.
\end{theorem}

\begin{proof}
\emph{Locality.} Whether $q$ is the $r$-th member of a train of length exactly $m$ depends on whether the positions $q+3t$, $-r\le t\le m-r+1$, are cars, which is decided by $c_{[q-3r-2,\,q+3(m-r+1)+2]}$; the value written at $q$ is one of the cars of that train. So $\Psi_m(\sigma)$ is a block map of radius at most $3m+2$, and it commutes with $\sh$.

\emph{Commutation.} Let $c'=\Psi_m(\sigma)(c)$. By Lemma~\ref{hw:cars219}(b), $c'$ has the bonds, cars, wall symbols and trains of $c$, and it differs from $c$ only at members of trains of length $m$. Let $q$ be the $r$-th member of such a train, with wall symbol $s$ and car vector $a$. Then $F(c')_q=\varphi_s(a_{\sigma^{-1}(r)})$ by Lemma~\ref{hw:cars219}(a), and the same train in $F(c)$ has car vector $(\varphi_s(a_1),\dots,\varphi_s(a_m))$, so $\Psi_m(\sigma)(F(c))_q=\varphi_s(a_{\sigma^{-1}(r)})$ as well. Let $q$ be any other position. Then $\Psi_m(\sigma)(F(c))_q=F(c)_q$, because $F(c)$ has the trains of $c$. If $q$ carries a bond in $c$, it carries one in $c'$, and $F(c')_q=\pi(c'_q)=\pi(c_q)=F(c)_q$ by Lemma~\ref{hw:bonds}(b). If $q$ carries no bond, its neighbours $q\pm1$ are not cars, because a neighbour of a car carries a bond; so $c'$ agrees with $c$ on $[q-1,q+1]$ and $F(c')_q=F(c)_q$. Hence $F\Psi_m(\sigma)=\Psi_m(\sigma)F$.

\emph{Group law.} Since $\Psi_m(\rho)$ preserves trains and replaces each car vector $a$ of length $m$ by $a\circ\rho^{-1}$, $\Psi_m(\sigma)\Psi_m(\rho)$ replaces it by $a\circ\rho^{-1}\circ\sigma^{-1}=a\circ(\sigma\rho)^{-1}$; so $\Psi_m$ is a homomorphism, each $\Psi_m(\sigma)$ is invertible with inverse $\Psi_m(\sigma^{-1})$, and $\Psi_m(\sigma)\in\Aut(X_F)$ by Lemma~\ref{lem:spacetime}. For $m\ne m'$, $\Psi_m(\sigma)$ and $\Psi_{m'}(\rho)$ change disjoint sets of cells and preserve each other's trains and car vectors, so they commute, and the $\Psi_m$ assemble to a homomorphism $\Psi$ of $\bigoplus_mS_m$.

\emph{Injectivity.} Let $g=\Psi((\sigma_{m'})_{m'})$ with $\sigma_m\ne\mathrm{id}$, and choose $r$ with $\sigma_m(r)\ne r$. Let $c$ be the configuration with $c_{3t}=1$ for $t=r$, $c_{3t}=2$ for $t\in\{1,\dots,m\}\setminus\{r\}$, and $c_q=0$ at every other position $q$. Every cell $0$ of $c$ carries a bond, so the cars of $c$ are exactly $3,6,\dots,3m$, which form one train of length $m$ with wall symbol $0$. Then $g(c)$ has the value $1$ at $3\sigma_m(r)\ne3r$, so $g(c)\ne c$. The last statement follows as in the proof of Theorem~\ref{thm:54}.
\end{proof}

Neither Theorem~\ref{thm:walls} nor Corollary~\ref{hw:periodic} applies to $219$R as stated. The walls are not fixed, since their symbol rotates, and the words $s\,s\,a\,s\,s$ with $a\ne s$, whose images are context-free by Lemma~\ref{hw:bonds}, have orbits of period $12$ under $F$ in which no word differs from the word half a period later at the car alone.

\begin{proposition}\label{hw:219D}
Let $F=F_{219}$, and for $s\in\A$ let $R_s=\varphi_{\pi^3(s)}\circ\varphi_{\pi^2(s)}\circ\varphi_{\pi(s)}\circ\varphi_s$, the permutation of $\A\setminus\{s\}$ by which $F^4$ acts on a car with wall symbol $s$. Let $D$ be the map that replaces the value $a$ of every car with wall symbol $s$ by $R_s(a)$ and leaves every other cell unchanged. Then $D$ is an automorphism of radius $2$ commuting with $F$, $D^3=\mathrm{id}\ne D$, and $D\notin\langle\sh,F\rangle\times\langle\kappa\rangle$.
\end{proposition}

\begin{proof}
$D$ has radius $2$, since being a car and the wall symbol are read from $c_{[p-2,p+2]}$. Since $\pi^4=\mathrm{id}$, $R_{\pi(s)}\circ\varphi_s=\varphi_s\circ R_s$, so $F$ intertwines the action of $D$ on a car with wall symbol $s$ and on the same car, with wall symbol $\pi(s)$, in $F(c)$; at all other cells the commutation argument in the proof of Theorem~\ref{hw:219} applies verbatim, with Lemma~\ref{hw:cars219}(b). A direct computation gives $R_0=(1\mapsto3\mapsto2\mapsto1)$, $R_1=(0\mapsto2\mapsto3\mapsto0)$, $R_2=(0\mapsto3\mapsto1\mapsto0)$ and $R_3=(0\mapsto1\mapsto2\mapsto0)$, all $3$-cycles; so $D^3=\mathrm{id}\ne D$, and $D$ is invertible. By Proposition~\ref{cl:prop:kappa}, $\langle\sh,F\rangle\times\langle\kappa\rangle\cong\Z^2\times\Z/2$, which has no element of order $3$.
\end{proof}

The map $D$ is the analogue of the relabelling $D$ of $54$R. A variant of $\Psi_2$ with trains bounded by runs of at least three wall cells, which acts on the windows $s\,s\,s\,a\,s\,s\,b\,s\,s\,s$ ($a,b\ne s$) by exchanging $a$ and $b$, has radius $6$; it commutes with $F$ on all $4^{15}$ windows and is an involution (computer-assisted; Appendix~\ref{app:stated}).

For rule $126$ the relation is $c_{i+1}=\kappa(c_i)$, and the walls are runs that alternate between a symbol and its complement. The two rules are related by the correspondence of \cite[(A.5)--(A.7)]{Takesue1989}: the cellwise map $\theta$ that complements the $y$-bit of the cells at odd sites and the $x$-bit of the cells at even sites satisfies $\theta F_{126}=F_{219}\,\sh\theta\sh^{-1}$, and it maps the pairs $b\,\kappa(b)$ to pairs of equal cells. It does not commute with $\sh$, so it does not identify the two automorphism groups, and the two rules are treated separately. Here Theorem~\ref{hw:general} applies verbatim.

\begin{proposition}\label{hw:126}
Let $F=F_{126}$ and
\[
W=\{\,b\;\kappa(b)\;q\;b\;\kappa(b)\ :\ b\in\A,\ q\in\A\setminus\{b,\kappa(b)\}\,\}\subsetneq\A^5,
\]
with $j=2$ and $\iota$ replacing the letter $q$ at position $2$ by $\kappa(q)$. Then \textup{(W1)}, \textup{(H1)}, \textup{(H2)}, \textup{(H4)} and \textup{(H3$'$)} hold, with the class of an occurrence $w=b\,\kappa(b)\,q\,b\,\kappa(b)$ equal to $b\in K=\A$, $\theta=\pi$, bit $\beta(w)=[q=\pi(b)]$ and $\varepsilon\equiv0$. Hence $\bigoplus_{m\ge1}S_m$ embeds in $\Aut(X_F)$. The flip-all map $E=\tau_W$ of Corollary~\ref{hw:flipall}, which replaces $c_p$ by $\kappa(c_p)$ whenever $c_{[p-2,p+2]}\in W$, is an involutive automorphism of radius $2$, and $E\notin\langle\sh,F\rangle\times\langle\kappa\rangle$.
\end{proposition}

\begin{proof}
The rule $126$ vanishes exactly on $000$ and $111$. In every cell of $w\in W$ the $y$-bit differs from that of a neighbour inside $w$: for the cells $b$ and $\kappa(b)$ because $y(\kappa(b))=1-y(b)$, and for $q$ because its two neighbours $\kappa(b)$ and $b$ have different $y$-bits. So $f=1$ at every cell of an occurrence, whatever the outside neighbours, and $F$ maps the occurrence $w$ to the cellwise image $\pi(w)=\pi(b)\,\kappa(\pi(b))\,\pi(q)\,\pi(b)\,\kappa(\pi(b))$, which lies in $W$ because $\pi$ commutes with $\kappa$. This is (W1), and it also gives the class map $\theta=\pi$. For (H1), $\iota_0c$ again has an occurrence at $0$, which $F$ maps to $\pi(\iota(w))=\iota(\pi(w))$; outside $[0,5)$ the configurations $F(\iota_0c)$ and $F(c)$ agree, since $\iota_0$ changes only cell $2$. For (H3$'$), $q\in\A\setminus\{b,\pi^2(b)\}=\{\pi(b),\pi^3(b)\}$, and $\kappa=\pi^2$ exchanges these two values, so $\beta\circ\iota=\beta+1$; and $[\pi(q)=\pi(\pi(b))]=[q=\pi(b)]$, so $\varepsilon\equiv0$. Two occurrences cannot lie at distance $1$ or $2$: at distance $1$ the second would force $q=b$, and at distance $2$ it would force $q=\kappa(b)$. As in the proof of Corollary~\ref{hw:periodic}, this gives (H2) and (H4). Theorem~\ref{hw:general} and Corollary~\ref{hw:flipall} apply. Finally, $E\ne\mathrm{id}$, and $E\ne\kappa$ since $E$ fixes $(0,0)^{\Z}$, which contains no occurrence. By Proposition~\ref{cl:prop:kappa}, $\langle\sh,F\rangle\times\langle\kappa\rangle\cong\Z^2\times\Z/2$, whose only element of order $2$ is $\kappa$.
\end{proof}

\begin{sloppypar}
The trains of Theorem~\ref{hw:general} here consist of occurrences $b\,\kappa(b)\,q\,b\,\kappa(b)$ in the same phase $b$, and the walls between their cars are alternating runs, which $F$ rotates by $\pi$ in every context. In both rules the commuting block maps of radius $2$ are too many to enumerate, so the census of \cite{PaperD} does not list these automorphisms of radius $2$.
\end{sloppypar}

\subsection{Blocking words}\label{hw:sec:kurka}

\begin{remark}\label{hw:kurka}
For a cellular automaton of radius $1$, a word $u$ is \emph{blocking} \cite{BT2000} if there is a position $q$ in $u$ such that $F^n(x)_q=F^n(x')_q$ for all $n\ge0$ and all $x,x'$ that contain $u$ at the same place; such a column separates the dynamics on its two sides, and by the argument in the proof of \cite[Theorem~3]{Kurka1997} a cellular automaton has a blocking word if and only if it has equicontinuity points, that is, if and only if it is not sensitive \cite[Corollary~2]{Kurka1997}. For the second-order elementary automata Takesue observed that the two sides of a wall cannot communicate and that a system with walls is not mixing \cite[Section~4]{Takesue1989}. Every word of $X_A\cup X_B$ is blocking for $F_{59}$, with $q$ the left cell (or the right cell), and also for $F_{59}^{-1}$, by Lemma~\ref{hw:59walls}. There are even blocking words of length $3$: the six words $a\,p\,p$ with $a\in\{p,3\}$, which are the walls of \cite[(4.3)]{Takesue1989} on three cells, form a family invariant under $F_{59}$ and $F_{59}^{-1}$ in which the last two cells have context-free images $p+1$, $p+1$, as follows from \eqref{hw:clock} (each of the two in-phase cells has its left neighbour in the same phase, hence is never triggered). The time-periodic walls of Proposition~\ref{hw:tp} are blocking as well, since their cells have context-free trajectories, and so are the words of $W_{27}$ for $F_{27}$, whose cells $1,\dots,5$ have context-free trajectories by the proof of Theorem~\ref{hw:27}. Words fixed in every context are blocking, but the converse fails, and the rules of this section show that the gap matters for automorphisms: rules $27$, $59$, $127$, $179$ and $223$ have no word of length at most $3$ fixed in every context (Table~\ref{tab:walls}), and rule $59$ has none of any length (Lemma~\ref{hw:nofixed}), yet they have blocking words and their automorphism groups contain every finite group. This is why Question~\ref{q:dichotomy} is phrased in terms of blocking words rather than fixed words. In the nonlinear examples of this paper, the automorphisms outside $\langle\sh,F\rangle$ all come from blocking words with a degree of freedom next to the blocking column: cars between fixed walls in Theorem~\ref{thm:walls}, a hidden bit inside a moving wall in Theorems~\ref{hw:general} and~\ref{hw:59}, and cars between rotating walls in Theorem~\ref{hw:219} and Proposition~\ref{hw:126}. For the affine rules $90$ and $165$ this fails: they have no blocking word, yet their automorphism groups contain every finite group (Corollary~\ref{lin:noblock}).
\end{remark}

\section{The linear rules 90 and 150}\label{sec:linear}

For a linear rule $f(a,b,c)=\alpha a\oplus\beta b\oplus\gamma c$ the automaton $F_f$ is $\F_2$-linear, and in terms of the ring $\Lam=\F_2[s,s^{-1}]$ acting on $\F_2^{\Z}$ by $s=\sh$ it is the matrix
\[
M=\begin{pmatrix}0&1\\1&p\end{pmatrix},\qquad p=\alpha s^{-1}+\beta+\gamma s,
\]
acting on column vectors $(x,y)^{\mathsf T}$; indeed $f(y_{i-1},y_i,y_{i+1})=(py)_i$. Rule $90$ has $p=s+s^{-1}$ and rule $150$ has $p=s+1+s^{-1}$. The characteristic polynomial of $M$ is $t^2+pt+1$.

\begin{lemma}\label{lem:commutant}
Let $K=\F_2(s)$. If $t^2+pt+1$ has distinct roots in $K$, or is irreducible over $K$, then the matrices in $M_2(\Lam)$ commuting with $M$ are exactly $aI+bM$ with $a,b\in\Lam$. Suppose moreover that $p=s^{-1}+\beta+s$ with $\beta\in\F_2$, as for rules $90$ and $150$, and let $R\ge1$. Then the matrices $aI+bM$ whose entries all lie in the span of $s^{-R},\dots,s^{R}$ are those with $a\in\operatorname{span}(s^{-R},\dots,s^R)$ and $b\in\operatorname{span}(s^{-R+1},\dots,s^{R-1})$, an $\F_2$-space of dimension $4R$. For $R=0$ they are $0$ and $I$.
\end{lemma}

\begin{proof}
In both cases $M$ is not scalar and its minimal polynomial over $K$ has degree $2$, so its commutant in $M_2(K)$ is $K[M]=\{aI+bM:a,b\in K\}$. If $aI+bM=\begin{psmallmatrix}a&b\\b&a+bp\end{psmallmatrix}$ has entries in $\Lam$ then $a,b\in\Lam$. For the radius count, the entries $a$ and $b$ must lie in the span of $s^{-R},\dots,s^R$, and then $a+bp$ lies in that span if and only if $bp$ does. Since $p$ has the terms $s$ and $s^{-1}$, the terms of $bp$ of degree $R+1$ and $-(R+1)$ are those of $b$ of degree $R$ and $-R$ multiplied by $s$ and $s^{-1}$, so the condition is that $b$ has no term $s^{\pm R}$. For $R\ge1$ this leaves $2R-1$ coefficients of $b$ and $2R+1$ of $a$, all free; for $R=0$ it forces $b=0$, and $a\in\{0,1\}$.
\end{proof}

For rule $90$, $t^2+(s+s^{-1})t+1=(t+s)(t+s^{-1})$ has the distinct roots $s^{\pm1}$. For rule $150$ the polynomial is irreducible over $K$: a root would be integral over $\Lam$, hence in $\Lam$, and the product of the two roots being $1$ makes it a unit $s^k$, but then $s^k+s^{-k}=s+1+s^{-1}$ is impossible.

For the linear rules the centralizer is a vector space and can be computed by Gaussian elimination over $\F_2$.

\begin{theorem}\label{thm:90}
Let $F=F_{90}$, let $\Lam=\F_2[s,s^{-1}]$ act on $\F_2^{\Z}$ through $s=\sh$, and let $M=\begin{psmallmatrix}0&1\\1&p\end{psmallmatrix}$ with $p=s+s^{-1}$, so that $F(x,y)=M(x,y)^{\mathsf T}$.
\begin{enumerate}
\item[(a)] The affine block maps commuting with $F$ are the maps $c\mapsto (aI+bM)c+k$ with $a,b\in\Lam$ and $k\in\{(0,0)^{\Z},(1,1)^{\Z}\}$. For $R\ge1$, those of radius at most $R$ form an $\F_2$-space of dimension $4R+1$; those of radius $0$ are the four maps $c\mapsto ac+k$ with $a\in\{0,1\}$.
\item[(b)] The maps $u=x+s^{-1}y$ and $w=x+sy$ are factor maps onto the binary full shift with $u\circ F=\sh^{-1}\circ u$ and $w\circ F=\sh\circ w$; the maps $\iota(v)=(v,s^{-1}v)$ and $\iota'(v)=(v,sv)$ are embeddings with $F\circ\iota=\iota\circ\sh^{-1}$ and $F\circ\iota'=\iota'\circ\sh$. For every block map $h$ of the binary full shift, $\iota\circ h\circ u$ and $\iota'\circ h\circ w$ commute with $F$.
\item[(c)] For every $R\ge1$ the $\F_2$-span of the affine commuting maps of radius at most $R$ and of the maps in (b) of radius at most $R$ has dimension exactly $2^{2R}+1$. For $R\le3$ this span is the whole space of block maps of radius at most $R$ commuting with $F$, so that space has dimension $5$, $17$ and $65$ for $R=1,2,3$.
\end{enumerate}
\end{theorem}

\begin{proof}
(a) An affine block map has the form $c\mapsto Nc+k$ with $N$ a linear block map and $k=g(0)$ a constant configuration, $k\in\{0^{\Z},1^{\Z},2^{\Z},3^{\Z}\}$. It commutes with $F=M$ if and only if $NM=MN$ and $Mk=k$. For a constant $k=(x,y)^{\Z}$ one has $F(k)=(y,x\oplus f(y,y,y))^{\Z}=(y,x)^{\Z}$ since $f(y,y,y)=0$ for rule $90$, so $k$ is fixed if and only if $x=y$, that is, $k\in\{(0,0)^{\Z},(1,1)^{\Z}\}$. Combined with Lemma~\ref{lem:commutant} this gives dimension $4R+1$ for $R\ge1$, and for $R=0$ the linear parts $0$ and $I$ with the two constants.

(b) The row vector $(1,s^{-1})$ satisfies $(1,s^{-1})M=(s^{-1},1+s^{-1}p)=(s^{-1},s^{-2})=s^{-1}(1,s^{-1})$, which is $u\circ F=\sh^{-1}\circ u$ for $u(x,y)=x+s^{-1}y$; similarly $(1,s)M=s(1,s)$. The column vector $(1,s^{-1})^{\mathsf T}$ satisfies $M(1,s^{-1})^{\mathsf T}=(s^{-1},1+ps^{-1})^{\mathsf T}=s^{-1}(1,s^{-1})^{\mathsf T}$, which is $F\circ\iota=\iota\circ\sh^{-1}$, and similarly for $\iota'$. Both $u$ and $w$ are surjective onto $\F_2^{\Z}$ (choose $y=0$). For a block map $h$ of the binary full shift, which commutes with $\sh^{\pm1}$,
\[
(\iota h u)F=\iota h\sh^{-1}u=\iota\sh^{-1}hu=F\iota hu,
\]
and likewise for $\iota'hw$.

(c) Fix $R\ge1$ and let $A_R$ be the space in (a), $U_R$ the span of the maps $\iota hu$ with $h$ a polynomial in the variables $u_j$, $j\in J_U=[-(R-2),R]$, and $W_R$ the span of the $\iota'hw$ with $h$ a polynomial in $w_k$, $k\in J_W=[-R,R-2]$; the index ranges are those for which $\iota hu$ and $\iota'hw$ have radius at most $R$, since $(\iota hu)(c)_i=(h(u)_i,h(u)_{i-1})$ and $u_{i+j}=x_{i+j}\oplus y_{i+j-1}$. The linear forms $u_j$ ($j\in J_U$) and $w_k$ ($k\in J_W$) on the $2(2R+1)$ coordinates of a window are jointly linearly independent: a relation $\sum a_ju_j+\sum b_kw_k=0$ gives $a_m=b_m$ from the coefficient of $x_m$ and $a_{m+1}=b_{m-1}$ from that of $y_m$, with $a,b$ extended by zero outside their ranges; so $a_{m+1}=a_{m-1}$ for all $m\in\Z$. Since $R\notin J_W$ and $R-1\notin J_W$, we have $b_R=b_{R-1}=0$, hence $a_R=b_R=0$ and $a_{R-1}=b_{R-1}=0$; descending with $a_{m-1}=a_{m+1}$ gives $a_m=0$ for all $m\le R$, hence for all $m$, and then $b_m=a_m=0$ for all $m$. Consequently distinct monomials in the $u_j$ give linearly independent functions, $\dim U_R=2^{|J_U|}=2^{2R-1}$, likewise $\dim W_R=2^{2R-1}$, and a polynomial in the $u$'s equals a polynomial in the $w$'s only if both are constant, so $U_R\cap W_R$ is spanned by the constant map onto $(1,1)^{\Z}$ and $\dim(U_R+W_R)=2^{2R}-1$. An element $\iota hu+\iota'h'w$ is affine if and only if $h$ and $h'$ are affine, again by independence of the variables, so $A_R\cap(U_R+W_R)=(A_R\cap U_R)+(A_R\cap W_R)$, where $A_R\cap U_R$ consists of the $\iota hu$ with $h$ affine and has dimension $|J_U|+1=2R$, likewise $\dim(A_R\cap W_R)=2R$, and the two meet in the constants. Therefore
\[
\dim(A_R+U_R+W_R)=(4R+1)+(2^{2R}-1)-(4R-1)=2^{2R}+1 .
\]
This is a lower bound for $\dim\Cent_R(F)$, since all three families commute with $F$. The Gaussian elimination of Section~\ref{sec:method} gives $\dim\Cent_R(F)=5,17,65$ for $R=1,2,3$, and a direct rank computation on the tables of the three families, each table re-verified to commute, gives the same numbers, so the span is everything for $R\le3$.
\end{proof}

\begin{conjecture}
$\dim_{\F_2}\Cent_R(F_{90})=2^{2R}+1$ for every $R\ge1$; equivalently, every block map commuting with $F_{90}$ is a sum of an affine commuting map, a map $\iota hu$ and a map $\iota'h'w$.
\end{conjecture}

The invertible elements are few. At $R=1$ the $32$ commuting maps contain $10$ invertible ones, all affine, namely the $5$ maps $\sh^aF^b$ and the $5$ further maps $(aI+bM)c+k$ whose determinant $a^2+abp+b^2$ is a unit of $\Lam$. The maps $\iota hu$ and $\iota'hw$ are never injective, since they factor through $u$ or $w$. (The maps $\mathrm{id}+\iota hu$ can be injective; see Section~\ref{lin:sec:movers}.) At $R=2$ the enumeration of all $2^{17}$ commuting maps, with the pair-graph test applied to the $98{,}222$ balanced ones, finds exactly $26$ invertible maps, all affine: the $13$ maps $\sh^aF^b$ and the $13$ further maps $(aI+bM)c+k$ with unit determinant. So every automorphism of $X_{F_{90}}$ of radius at most $2$ is affine. This does not persist: by Theorem~\ref{lin:movers}, $X_{F_{90}}$ has non-affine automorphisms, and its automorphism group contains every finite group.

\begin{proposition}\label{prop:150}
For $1\le R\le3$, $\Cent_R(F_{150})$ consists exactly of the $2^{4R}$ linear maps $(aI+bM)c$ of Lemma~\ref{lem:commutant}.
\end{proposition}

\begin{proof}
Lemma~\ref{lem:commutant} gives $2^{4R}$ linear commuting maps, and the only constant configuration fixed by $F_{150}$ is $(0,0)^{\Z}$: since $f(y,y,y)=y$, $F_{150}$ maps $(x,y)^{\Z}$ to $(y,x\oplus y)^{\Z}$, which equals $(x,y)^{\Z}$ only for $x=y=0$. So there is no further constant. Gaussian elimination gives $\dim\Cent_R(F_{150})=4,8,12$ for $R=1,2,3$.
\end{proof}

So the non-affine commuting maps of $F_{90}$ are explained entirely by the two eigen-factors onto the binary shift, and the rule without eigen-factors has none, at least through radius $3$. Both results also confirm the encoding: for the linear rules the SAT enumeration at $R\le2$ and the elimination agree, and the $2^{2R}+1$ and $4R$ patterns are exactly what the algebra predicts.

\begin{table}[t]
\caption{Block maps of radius at most $R$ commuting with $F_{90}$ and $F_{150}$. ``Named'' are the maps $\sh^aF^b$ with $|a|+|b|\le R$; ``other'' are maps that are neither named nor constant; the last two columns count the invertible commuting maps, that is, the automorphisms of radius at most $R$, and those outside $\langle\sh,F\rangle$; for $150$R, $h=\mathrm{id}+\sh F$ as in Proposition~\ref{prop:150Z3}.}
\label{tab:census}
\centering
\begin{tabular}{llrrrrrl}
\toprule
rule & $R$ & commuting & named & constants & other & invertible & outside $\langle\sh,F\rangle$\\
\midrule
90R & 1 & $2^{5}$ & 5 & 2 & 25 & 10 & $5$, all affine\\
90R & 2 & $2^{17}$ & 13 & 2 & rest & 26 & $13$, all affine\\
90R & 3 & $2^{65}$ & 25 & 2 & rest & --- & ---\\
150R & 1 & $2^{4}$ & 5 & 1 & 10 & 9 & $4$, all in $\langle\sh,F,h\rangle$\\
150R & 2 & $2^{8}$ & 13 & 1 & rest & 43 & $30$, all in $\langle\sh,F,h\rangle$\\
150R & 3 & $2^{12}$ & 25 & 1 & rest & 115 & $90$, all in $\langle\sh,F,h\rangle$\\
\bottomrule
\end{tabular}
\end{table}

The commuting maps of $F_{150}$ are all linear through radius $3$, but its automorphism group is still larger than $\langle\sh,F_{150}\rangle$. In $M_2(\Lam)$ the shift is the scalar matrix $sI$ and $F_{150}$ is $M$; under the ring isomorphism $\Lam[M]\cong\Lam[t]/(t^2+pt+1)$ the commutant of Lemma~\ref{lem:commutant} becomes a subring of the field $L=K[t]/(t^2+pt+1)$, a quadratic extension of $K=\F_2(s)$, and composition of commuting linear maps becomes multiplication in $L$.

\begin{proposition}\label{prop:150Z3}
Let $F=F_{150}$ and $h=\mathrm{id}+\sh F$. Then $h\in\Aut(X_F)$, $h\notin\langle\sh,F\rangle$, and $(a,b,e)\mapsto\sh^aF^bh^e$ is an isomorphism $\Z^3\to\langle\sh,F,h\rangle$. Moreover $F$ has no blocking word. In particular $\Aut(X_{F_{150}})$ is not virtually $\Z^2$.
\end{proposition}

\begin{proof}
The map $h$ has matrix $I+sM$, so it is a block map commuting with $F$ and $\sh$. It is invertible, because $\det(I+sM)=1+sp+s^2=1+s(s+1+s^{-1})+s^2=s$ is a unit of $\Lam$, so the inverse matrix has entries in $\Lam$. By Lemma~\ref{lem:spacetime}, $h\in\Aut(X_F)$.

For independence, identify $\sh^aF^bh^e$ with the element $s^at^b(1+st)^e$ of $L^\times$, where $t$ is the image of $M$; the identification is an injective homomorphism on the commuting linear invertible maps, since $I$ and $M$ are linearly independent over $\Lam$. Suppose $s^at^b(1+st)^e=1$. Let $N\colon L^\times\to K^\times$ be the norm, which is the determinant: $N(s)=s^2$, $N(t)=\det M=1$ and $N(1+st)=s$. So $s^{2a+e}=1$, and $e=-2a$. Put $u=(1+st)^2s^{-1}$; the relation becomes $t^b=u^a$. Since $t^2=pt+1$, we have $t^{-1}=t+p$, and a direct computation with $sp=s^2+s+1$ and $(s^2+s+1)p=s^3+s+s^{-1}$ gives
\[
u=s^{-1}+s+(s^2+s+1)t=s^3+(s^2+s+1)t^{-1},\qquad ut=s^3t+s^2+s+1 .
\]
Let $v$ be a valuation of $L$ extending a valuation of $K$ with $v(s)=e_0>0$. Then $v(p)=-e_0$, and since $t\cdot t^{-1}=1$ and $t+t^{-1}=p$, one of $v(t)$, $v(t^{-1})$ is $-e_0$ and the other $e_0$. If $v(t)=-e_0$, the second expression gives $v(u)=\min(3e_0,e_0)=e_0$; if $v(t)=e_0$, the first gives $v(u)=-e_0$, since $v(s^{-1})=-e_0$ and the other terms have positive valuation. In both cases $v(u)=-v(t)$, and $t^b=u^a$ gives $b=-a$, so $(ut)^a=1$. The element $ut$ is not a root of unity: a root of unity in $L$ is algebraic over $\F_2$, so it lies in the algebraic closure of $\F_2$ in $L$, a finite field $\F_q$ with $[\F_q:\F_2]\le[L:K]=2$; every element of $\F_4$ has trace over $K$ in $\F_2$, while the trace of $ut$ is $s^3(t+t^{-1})=s^3p=s^4+s^3+s^2$. Hence $a=0$, and then $b=e=0$. So the homomorphism $\Z^3\to\langle\sh,F,h\rangle$ is injective. In particular $h\notin\langle\sh,F\rangle$.

For the blocking words, let $u$ be a word, $q$ a position in it, and $x$ a configuration containing $u$ on $[0,|u|)$. For $N\ge|u|$ let $x'$ be $x$ with the $y$-bit of the cell $N$ complemented, and let $\Delta=Y_x\oplus Y_{x'}$. By linearity $\Delta$ satisfies $\Delta(t+1,i)=\Delta(t-1,i)\oplus\Delta(t,i-1)\oplus\Delta(t,i)\oplus\Delta(t,i+1)$, with $\Delta(-1,\cdot)=0$ and $\Delta(0,\cdot)$ the unit vector at $N$. By induction on $t\ge0$, $\Delta(t,i)=0$ for $i<N-t$ and $\Delta(t,N-t)=1$: at $i<N-t-1$ all four terms vanish, and at $i=N-t-1$ only $\Delta(t,N-t)=1$ is nonzero. For $t=N-q$ this gives $F^{N-q}(x)_q\ne F^{N-q}(x')_q$, so $u$ is not blocking with column $q$.

Finally, if $\Aut(X_F)$ had a finite-index subgroup isomorphic to $\Z^2$, it would meet $\langle\sh,F,h\rangle\cong\Z^3$ in a subgroup of finite index in $\Z^3$, hence of rank $3$, inside $\Z^2$; this is impossible.
\end{proof}

In fact $\Aut(X_{F_{150}})=\langle\sh,F,h\rangle$: every automorphism of $X_{F_{150}}$ is affine (Lemma~\ref{aff:lem:R}), so it is a map $aI+bM$ whose determinant $a^2+abp+b^2$ is a unit of $\Lam$, and these maps are exactly the elements $s^at^b(1+st)^e$ (Lemma~\ref{aff:lem:units} and Theorem~\ref{thm:affine}(4)).

\begin{remark}\label{lin:rem:window}
For $M\ge0$ let $\Lambda_M$ be the light-speed window graph of $F$: its vertices are the words $\omega\in\A^{M+1}$, and for every configuration $\tilde c$ with $\tilde c_{[p,p+M]}=\omega$ it has an edge from $\omega$ to the word $F(\tilde c)_{[p+1,p+M+1]}$, which depends only on $\omega$ and on $\tilde c_{p+M+1},\tilde c_{p+M+2}$. Its strong connectivity for every $M$ is one of the hypotheses of the rigidity theorems of \cite{PaperB}. A direct construction of $\Lambda_M$ from the local rule (Appendix~\ref{app:stated}) finds it strongly connected for $M\le4$ for the affine rules $15$, $60$, $105$, $150$, $195$ and $240$, and not strongly connected for $90$ and $165$. So for $150$ the extra automorphisms of Proposition~\ref{prop:150Z3} are not detected by these graphs; they come from linearity.
\end{remark}

\subsection{Every finite group for 90R and 165R}\label{lin:sec:movers}

The maps $\iota hu$ of Theorem~\ref{thm:90}(b) are never invertible, but a map $\mathrm{id}+\iota hu$ can be. Here we show that such maps carry two commuting copies of $\bigoplus_mS_m$, for rule $90$ and for its complement $165$.

Throughout this subsection $F$ is $F_{90}$ or $F_{165}$, and $\varepsilon=0$ or $1$ accordingly. Since $165(a,b,c)=1\oplus a\oplus c$, we have $F_{165}(c)=F_{90}(c)+\varepsilon(0,1)^{\Z}$, with cellwise addition in $\F_2^2$. Binary configurations $v\in\F_2^{\Z}$ are added cellwise, $\mathbf 1=1^{\Z}$, $\bar v=v+\mathbf 1$, and $(\sh v)_n=v_{n+1}$. Define
\[
d(c)_n=x_{n+1}+y_{n+2},\qquad e(c)_n=y_n+x_{n+1},\qquad \iota(v)_i=(v_{i-1},v_i),\qquad \iota'(v)_i=(v_{i+1},v_i);
\]
thus $\iota$ is the embedding of Theorem~\ref{thm:90}(b), and $d=\sh w$, $e=\sh u$ for its factor maps $w$ and $u$.

\begin{lemma}\label{lin:coords}
\begin{enumerate}
\item[(a)] $d$ and $e$ are $\F_2$-linear block maps $\A^{\Z}\to\F_2^{\Z}$ commuting with $\sh$, $d$ reads the cells $n+1,n+2$ and $e$ the cells $n,n+1$, and both are surjective.
\item[(b)] $d(F(c))=\sh d(c)+\varepsilon\mathbf 1$ and $e(F(c))=\sh^{-1}e(c)+\varepsilon\mathbf 1$.
\item[(c)] $F_{90}\circ\iota=\iota\circ\sh$ and $F_{90}\circ\iota'=\iota'\circ\sh^{-1}$.
\item[(d)] $d\circ\iota=\mathrm{id}+\sh^2$, $e\circ\iota=0$, $d\circ\iota'=0$ and $e\circ\iota'=\mathrm{id}+\sh^2$.
\end{enumerate}
\end{lemma}

\begin{proof}
(a) is clear; for surjectivity take $y=0$. Write $F(c)_i=(y_i,\,x_i+y_{i-1}+y_{i+1}+\varepsilon)$. Then $d(F(c))_n=y_{n+1}+x_{n+2}+y_{n+1}+y_{n+3}+\varepsilon=d(c)_{n+1}+\varepsilon$, and $e(F(c))_n=x_n+y_{n-1}+y_{n+1}+\varepsilon+y_{n+1}=e(c)_{n-1}+\varepsilon$, which is (b). For (c), $F_{90}(\iota v)_i=(v_i,\,v_{i-1}+v_{i-1}+v_{i+1})=(v_i,v_{i+1})=\iota(\sh v)_i$ and $F_{90}(\iota'v)_i=(v_i,\,v_{i+1}+v_{i-1}+v_{i+1})=(v_i,v_{i-1})=\iota'(\sh^{-1}v)_i$. For (d), $d(\iota v)_n=v_n+v_{n+2}$, $e(\iota v)_n=v_n+v_n=0$, $d(\iota'v)_n=v_{n+2}+v_{n+2}=0$ and $e(\iota'v)_n=v_n+v_{n+2}$.
\end{proof}

\begin{definition}\label{lin:def:integral}
Let $\mathcal G$ be the set of automorphisms $\psi$ of the binary full shift with $\psi(0^{\Z})=0^{\Z}$ for which there is a block map $J\colon\F_2^{\Z}\to\F_2^{\Z}$ with $J(0^{\Z})=0^{\Z}$ and
\[
\psi(z)=z+J(z)+\sh^2J(z)\qquad(z\in\F_2^{\Z}),
\]
that is, $\psi=\mathrm{id}+(\mathrm{id}+\sh^2)J$. We call $J=J_\psi$ the \emph{integral} of $\psi$. Let $\mathcal G^{\kappa}\subseteq\mathcal G$ consist of the $\psi$ with $J_\psi(\bar z)=J_\psi(z)$ for all $z$; such $\psi$ satisfy $\psi(\bar z)=\overline{\psi(z)}$.
\end{definition}

\begin{lemma}\label{lin:integral}
The integral is unique. $\mathcal G$ is a group, with $J_{\psi\chi}=J_\chi+J_\psi\circ\chi$ and $J_{\psi^{-1}}=J_\psi\circ\psi^{-1}$, and $\mathcal G^{\kappa}$ is a subgroup.
\end{lemma}

\begin{proof}
Let $K$ be a block map of radius $r$ with $K(0^{\Z})=0^{\Z}$ and $(\mathrm{id}+\sh^2)K=0$, so that every $K(z)$ is invariant under $\sh^2$. Given $z$ and $n$, let $z'$ agree with $z$ on $[n-r,n+r]$ and vanish on $[n+2N-r,n+2N+r]$ for some large $N$. Then $K(z)_n=K(z')_n=K(z')_{n+2N}=K(0^{\Z})_n=0$. So $K=0$, which gives uniqueness. For $\psi,\chi\in\mathcal G$, $(\mathrm{id}+\sh^2)(J_\chi+J_\psi\circ\chi)=(\chi+\mathrm{id})+(\psi\chi+\chi)=\psi\chi+\mathrm{id}$, and $J_\chi+J_\psi\circ\chi$ is a block map vanishing at $0^{\Z}$; so $\psi\chi\in\mathcal G$ with the stated integral. Similarly $(\mathrm{id}+\sh^2)(J_\psi\circ\psi^{-1})=(\psi+\mathrm{id})\circ\psi^{-1}=\mathrm{id}+\psi^{-1}$, so $\psi^{-1}\in\mathcal G$. If $J_\psi$ and $J_\chi$ are invariant under $z\mapsto\bar z$, then $\psi$ and $\chi$ commute with it, since $\psi(\bar z)=\bar z+(\mathrm{id}+\sh^2)J_\psi(z)=\overline{\psi(z)}$, and hence so are $J_\chi+J_\psi\circ\chi$ and $J_\psi\circ\psi^{-1}$.
\end{proof}

\begin{theorem}\label{lin:movers}
For $\psi\in\mathcal G$ define
\[
\Phi_\psi(c)=c+\iota\bigl(J_\psi(d(c))\bigr),\qquad \Phi'_\psi(c)=c+\iota'\bigl(J_\psi(e(c))\bigr).
\]
Let $H=\mathcal G$ if $F=F_{90}$ and $H=\mathcal G^{\kappa}$ if $F=F_{165}$. Then $(\psi,\chi)\mapsto\Phi_\psi\Phi'_\chi$ is an injective homomorphism $H\times H\to\Aut(X_F)$. The group $\bigoplus_{m\ge1}S_m$ embeds in $\mathcal G^{\kappa}$ (Lemma~\ref{lin:markers}), so $(\bigoplus_{m\ge1}S_m)^2$ embeds in $\Aut(X_{F_{90}})$ and in $\Aut(X_{F_{165}})$. In particular both groups contain every finite group, are nonabelian, and are not virtually $\Z^n$ for any $n$.
\end{theorem}

\begin{proof}
\emph{Locality.} $\Phi_\psi$ and $\Phi'_\psi$ are composites and sums of block maps commuting with $\sh$, so they are block maps.

\emph{The coordinates.} By linearity of $d$ and $e$ and Lemma~\ref{lin:coords}(d),
\[
d\circ\Phi_\psi=\psi\circ d,\qquad e\circ\Phi_\psi=e,\qquad d\circ\Phi'_\psi=d,\qquad e\circ\Phi'_\psi=\psi\circ e .
\]
For instance $d(\Phi_\psi(c))=d(c)+(\mathrm{id}+\sh^2)J_\psi(d(c))=\psi(d(c))$.

\emph{Commutation.} Write $F=F_{90}+k$ with $k=\varepsilon(0,1)^{\Z}$, and recall that $J_\psi$ commutes with $\sh$. By Lemma~\ref{lin:coords}(c), $F(\Phi_\psi(c))=F(c)+F_{90}\iota J_\psi(d(c))=F(c)+\iota\sh J_\psi(d(c))$. By Lemma~\ref{lin:coords}(b), $\Phi_\psi(F(c))=F(c)+\iota J_\psi(\sh d(c)+\varepsilon\mathbf 1)$, and $J_\psi(\sh d(c)+\varepsilon\mathbf 1)=J_\psi(\sh d(c))=\sh J_\psi(d(c))$: for $\varepsilon=0$ trivially, and for $\varepsilon=1$ because $\psi\in\mathcal G^{\kappa}$. So $F\Phi_\psi=\Phi_\psi F$. The same computation with $\iota'$, $e$ and $\sh^{-1}$ gives $F\Phi'_\psi=\Phi'_\psi F$.

\emph{Homomorphism and invertibility.} By the coordinate identities, the linearity of $\iota$ and Lemma~\ref{lin:integral},
\[
\Phi_\psi(\Phi_\chi(c))=c+\iota J_\chi(d(c))+\iota J_\psi(\chi(d(c)))=c+\iota J_{\psi\chi}(d(c))=\Phi_{\psi\chi}(c),
\]
and $\Phi_{\mathrm{id}}=\mathrm{id}$ since $J_{\mathrm{id}}=0$. So $\Phi_\psi$ is invertible with inverse $\Phi_{\psi^{-1}}$; here the invertibility of $\psi=\mathrm{id}+(\mathrm{id}+\sh^2)J_\psi$ is what makes $\Phi_\psi$ invertible. With the commutation, $\Phi_\psi\in\Aut(X_F)$ by Lemma~\ref{lem:spacetime}, and $\psi\mapsto\Phi_\psi$ is a homomorphism $H\to\Aut(X_F)$. The same holds for $\psi\mapsto\Phi'_\psi$.

\emph{The two movers commute.} Since $d\circ\Phi'_\chi=d$ and $e\circ\Phi_\psi=e$,
\[
\Phi_\psi\Phi'_\chi(c)=c+\iota'J_\chi(e(c))+\iota J_\psi(d(c))=\Phi'_\chi\Phi_\psi(c).
\]
So $(\psi,\chi)\mapsto\Phi_\psi\Phi'_\chi$ is a homomorphism $H\times H\to\Aut(X_F)$.

\emph{Injectivity.} If $\Phi_\psi\Phi'_\chi=\mathrm{id}$, then applying $d$ gives $\psi(d(c))=d(c)$ for every $c$, so $\psi=\mathrm{id}$ because $d$ is surjective; applying $e$ gives $\chi=\mathrm{id}$.

The remaining statements follow from Lemma~\ref{lin:markers}, with $\mathcal G^{\kappa}\subseteq\mathcal G$, and from the argument in the proof of Theorem~\ref{thm:54}.
\end{proof}

It remains to produce automorphisms of the binary full shift with integrals. They come from marker permutations whose changes have zero sum on the even and on the odd positions of each segment.

\begin{lemma}\label{lin:markers}
For $L\ge1$ let $U_L$ be the set of words $u\in\{0,1\}^L$ with $u_0=u_{L-1}=0$, with no factor $11$ and no factor $000$, and with $\sum_{k\ \mathrm{even}}u_k\equiv\sum_{k\ \mathrm{odd}}u_k\equiv0\pmod2$. A \emph{frame} of $z\in\F_2^{\Z}$ of length $L$ at $p$ is a position $p$ with $z_{[p,p+L+6)}=111\,u\,111$ (a positive frame) or $z_{[p,p+L+6)}=000\,\bar u\,000$ (a negative frame) for some $u\in U_L$, its \emph{content}; its \emph{segment} is $[p+3,p+L+2]$. For a family $\rho=(\rho_L)_{L\ge1}$ of permutations $\rho_L\in\Sym(U_L)$, all but finitely many equal to the identity, let $\psi_\rho(z)$ be obtained from $z$ by replacing the content $u$ of every frame of length $L$ by $\rho_L(u)$, that is, writing $\rho_L(u)$ on the segment of a positive frame and $\overline{\rho_L(u)}$ on that of a negative one. Then $\rho\mapsto\psi_\rho$ is an injective homomorphism $\bigoplus_L\Sym(U_L)\to\mathcal G^{\kappa}$. Moreover $|U_{4N+11}|\ge N+1$ for $N\ge0$, so $\bigoplus_{m\ge1}S_m$ embeds in $\mathcal G^{\kappa}$.
\end{lemma}

\begin{proof}
(i) \emph{No run of three equal symbols meets a segment.} In a positive frame every $1$ of the segment has two neighbours $0$, since $u$ has no $11$ and starts and ends with $0$; and the cells next to the segment are $1$, so a run $000$ meeting the segment would lie inside it, which $u$ excludes. Negative frames are the complements.

(ii) \emph{Segments of distinct frames are disjoint, at least three cells apart.} Two frames at the same position have the same sign, given by $z_p$, and if their lengths were $L<L'$, then the cell $p+L+3$, which lies in a run of three, would lie in the segment $[p+3,p+L'+2]$ of the longer one, contradicting (i). Let $p<q$ be frames, of lengths $L_p$ and $L_q$. The run $z_{[q,q+3)}$ is disjoint from the segment of $p$ by (i), and $q+2>p+2$, so $q\ge p+L_p+3$. Hence the segment of $q$ starts at $q+3\ge p+L_p+6$, at least three cells after the end $p+L_p+2$ of the segment of $p$.

(iii) \emph{$\psi_\rho$ preserves frames.} Let $z'=\psi_\rho(z)$. By (i) and (ii), $\psi_\rho$ changes only cells of segments and leaves the runs of three equal symbols in place, so every frame of $z$ is a frame of $z'$, of the same length and sign, with content $\rho_L(u)$; note that $\rho_L(u)\in U_L$ again satisfies (i). Conversely let $p$ be a frame of $z'$. If its segment meets the segment of a frame $q$ of $z$, then $q$ is also a frame of $z'$, and $p=q$ by (ii) applied to $z'$. Otherwise $z=z'$ on the segment of $p$, and also on the two runs of $p$: a cell of one of them that lay in a segment of some frame of $z$ would lie in a segment of the same frame of $z'$, contradicting (i) for $z'$. So $p$ is a frame of $z$. Hence $z$ and $z'$ have the same frames, with the same lengths and signs, and the contents of $z'$ are the images under $\rho$ of those of $z$.

(iv) \emph{Group law.} By (iii), $\psi_\rho\psi_{\rho'}=\psi_{\rho\rho'}$ and $\psi_{\mathrm{id}}=\mathrm{id}$; and $\psi_\rho$ is a block map commuting with $\sh$, since a cell reads the frame containing it in its segment, of length at most $\max\{L:\rho_L\ne\mathrm{id}\}+6$. So each $\psi_\rho$ is an automorphism. It fixes $0^{\Z}$, which has no frames, and if $\rho_L(u)\ne u$ then $\psi_\rho$ moves the configuration $0^{\Z}$ with $111\,u\,111$ written at $0$; so $\rho\mapsto\psi_\rho$ is injective. The frames of $\bar z$ are those of $z$ with the signs exchanged and the same contents, so $\psi_\rho(\bar z)=\overline{\psi_\rho(z)}$.

(v) \emph{The integral.} Let $\delta=z+\psi_\rho(z)$. It vanishes outside segments, and on the segment $[a,a+L-1]$ of a frame with content $u$ it is $\delta_{a+k}=u_k+\rho_L(u)_k$, for either sign. Define $J(z)_{a+k}=\sum\{\delta_{a+k'}:\ k\le k'\le L-1,\ k'\equiv k\pmod 2\}$ for $0\le k\le L-1$ on each segment, and $J(z)_n=0$ outside segments. This is a block map with $J(0^{\Z})=0^{\Z}$, and $J(\bar z)=J(z)$, since $\bar z$ has the same segments and the same $\delta$. We check $J(z)_n+J(z)_{n+2}=\delta_n$ for every $n$. If $n$ and $n+2$ lie in one segment, this is the definition. If $n=a+k$ lies in a segment and $n+2$ does not, then $k\ge L-2$, $J(z)_n=\delta_n$, and $n+2$ lies in no segment by (ii), so $J(z)_{n+2}=0$. If $n$ lies in no segment, then $\delta_n=J(z)_n=0$, and either $n+2$ lies in no segment, or by (ii) $n+2\in\{a,a+1\}$ for the start $a$ of a segment; then $J(z)_{n+2}$ is the sum of $\delta$ over the even, respectively odd, positions of that segment, which is $0$ because $u$ and $\rho_L(u)$ both lie in $U_L$. Thus $\psi_\rho=\mathrm{id}+(\mathrm{id}+\sh^2)J$, and $\psi_\rho\in\mathcal G^{\kappa}$.

(vi) \emph{Counting.} For $0\le i\le N$ the word $0\,(10)^{2i+1}\,100\,10\,100\,(10)^{2(N-i)}$ has length $4N+11$, starts and ends with $0$, and has no factor $11$ or $000$. Its $1$'s lie at the odd positions $1,3,\dots,4i+1$, then at $4i+3$, $4i+6$ and $4i+8$, and then at the $2(N-i)$ odd positions $4i+11,\dots$; so $4N+2$ of them are at odd and $2$ at even positions, and the word lies in $U_{4N+11}$. These $N+1$ words are distinct. Embedding $S_m$ in $\Sym(U_{4m+7})$ for each $m$ gives $\bigoplus_mS_m\hookrightarrow\bigoplus_L\Sym(U_L)\hookrightarrow\mathcal G^{\kappa}$.
\end{proof}

\begin{sloppypar}
\begin{example}\label{lin:ex:S3}
A smaller instance uses frames $111\,u\,111$ and $000\,\bar u\,000$ with $u$ in $U=\{0010110,\,0101010,\,0110100\}$, which are not in $U_7$ but all have even-position sum $0$ and odd-position sum $1$. For $\rho\in S_3$, $\psi_\rho$ replaces $u_i$ by $u_{\rho(i)}$, and $J_\rho$ is defined on each segment by the formula in Lemma~\ref{lin:markers}(v). Exhaustive checks on the binary full shift (computer-assisted; Appendix~\ref{app:stated}) show that distinct frames share no segment cell (two frames can overlap only in a shared run, at offset $10$), that $\psi_\rho$ preserves the frames and permutes the contents, on all $2^{31}$ windows of length $31$, that $d+J_\rho(d)+\sh^2J_\rho(d)=\psi_\rho(d)$ on all $2^{21}$ windows, and that $J_\rho(\bar d)=J_\rho(d)$. So $\psi_\rho\in\mathcal G^{\kappa}$, and $\rho\mapsto\Phi_{\psi_\rho}$ embeds $S_3$ in $\Aut(X_{F_{90}})$ and in $\Aut(X_{F_{165}})$ by maps of cell radius $11$. The maps $\Phi_{\psi_{(01)}}$ and $\Phi_{\psi_{(12)}}$ do not commute.
\end{example}
\end{sloppypar}

The maps $\Phi_\psi$ have the form $\mathrm{id}+\iota\circ h\circ u'$ with $u'=d$ a factor map and $h=J_\psi$. In the classification of \cite{PaperD}, rules $90$ and $165$, which are affine, are listed among the classes whose automorphism group contains every finite group.

\begin{corollary}\label{lin:noblock}
For $f\in\{90,165\}$, $F_f$ has no blocking word, and in particular no word fixed in every context, and $\Aut(X_{F_f})$ is not virtually $\Z^2$.
\end{corollary}

\begin{proof}
Let $u$ be a word, $q$ a position in it, and $x$ a configuration containing $u$ on $[0,|u|)$. For $N\ge|u|+q$ let $e_N\in\F_2^{\Z}$ be the unit vector at $N$, and $x'=x+\iota(e_N)$, which differs from $x$ only at the cells $N$ and $N+1$ and so also contains $u$ on $[0,|u|)$. Since $F^n$ differs from $F_{90}^n$ by a constant, Lemma~\ref{lin:coords}(c) gives $F^n(x')-F^n(x)=F_{90}^n(\iota(e_N))=\iota(\sh^ne_N)=\iota(e_{N-n})$, whose $y$-bit at $q$ is $1$ for $n=N-q$. So $F^{N-q}(x')_q\ne F^{N-q}(x)_q$, and $u$ is not blocking with column $q$. A word fixed in every context is blocking. The last statement is Theorem~\ref{lin:movers}.
\end{proof}


\section{Affine rules}\label{sec:affine}

The $16$ affine elementary rules $f(a,b,c)=\varepsilon\oplus\alpha a\oplus\beta b\oplus\gamma c$ form nine classes:
\[
\{0\},\ \{255\},\ \{51,204\},\ \{90\},\ \{165\},\qquad \{15,85,170,240\},\ \{60,102\},\ \{153,195\},\ \{105,150\}.
\]
In each of the first five, $\bigoplus_mS_m$ embeds in $\Aut(X_{F_f})$: for $0$ and $51$ through fixed walls (Table~\ref{tab:walls}), for $90$ and $165$ by Theorem~\ref{lin:movers}, and for $255$ by Proposition~\ref{aff:prop:255} below. For the last four we prove Theorem~\ref{thm:affine}: every automorphism is affine, and the automorphism group is determined. The proof uses the isomorphism rigidity theorem of Kitchens and Schmidt \cite{KS2000}, which we state in Section~\ref{aff:sec:KS}; everything else is proved here.

\subsection{Affine commuting maps}\label{aff:sec:setting}

Let $f$ be affine, with $\Lam=\F_2[s,s^{-1}]$, $p=\alpha s^{-1}+\beta+\gamma s$ and $M=\begin{psmallmatrix}0&1\\1&p\end{psmallmatrix}$ as in Section~\ref{sec:linear}. Then $F_f(c)=Mc+k_0$ with $k_0=(0,\varepsilon)^{\Z}$. Put $q=t^2+pt+1\in\Lam[t]$ and
\[
\mathcal O=\Lam[t]/(q).
\]
Since $t(t+p)=1$ in $\mathcal O$, $t$ is a unit, and $\mathcal O$ is also the quotient of $\F_2[s^{\pm1},t^{\pm1}]$ by $q$. The ring $\Lam[M]=\{aI+bM:a,b\in\Lam\}$ is free over $\Lam$ with basis $I,M$, since the entries of $aI+bM$ are $a$, $b$, $b$, $a+bp$, and $M^2=pM+I$; so $a+bt\mapsto aI+bM$ is a ring isomorphism $\mathcal O\cong\Lam[M]$, under which $\sh$ corresponds to $s$ and the linear part of $F_f$ to $t$. We use it to identify the two rings. The proof of Lemma~\ref{lem:commutant} uses only that $M$ is not scalar, so for every affine rule the matrices in $M_2(\Lam)$ commuting with $M$ are the elements of $\Lam[M]$; in particular they commute with each other. On a constant configuration $k$ an element $a\in\Lam$ acts as the sum $a(1)$ of its coefficients, and $N\in\Lam[M]$ acts as the $2\times2$ matrix $N(1)$ over $\F_2$ obtained by putting $s=1$.

\begin{lemma}\label{aff:lem:comm}
Let $f$ be affine, $N\in M_2(\Lam)$ and $k$ a constant configuration. The block map $\varphi(c)=Nc+k$ commutes with $F_f$ if and only if $N\in\Lam[M]$ and $(N+I)k_0=(M+I)k$. The automorphisms of $X_{F_f}$ that are affine, that is, of this form, form a group $\Aut_{\mathrm{aff}}(X_{F_f})$, and $\varphi\mapsto N$ is a homomorphism from it to the abelian group $\Lam[M]^{\times}$ whose kernel consists of the translations $c\mapsto c+k$ with $(M+I)k=0$, at most four maps. In particular $\Aut_{\mathrm{aff}}(X_{F_f})$ is metabelian.
\end{lemma}

\begin{proof}
$\varphi(F_f(c))=NMc+Nk_0+k$ and $F_f(\varphi(c))=MNc+Mk+k_0$; these agree for all $c$ if and only if $NM=MN$ (compare the parts depending on $c$) and $Nk_0+k=Mk+k_0$. If $\varphi$ is invertible, then so is the continuous additive map $c\mapsto Nc=\varphi(c)+k$, which commutes with $\sh$; its inverse is again continuous, additive and commutes with $\sh$, so by Curtis--Hedlund--Lyndon it is a block map with an $\F_2$-linear local rule, that is, a matrix $N'\in M_2(\Lam)$ with $NN'=N'N=I$. So $N$ is invertible in $M_2(\Lam)$, the inverse $c\mapsto N'(c+k)$ of $\varphi$ is again affine and commutes with $F_f$, and $N'=N^{-1}\in\Lam[M]$. The composite of $c\mapsto Nc+k$ and $c\mapsto N'c+k'$ has linear part $NN'$, and the map with linear part $I$ and constant $k$ commutes with $F_f$ if and only if $(M+I)k=0$.
\end{proof}

So a class with $\bigoplus_mS_m\le\Aut(X_{F_f})$ always has automorphisms that are not affine, since $\bigoplus_mS_m$ is not solvable. This applies to the first five classes above.

\subsection{The rigidity theorem}\label{aff:sec:KS}

Let $\beta$ be an action of $\Z^d$ by continuous automorphisms of a compact abelian group $G$, with normalised Haar measure $\lambda_G$. The action is \emph{irreducible} if every closed $\beta$-invariant subgroup $H\subsetneq G$ is finite, and \emph{mixing} if it is mixing with respect to $\lambda_G$. A \emph{measurable conjugacy} between two such actions $(G,\beta)$ and $(G',\beta')$ is a measure-preserving map $\varphi\colon(G,\lambda_G)\to(G',\lambda_{G'})$, invertible modulo null sets, with $\varphi\circ\beta^{v}=\beta'^{v}\circ\varphi$ almost everywhere for every $v\in\Z^d$.

\begin{theorem}[Kitchens--Schmidt \cite{KS2000}; see also {\cite[Theorem~8.4]{Schmidt2002}}]\label{aff:thm:KS}
Let $d\ge2$, and let $(G,\beta)$ and $(G',\beta')$ be irreducible and mixing actions of $\Z^d$ by continuous automorphisms of compact zero-dimensional abelian groups. Then every measurable conjugacy $\varphi$ between them is affine almost everywhere: there are a continuous surjective group homomorphism $A\colon G\to G'$ and $g'\in G'$ with $\varphi(g)=A(g)+g'$ for $\lambda_G$-almost every $g$.
\end{theorem}

We apply it with $d=2$ and $G=G'=\A^{\Z}=(\F_2^2)^{\Z}$. The dual group $\widehat G$ consists of the characters $\chi_{\xi,\eta}(x,y)=(-1)^{\sum_i(\xi_ix_i+\eta_iy_i)}$ with $\xi,\eta\in\F_2^{\Z}$ finitely supported, and we identify $\xi$ with $\sum_i\xi_is^i\in\Lam$.

\begin{lemma}\label{aff:lem:dual}
Let $f$ be affine with linear part $M$. Under $\chi\mapsto\chi\circ\sh$ and $\chi\mapsto\chi\circ M$, the group $\widehat G$ is an $\F_2[s^{\pm1},t^{\pm1}]$-module, with $s$ acting as $\chi\mapsto\chi\circ\sh$ and $t$ as $\chi\mapsto\chi\circ M$, and it is isomorphic to $\mathcal O$ with $s$ and $t$ acting by multiplication.
\end{lemma}

\begin{proof}
For finitely supported $\xi$ and any $x$, $\sum_i\xi_ix_{i+1}=\sum_i\xi_{i-1}x_i$, so $\chi_{\xi,\eta}\circ\sh=\chi_{s\xi,s\eta}$; and since $(py)_i=\alpha y_{i-1}+\beta y_i+\gamma y_{i+1}$, $\sum_i\eta_i(py)_i=\sum_i(p\eta)_iy_i$. Hence $\chi_{\xi,\eta}(M(x,y))=\chi_{\xi,\eta}(y,x+py)=\chi_{\eta,\xi+p\eta}(x,y)$, so $\chi\circ M$ corresponds to the matrix $M$ acting on the column $(\xi,\eta)^{\mathsf T}\in\Lam^2$. The vector $e_1=(1,0)^{\mathsf T}$ has $Me_1=(0,1)^{\mathsf T}$, so $a+bt\mapsto(aI+bM)e_1=(a,b)^{\mathsf T}$ is a bijection $\mathcal O\to\Lam^2$ compatible with $s$ and $t$.
\end{proof}

\subsection{Every automorphism is affine}\label{aff:sec:R}

\begin{lemma}\label{aff:lem:R}
Let $f$ be affine with linear part $M$, and suppose:
\begin{enumerate}
\item[(R1)] $q=t^2+pt+1$ is irreducible over $K=\F_2(s)$;
\item[(R2)] $s^at^b\ne1$ in $\mathcal O$ for $(a,b)\ne(0,0)$;
\item[(R3)] there is $m\ge1$ with $F_f^m=M^m$ such that the subring $S=\Lam[t^m,t^{-m}]$ of $\mathcal O$ has finite index in $\mathcal O$ as an additive group.
\end{enumerate}
Then every automorphism of $X_{F_f}$ is affine: it has the form $c\mapsto Nc+k$ with $N\in\Lam[M]^{\times}=\mathcal O^{\times}$ and $k$ a constant configuration with $(N+I)k_0=(M+I)k$. Consequently $\Aut(X_{F_f})=\Aut_{\mathrm{aff}}(X_{F_f})$.
\end{lemma}

For a linear rule (R3) holds with $m=1$ and $S=\mathcal O$.

\begin{proof}
Let $\beta$ be the $\Z^2$-action on $G=\A^{\Z}$ generated by $\sh$ and $F_f^m=M^m$, which act by continuous automorphisms. By Lemma~\ref{aff:lem:dual}, $\widehat G$ is $\mathcal O$, with the generators of $\beta$ acting by multiplication by $s$ and $t^m$. By (R1), $\mathcal O$, which is free over $\Lam$ with basis $1,t$, embeds in the field $K[t]/(q)$, so it is a domain.

\emph{Irreducibility.} The map $t\mapsto t+p$ is a ring automorphism of $\mathcal O$, since $t+p$ is again a root of $q$. For $n=a+bt\ne0$ put $\bar n=a+b(t+p)$; then $n\bar n=a^2+abp+b^2=:\nu$ lies in $\Lam$ and is nonzero, as $\mathcal O$ is a domain. So $\nu\mathcal O\subseteq n\mathcal O$, and $\mathcal O/\nu\mathcal O\cong(\Lam/\nu\Lam)^2$ is finite, since $\Lam/\nu\Lam\cong\F_2[s]/(\nu_0)$ for the polynomial $\nu_0$ obtained from $\nu$ by removing powers of $s$. Hence $\mathcal O/n\mathcal O$ is finite, and so is $n\mathcal O/nS\cong\mathcal O/S$ by (R3); so $\mathcal O/nS$ is finite. Now let $H\subsetneq G$ be closed and $\beta$-invariant. Its annihilator $H^\perp\subseteq\widehat G$ is nonzero, and it is invariant under $s^{\pm1}$ and $t^{\pm m}$, hence an $S$-submodule of $\mathcal O$. Choosing $0\ne n\in H^\perp$ we get $nS\subseteq H^\perp$, so $\widehat H\cong\mathcal O/H^\perp$ is finite, and $H$ is finite.

\emph{Mixing.} For characters $\chi,\psi$, $\int(\chi\circ\beta^v)\,\overline{\psi}\,d\lambda_G$ is $1$ if $\chi\circ\beta^v=\psi$ and $0$ otherwise. If $\chi\ne1$ corresponds to $0\ne n\in\mathcal O$ and $\chi\circ\beta^{(a,b)}=\chi\circ\beta^{(a',b')}$, then $(s^{a-a'}t^{m(b-b')}-1)s^{a'}t^{mb'}n=0$, so $s^{a-a'}t^{m(b-b')}=1$ and $(a,b)=(a',b')$ by (R2). So for $\chi\ne1$ the equation $\chi\circ\beta^v=\psi$ has at most one solution $v$, while for $\chi=1$ it holds for all $v$ or none according as $\psi=1$. Hence $\int(\chi\circ\beta^v)\overline\psi\,d\lambda_G\to\int\chi\,d\lambda_G\int\overline\psi\,d\lambda_G$ as $|v|\to\infty$, for all characters. By linearity this holds for trigonometric polynomials, which are dense in $L^2(\lambda_G)$, and since each $\beta^v$ preserves $\lambda_G$ it holds for all $L^2$ functions; this is mixing.

\emph{Haar measure.} Haar measure on $G$ is the uniform Bernoulli measure. Let $\varphi\in\Aut(X_{F_f})$, a block map of some radius $r$ that is onto. By Hedlund \cite[Theorem~5.4]{Hedlund1969}, for an onto block map of the full $4$-shift every word of length $n$ has exactly $4^{2r}$ preimages of length $n+2r$ under the local rule; so the preimage of a cylinder of length $n$ is a union of $4^{2r}$ cylinders of length $n+2r$ and has the same measure $4^{-n}$. Hence $\varphi$ preserves $\lambda_G$. It commutes with $\sh$ and $F_f$, hence with $\beta$, so it is a measurable conjugacy of $(G,\beta)$ with itself.

\emph{Conclusion.} By Theorem~\ref{aff:thm:KS} there are a continuous surjective homomorphism $A\colon G\to G$ and $k\in G$ with $\varphi=A+k$ almost everywhere. Both sides are continuous, and a nonempty open set has positive measure, so $\varphi=A+k$ everywhere. Then $k=\varphi(0^{\Z})$ is fixed by $\sh$, hence constant, and $A=\varphi-k$ commutes with $\sh$ and is bijective, since $\varphi$ is. A continuous additive map of $G$ commuting with $\sh$ is a block map with an $\F_2$-linear local rule (Curtis--Hedlund--Lyndon), that is, a matrix $N\in M_2(\Lam)$. Lemma~\ref{aff:lem:comm} gives $N\in\Lam[M]$, $N$ invertible there, and the condition on $k$.
\end{proof}

\subsection{The unit groups}\label{aff:sec:units}

\begin{lemma}\label{aff:lem:units}
\begin{enumerate}
\item For $p=s^{\pm1}$ \textup{(}rules $170$ and $240$\textup{)}, $\mathcal O\cong\F_2[t,t^{-1},(1+t)^{-1}]$ with $s^{\pm1}=(1+t)^2t^{-1}$, and $\mathcal O^{\times}=t^{\Z}(1+t)^{\Z}\cong\Z^2$.
\item For $p=1+s^{\pm1}$ \textup{(}rules $102$ and $60$\textup{)}, $\mathcal O\cong\F_2[t,t^{-1},(1+t+t^2)^{-1}]$ with $s^{\pm1}=(1+t+t^2)t^{-1}$, and $\mathcal O^{\times}=s^{\Z}t^{\Z}\cong\Z^2$.
\item For $p=s+1+s^{-1}$ \textup{(}rule $150$\textup{)}, $\mathcal O^{\times}=s^{\Z}t^{\Z}(1+st)^{\Z}\cong\Z^3$.
\end{enumerate}
In all three cases \textup{(R1)} and \textup{(R2)} hold.
\end{lemma}

\begin{proof}
(1) Take $p=s$; the case $p=s^{-1}$ follows by the automorphism $s\mapsto s^{-1}$ of $\Lam$. Let $B=\F_2[t,t^{-1},(1+t)^{-1}]$ and $u=(1+t)^2t^{-1}\in B$. The homomorphism $\Lam[t]\to B$ with $s\mapsto u$, $t\mapsto t$ is onto, since $(1+t)^{-1}=(1+t)t^{-1}u^{-1}$, and it kills $q$, since $t^2+ut+1=t^2+(1+t)^2+1=0$. The induced map $\mathcal O\to B$ sends $a+bt$ to $a(u)+b(u)t$, and it is injective: $u$ is transcendental over $\F_2$ and $[\F_2(t):\F_2(u)]=2$, so $1$ and $t$ are linearly independent over $\F_2(u)$. So $\mathcal O\cong B$. The ring $B$ is a localisation of $\F_2[t]$, hence a unique factorisation domain whose units are the products of powers of the inverted primes $t$ and $1+t$. Since $\mathcal O$ is a domain, $q$ is prime in $\Lam[t]$, and being monic it is irreducible over $K$ by Gauss's lemma; this is (R1). In $B$, $s^at^b=(1+t)^{2a}t^{b-a}$, which is $1$ only for $a=b=0$; this is (R2).

(2) The same argument with $u=(1+t+t^2)t^{-1}$, which satisfies $t^2+(1+u)t+1=0$, and the prime $1+t+t^2$ of $\F_2[t]$; here $(1+t+t^2)^{-1}=t^{-1}u^{-1}$, the units are $t^{\Z}(1+t+t^2)^{\Z}=t^{\Z}(st)^{\Z}=s^{\Z}t^{\Z}$, and $s^at^b=(1+t+t^2)^at^{b-a}$.

(3) Irreducibility (R1) was shown in Section~\ref{sec:linear}. Multiplying $q$ by $s$, $\mathcal O$ is the coordinate ring of the affine curve $s^2t+st^2+st+s+t=0$ in the torus $st\ne0$. Let $E\subset\mathbb P^2$ be the cubic
\[
s^2t+st^2+stz+sz^2+tz^2=0 .
\]
Its partial derivatives in characteristic $2$ are $t^2+tz+z^2$, $s^2+sz+z^2$ and $st$; if all vanish, then $s=0$ or $t=0$, and either case forces $s=t=z=0$. So $E$ is smooth, hence an elliptic curve over $\F_2$, geometrically irreducible, with function field the fraction field of $\mathcal O$; in particular $\F_2$ is algebraically closed in it, and a function on $E$ without poles is a constant in $\F_2$. The points of $E$ over $\overline{\F}_2$ outside the affine torus $\{z\ne0,\ s\ne0,\ t\ne0\}$ are, since $s=0$ forces $tz^2=0$, $t=0$ forces $sz^2=0$, and $z=0$ forces $st(s+t)=0$,
\[
P_0=(0:0:1),\quad P_1=(1:0:0),\quad P_2=(0:1:0),\quad P_3=(1:1:0);
\]
these are the four points of $E(\F_2)$. Put $g=1+st$. Its norm is $(1+st)(1+st^{-1})=1+sp+s^2=s$, so $g\in\mathcal O^{\times}$. The divisors of $s=s/z$, $t=t/z$ and $g=(z^2+st)/z^2$ are
\[
\operatorname{div}s=P_0+P_2-P_1-P_3,\qquad \operatorname{div}t=P_0+P_1-P_2-P_3,\qquad \operatorname{div}g=P_1+P_2-2P_3 .
\]
Indeed, the line $s=0$ meets $E$ where $tz^2=0$, that is at $P_0$ simply and at $P_2$ doubly; the line $t=0$ meets $E$ at $P_0$ simply and at $P_1$ doubly; the line $z=0$ meets $E$ simply at $P_1,P_2,P_3$. For $g$: $g$ has no zero or pole on the affine part, it equals $1$ at $P_0$, and its numerator $z^2+st$ is $1$ at $P_3$, so $g$ has a double pole at $P_3$. At $P_1$, in the chart $s=1$ the curve is $t(1+t+z+z^2)=z^2$, so $z$ is a local parameter, $t=z^2w^{-1}$ with $w=1+t+z+z^2$, and $z^2+t=z^2(1+w)w^{-1}$ with $1+w=t+z+z^2$ of order $1$; so $g=(z^2+t)/z^2$ has a simple zero at $P_1$, and by the symmetry $s\leftrightarrow t$ of $E$ and $g$ also at $P_2$.

Let $\mathcal D$ be the group of divisors of degree $0$ supported on $\{P_0,\dots,P_3\}$, which is free of rank $3$, and define $\theta\colon\mathcal D\to\Z/4$ by $\theta(\sum n_iP_i)=2n_0+n_1-n_2$. It vanishes on the three divisors above (on $\operatorname{div}t$ the value is $4$), and it is onto. The three divisors span a subgroup $L\subseteq\ker\theta$ of index $|\det|=4$ in $\mathcal D$, in the basis $P_0-P_3,P_1-P_3,P_2-P_3$, where they are $(1,-1,1)$, $(1,1,-1)$ and $(0,1,1)$. Since $\ker\theta$ also has index $4$, $L=\ker\theta$. In particular $s,t,g$ generate a free abelian group $W$ of rank $3$, and $\operatorname{div}\colon W\to L$ is an isomorphism.

Now let $u\in\mathcal O^{\times}$. Both $u$ and $u^{-1}$ are regular on the affine curve, so $D=\operatorname{div}u\in\mathcal D$. The values $\theta(P_i-P_3)$ are $2,1,3,0$ for $i=0,1,2,3$, so there is an $i$ with $D-(P_i-P_3)\in\ker\theta=L$, say $D-(P_i-P_3)=\operatorname{div}w$ with $w\in W$. Then $P_i-P_3=\operatorname{div}(u/w)$ is principal, and on an elliptic curve a divisor $P-P'$ is principal only if $P=P'$ \cite[Cor.~III.3.5]{Silverman2009}. So $i=3$, $u/w$ has no zeros or poles, and $u/w\in\F_2^{\times}=\{1\}$. Hence $\mathcal O^{\times}=W$. Finally (R2) holds since $\operatorname{div}(s^at^b)=(a+b,\,b-a,\,a-b)$ in the basis above vanishes only for $a=b=0$.
\end{proof}

\subsection{The automorphism groups}\label{aff:sec:groups}

For $p(1)=1$, that is, when $f$ has an odd number of linear terms, conjugation by the cellwise complement $\kappa(c)=c+(1,1)^{\Z}$ relates the two members of a pair $\{f,255-f\}$: on constants $M(1)=\begin{psmallmatrix}0&1\\1&1\end{psmallmatrix}$, so $(M+I)(1,1)^{\Z}=(0,1)^{\Z}$, and hence
\begin{equation}\label{aff:eq:kappa}
F_{255-f}=\kappa\,F_f\,\kappa\qquad\text{when }f\text{ is linear and }p(1)=1 .
\end{equation}
Since $\kappa$ is a homeomorphism commuting with $\sh$, $\varphi\mapsto \kappa\varphi\kappa$ is an isomorphism $\Aut(X_{F_f})\to\Aut(X_{F_{255-f}})$. This relates $170$ to $85$, $240$ to $15$ and $150$ to $105$.

\begin{proof}[Proof of Theorem~\ref{thm:affine}]
(1) Let $f\in\{170,240\}$. By Lemma~\ref{aff:lem:units}(1), Lemma~\ref{aff:lem:R} applies with $m=1$. Since $p(1)=1$, $M(1)+I=\begin{psmallmatrix}1&1\\1&0\end{psmallmatrix}$ is invertible, so the only constant with $(M+I)k=0$ is $0^{\Z}$, and $\varphi\mapsto N$ is an isomorphism $\Aut(X_F)\to\mathcal O^{\times}=t^{\Z}(1+t)^{\Z}$, under which $F\mapsto t$ and $h=\mathrm{id}+F\mapsto1+t$. So $(a,b)\mapsto F^ah^b$ is an isomorphism $\Z^2\to\Aut(X_F)$. As $s^{\pm1}=(1+t)^2t^{-1}$, $h^2=\sh^{\pm1}F$, so $\langle\sh,F\rangle=\langle F,h^2\rangle$ has index $2$. For $f\in\{15,85\}$ conjugate by $\kappa$, using \eqref{aff:eq:kappa}: $\kappa h\kappa(c)=c+Mc+M(1,1)^{\Z}$, which is $c+F_f(c)+(1,1)^{\Z}$ because $F_f(c)=Mc+(0,1)^{\Z}$ and $M(1,1)^{\Z}=(1,0)^{\Z}$.

(2) Let $f\in\{60,102\}$. By Lemma~\ref{aff:lem:units}(2), Lemma~\ref{aff:lem:R} applies with $m=1$, and $\mathcal O^{\times}=s^{\Z}t^{\Z}$ consists of the linear parts of the maps $\sh^aF^b$. Since $p(1)=0$, $M(1)=\begin{psmallmatrix}0&1\\1&0\end{psmallmatrix}$, and the constants with $(M+I)k=0$ are $0^{\Z}$ and $(1,1)^{\Z}$. So $\varphi\mapsto N$ is onto $\mathcal O^\times$, with kernel $\{\mathrm{id},\kappa\}$, and it is split by $s^at^b\mapsto\sh^aF^b$. The involution $\kappa$ commutes with every automorphism, since $N(1,1)^{\Z}=(1,1)^{\Z}$ for $N=s^at^b$; so $\Aut(X_F)=\langle\sh,F\rangle\times\langle\kappa\rangle$.

(3) Let $f\in\{153,195\}$, with linear parts $102$ and $60$, so $\mathcal O$ and $\mathcal O^{\times}=s^{\Z}t^{\Z}$ are as in (2). On constants $M(1)$ is the transposition of the two bits and $k_0=(0,1)^{\Z}$, so $(I+M+M^2+M^3)k_0=0$ and $F^4=M^4$. In $\mathcal O$, $t^4=p^3t+p^2+1$ and $t^{-4}=(t+p)^4=p^3t+p^4+p^2+1$, so $S=\Lam[t^{4},t^{-4}]$ is contained in $\Lam\oplus\Lam p^3t$; the latter is a subring, as $(p^3t)^2=p^4\cdot p^3t+p^6$, and it is generated by $t^{\pm4}$ over $\Lam$, so $S=\Lam\oplus\Lam p^3t$ and $\mathcal O/S\cong\Lam/p^3\Lam$, a group of order $8$. So Lemma~\ref{aff:lem:R} applies with $m=4$. For $N=s^at^b$, $(N+I)k_0=(M(1)^b+I)(0,1)$ is $0$ for even $b$ and $(1,1)$ for odd $b$, and $(M+I)(x,y)^{\Z}=(x+y,x+y)^{\Z}$; so for each $N$ there are exactly two constants $k$, differing by $(1,1)^{\Z}$, and the automorphisms with linear part $s^at^b$ are $\sh^aF^b$ and $\kappa\sh^aF^b$. The involution $\kappa$ commutes with $F$ (Section~\ref{sec:prelim}), and $\kappa\notin\langle\sh,F\rangle$ since its linear part is $I$. As in (2), $\Aut(X_F)=\langle\sh,F\rangle\times\langle\kappa\rangle$.

(4) Let $f=150$. By Lemma~\ref{aff:lem:units}(3), Lemma~\ref{aff:lem:R} applies with $m=1$; as in (1) there is no nonzero constant with $(M+I)k=0$, so $\varphi\mapsto N$ is an isomorphism onto $\mathcal O^{\times}=s^{\Z}t^{\Z}(1+st)^{\Z}$, under which $\sh\mapsto s$, $F\mapsto t$ and $h=\mathrm{id}+\sh F\mapsto1+st$. For $f=105$ conjugate by $\kappa$: $\kappa h\kappa(c)=c+\sh(Mc)+(1,0)^{\Z}=c+\sh F_{105}(c)+(1,1)^{\Z}$.

In all four classes every automorphism is affine by Lemma~\ref{aff:lem:R}. The statements for the remaining members follow by reflection, which conjugates $F_{170}$ to $F_{240}$, $F_{85}$ to $F_{15}$, $F_{102}$ to $F_{60}$ and $F_{153}$ to $F_{195}$.
\end{proof}

\begin{remark}\label{aff:rem:ledrappier}
For $f=170$ the ring $\mathcal O\cong\F_2[t^{\pm1},(1+t)^{-1}]$ is also $\F_2[u_1^{\pm1},u_2^{\pm1}]/(1+u_1+u_2)$, with $u_1=t$ and $u_2=1+t$, the dual module of Ledrappier's shift $\{z\in\F_2^{\Z^2}:z_v+z_{v+e_1}+z_{v+e_2}=0\}$ \cite{Ledrappier1978}. By duality, $X_{F_{170}}$ with the action of $F$ and $h=\mathrm{id}+F$ is conjugate to Ledrappier's shift with its two coordinate shifts; explicitly, $c\mapsto(y(F^ah^b(c)_0))_{(a,b)\in\Z^2}$ satisfies the three-dot relation because $h^{b+1}=h^b+Fh^b$. Under this conjugacy $\langle\sh,F\rangle$ is the restriction of the Ledrappier action to the subgroup of index $2$ generated by $e_1$ and $2e_2-e_1$, and Theorem~\ref{thm:affine}(1) says that $\Aut(X_{F_{170}})$ is the full Ledrappier action.
\end{remark}

\begin{remark}\label{aff:rem:consist}
The hypotheses of Lemma~\ref{aff:lem:R} fail for the five other affine classes, and each of them has automorphisms that are not affine (Lemma~\ref{aff:lem:comm} and the remark after it). For $90$ and $165$, $q=(t+s)(t+s^{-1})$ is reducible; for $0$ and $255$, $p=0$ and $q=(t+1)^2$; for $51$ and $204$, $p=1$ and $q=t^2+t+1$ is irreducible over $K$, but $t^3=1$ in $\mathcal O$, so (R2) fails and the action is not mixing. So among the affine rules, (R1) and (R2) hold exactly for the classes whose automorphisms are all affine.
\end{remark}

\begin{remark}\label{aff:rem:endo}
Lemma~\ref{aff:lem:R} concerns automorphisms. Whether every block map commuting with $F_f$, invertible or not, is affine for the classes of Theorem~\ref{thm:affine} is open. By Lemma~\ref{aff:lem:comm} the affine commuting block maps of radius at most $R$, for $R\ge1$, form an affine $\F_2$-space of dimension $4R+1$ for $170$, $4R+2$ for $102$, $4R+1$ for $153$ and $4R$ for $150$: the linear parts $aI+bM$ of radius at most $R$ have $a$ of degree in $[-R,R]$ and $b$ of degree in $[-R,R-1]$ for $p=s$ and $p=1+s$, and in $[-R+1,R-1]$ for $150$, as a direct check of the entries $a$, $b$, $b$, $a+bp$ shows; the constants are $0^{\Z}$ only for $170$ and $150$, and $0^{\Z}$, $(1,1)^{\Z}$ for $102$; and for $153$ the condition $(N+I)k_0=(M+I)k$ cuts the linear parts by one affine condition, $a(1)+b(1)=1$, and leaves two constants for each. For $R\le3$ the elimination of Section~\ref{sec:method} gives exactly these dimensions for the space of all commuting block maps (Appendix~\ref{app:census}), so for these rules every commuting block map of radius at most $3$ is affine.
\end{remark}

\subsection{Rule 255}\label{aff:sec:255}

For $f=255$, $F(x,y)=(y,1\oplus x)$ permutes the cells as $(0,0)\mapsto(0,1)\mapsto(1,1)\mapsto(1,0)\mapsto(0,0)$. Relabel the cells by $\Z/4$, writing $0,1,2,3$ for these four values in this order; then a configuration becomes $x\in(\Z/4)^{\Z}$, and $F(x)=x+\mathbf 1$ adds $1$ in every cell. So $\Aut(X_{F_{255}})$ is the group of automorphisms of the full shift on $\Z/4$ that commute with adding the constant $\mathbf 1$, and $\langle\sh,F\rangle=\langle\sh\rangle\times\langle F\rangle\cong\Z\times\Z/4$. Let $\Delta\colon(\Z/4)^{\Z}\to(\Z/4)^{\Z}$, $\Delta(x)_i=x_{i+1}-x_i$, that is $\Delta=\sh-\mathrm{id}$. It is an onto block map whose fibres are the cosets of the four constant configurations.

\begin{definition}\label{aff:def:G4}
Let $\mathcal G_4$ be the set of automorphisms $\psi$ of the full shift $(\Z/4)^{\Z}$ for which there is a block map $J$ of $(\Z/4)^{\Z}$ with $\psi=\mathrm{id}+(\sh-\mathrm{id})J$.
\end{definition}

Every $\psi\in\mathcal G_4$ preserves the sum over a period of every spatially periodic configuration, since $(\sh-\mathrm{id})J(z)$ has sum $0$ over a period of $z$.

\begin{proposition}\label{aff:prop:255}
Let $F=F_{255}$. For $\varphi\in\Aut(X_F)$ there is a unique $\tilde\varphi\in\mathcal G_4$ with $\Delta\circ\varphi=\tilde\varphi\circ\Delta$, and a unique $\epsilon(\varphi)\in\Z/4$ with $\varphi(0^{\Z})=\epsilon(\varphi)^{\Z}$. The map $\varphi\mapsto(\epsilon(\varphi),\tilde\varphi)$ is an isomorphism
\[
\Aut(X_{F_{255}})\;\cong\;\Z/4\times\mathcal G_4 ,
\]
with $\Z/4=\langle F\rangle$. Moreover $\bigoplus_{m\ge1}S_m$ embeds in $\mathcal G_4$, so $\Aut(X_{F_{255}})$ contains every finite group.
\end{proposition}

\begin{proof}
\emph{The map $\varphi\mapsto\tilde\varphi$.} Since $\varphi(x+\mathbf c)=\varphi(x)+\mathbf c$ for every constant $\mathbf c$, the configuration $\varphi(x)-x$ depends only on the coset $x+\{\text{constants}\}$, that is on $\Delta x$: there is a map $J$ with $\varphi(x)=x+J(\Delta x)$. It commutes with $\sh$, and $J(z)_0=\varphi(x)_0-x_0$ depends on $x_{[-r,r]}$ only up to adding a constant, that is on $z_{[-r,r-1]}$; so $J$ is a block map. Then $\Delta\varphi(x)=\Delta x+(\sh-\mathrm{id})J(\Delta x)$, so $\tilde\varphi=\mathrm{id}+(\sh-\mathrm{id})J$ satisfies $\Delta\varphi=\tilde\varphi\Delta$, and it is unique because $\Delta$ is onto. It is onto because $\varphi$ is, and injective: if $\tilde\varphi(\Delta x)=\tilde\varphi(\Delta x')$, then $\varphi(x')=\varphi(x)+\mathbf c=\varphi(x+\mathbf c)$ for a constant $\mathbf c$, so $x'=x+\mathbf c$ and $\Delta x'=\Delta x$. So $\tilde\varphi\in\mathcal G_4$, and $\varphi\mapsto\tilde\varphi$ is a homomorphism, again because $\Delta$ is onto.

\emph{Kernel.} If $\tilde\varphi=\mathrm{id}$, then $(\sh-\mathrm{id})J=0$, so every $J(z)$ is constant. If $J$ has radius $r$, then for any $z,z'$ and large $n$ a configuration $z''$ agreeing with $z$ on $[-r,r]$ and with $z'$ on $[n-r,n+r]$ gives $J(z)_0=J(z'')_0=J(z'')_n=J(z')_n=J(z')_0$, so $J$ is a constant $\mathbf c$ and $\varphi=F^{c}$. So the kernel is $\langle F\rangle\cong\Z/4$.

\emph{Onto $\mathcal G_4$.} Let $\psi=\mathrm{id}+(\sh-\mathrm{id})J\in\mathcal G_4$ and $\varphi(x)=x+J(\Delta x)$. This is a block map commuting with $\sh$ and with adding constants, and $\Delta\varphi=\psi\Delta$. It is injective: $\varphi(x)=\varphi(x')$ gives $\psi(\Delta x)=\psi(\Delta x')$, so $x'=x+\mathbf c$, and then $\varphi(x')=\varphi(x)+\mathbf c$ forces $\mathbf c=0$. It is onto: given $w$, choose $x$ with $\Delta x=\psi^{-1}(\Delta w)$; then $\Delta\varphi(x)=\Delta w$, so $\varphi(x)=w+\mathbf c$ and $\varphi(x-\mathbf c)=w$. So $\varphi\in\Aut(X_F)$ and $\tilde\varphi=\psi$.

\emph{Splitting.} $\varphi(0^{\Z})$ is fixed by $\sh$, hence constant, and $\varphi\chi(0^{\Z})=\varphi(\epsilon(\chi)^{\Z})=\varphi(0^{\Z})+\epsilon(\chi)^{\Z}$, so $\epsilon$ is a homomorphism to $\Z/4$; it maps $F^c$ to $c$, so it is an isomorphism on the kernel $\langle F\rangle$, which is central. Hence $\Aut(X_F)=\langle F\rangle\times\ker\epsilon$, and $\varphi\mapsto\tilde\varphi$ maps $\ker\epsilon$ isomorphically onto $\mathcal G_4$.

\emph{Markers.} For $L\ge1$ let $V_L$ be the set of words $u\in\{1,2,3\}^L$ with $\sum_ku_k\equiv0\pmod4$. Call a \emph{segment} of $z\in(\Z/4)^{\Z}$ of length $L$ an interval $[a,a+L-1]$ on which $z$ has no $0$, with $z_{a-1}=z_{a+L}=0$; distinct segments are disjoint and separated by a $0$. For $\rho=(\rho_L)_L$ with $\rho_L\in\Sym(V_L)$, all but finitely many the identity, let $\psi_\rho(z)$ replace the content $u$ of every segment of length $L$ with $u\in V_L$ by $\rho_L(u)$. It changes only nonzero cells to nonzero values, so it preserves the set of zeros, hence the segments, their lengths and the property $u\in V_L$. So $\psi_\rho\psi_{\rho'}=\psi_{\rho\rho'}$, each $\psi_\rho$ is an automorphism of the full shift, a block map of radius at most $\max\{L:\rho_L\ne\mathrm{id}\}$, and $\rho\mapsto\psi_\rho$ is injective, since $\psi_\rho$ moves the configuration $0^{\infty}u\,0^{\infty}$ when $\rho_L(u)\ne u$. Let $\delta=\psi_\rho(z)-z$; on a segment $[a,a+L-1]$ it has sum $0$ modulo $4$, because $u$ and $\rho_L(u)$ lie in $V_L$, and it vanishes off segments. Put $J(z)_{a+j}=\sum_{0\le k<j}\delta_{a+k}$ for $0\le j\le L-1$ on every segment, and $J(z)_n=0$ off segments. This is a block map, and $J(z)_{n+1}-J(z)_n=\delta_n$ for every $n$: inside a segment by definition, at its last cell because the sum of $\delta$ over the segment is $0$, and at the $0$ before a segment because $J$ vanishes at the first cell of a segment. So $\psi_\rho=\mathrm{id}+(\sh-\mathrm{id})J\in\mathcal G_4$, and $\bigoplus_L\Sym(V_L)$ embeds in $\mathcal G_4$. The $2^m$ words of length $2m$ made of the blocks $13$ and $31$ lie in $V_{2m}$, so $S_m$ embeds in $\Sym(V_{2m})$, and $\bigoplus_mS_m$ embeds in $\mathcal G_4$.
\end{proof}

So $\Aut(X_{F_{255}})$ contains every finite group and is not virtually $\Z^n$ for any $n$. Through $\varphi\mapsto\tilde\varphi$, its automorphism group is, up to the central factor $\langle F\rangle$, the group $\mathcal G_4$ of automorphisms of the full shift on $\Z/4$ whose displacement $\psi-\mathrm{id}$ is a difference $(\sh-\mathrm{id})J$ of a block map.

\section{Open questions}\label{sec:questions}

\begin{enumerate}
\item Is $\Aut(X_{F_{54}})$ generated by $\sh$, $F$, the wall relabellings of Theorem~\ref{thm:54} and the wall-trading involutions of the kind in Theorem~\ref{thm:tau}, of all radii? At radius $2$ the relabellings suffice (Theorem~\ref{thm:54census}); at radius $3$ they do not, by Theorem~\ref{thm:tau}. A complete count of $\Cent^{\times}_3(F_{54})$ would be a first test.
\item Determine $\Aut(X_{F_{90}})$. Its elements of radius at most $2$ are affine, but it contains $(\bigoplus_mS_m)^2$ (Theorem~\ref{lin:movers}). Is it generated by the affine automorphisms and the maps $\Phi_\psi$, $\Phi'_\chi$ of that theorem?
\item For the affine classes of Theorem~\ref{thm:affine}, is every block map commuting with $F_f$ affine? Every automorphism is (Lemma~\ref{aff:lem:R}), and every commuting block map of radius at most $3$ is (Remark~\ref{aff:rem:endo}).
\item In every example of this paper with a nonlinear rule, the automorphisms outside $\langle\sh,F\rangle$ come from a blocking word with a degree of freedom next to the blocking column. For a rule that is not affine, the absence of blocking words does not force the automorphism group to be virtually $\Z^2$: the class of $95$ is a counterexample \cite{PaperE}. Question~\ref{q:dichotomy} asks whether it is the only one among the elementary rules.
\end{enumerate}


\appendix
\section{Computations}\label{app:comp}

This appendix describes the finite computations used in the paper. Every result marked as computer-assisted in the text is the output of a terminating exact search over a finite set; where random sampling is used, this is said, and such runs serve only as checks. The code and its output accompany the paper. Section~\ref{app:census} describes the censuses of commuting block maps, Section~\ref{app:proofs} the computations on which proofs rely, Section~\ref{app:stated} the computations whose results are stated in the text without a proof by hand, and Section~\ref{app:confirm} computational confirmations of results that are proved by hand. Section~\ref{app:repro} lists running times.

\subsection{The censuses}\label{app:census}

The censuses of Theorem~\ref{thm:54census} and Proposition~\ref{hw:census}, and the dimension counts of Theorem~\ref{thm:90}(c), Proposition~\ref{prop:150} and Remark~\ref{aff:rem:endo}, use the encoding, the enumeration, the re-verification and the invertibility test of Section~\ref{sec:method}.

\emph{Rule $54$ at radius $2$.} The full centralizer $\Cent_2(F_{54})$ exceeds the enumeration cap of $2\times10^4$ solutions, so the balanced commuting maps were enumerated, with the balance condition as native cardinality constraints; there are $23$, and all pass the pair-graph test.

\emph{Rule $54$ at radius $3$.} The enumeration of $\Cent^{\times}_3(F_{54})$ described in Section~\ref{sec:54r3} adds to the CNF the injectivity of $g$ on all spatially periodic configurations of period at most $5$ and learns necessary conditions lazily (diamonds of length $13$, and collisions of periodic configurations of period at most $10$), shared among the worker processes; the balanced commuting maps number more than $2\times10^4$. The enumeration is not complete, and the paper uses from it only the qualitative statements of Section~\ref{sec:54r3}.

\emph{Rule $59$ at radius $2$} (Proposition~\ref{hw:census}). $\Cent^{\times}_2(F_{59})$ was enumerated as in Section~\ref{sec:method}; each of the $133$ solutions passed the pair-graph test, and each was matched against the tables of the maps (i)--(iii) of the proposition, computed by exact composition.

\emph{Rule $90$ at radius $2$.} The invertibility of the $2^{17}$ commuting maps was decided with the balance condition and the pair-graph test, in $5{,}851$ seconds: $98{,}222$ of them are balanced, and $26$ are invertible, all affine.

\emph{Linear rules.} The dimensions of $\Cent_R(F_{90})$ and $\Cent_R(F_{150})$ for $R\le3$ were computed by the elimination of Section~\ref{sec:method}, and for $R\le2$ cross-checked against the SAT enumeration of all solutions. For Theorem~\ref{thm:90}(c) the tables of the three families were built directly, each re-verified to commute, and their span was computed by a rank computation over $\F_2$.

\emph{Affine rules} (Remark~\ref{aff:rem:endo}). For $f\in\{170,102,153,150\}$ and $R\le3$ the commutation condition, an affine system over $\F_2$ in the $2\cdot4^{2R+1}$ table bits (with a constant term for $153$), was solved by elimination in an implementation independent of the SAT encoding. The solution spaces are consistent and have dimensions $5,9,13$ for $170$, $6,10,14$ for $102$, $5,9,13$ for $153$ and $4,8,12$ for $150$, at $R=1,2,3$, which are the dimensions $4R+1$, $4R+2$, $4R+1$ and $4R$ of the affine commuting maps. By reflection the same holds for $240$, $60$ and $195$, and by the conjugation \eqref{aff:eq:kappa} for $85$, $15$ and $105$.

\subsection{Computations used in proofs}\label{app:proofs}

\emph{Fixed words of all lengths} (Proposition~\ref{cl:prop:133}(3) and Corollary~\ref{cor:48}). For each of the $256$ rules and each $k\le40$, the words $w\in\{0,1\}^k$ that are fixed in every context are found as walks in the graph on the four pairs $(w_j,w_{j+1})$, as in the proof. Fixed words exist exactly for the $81$ rules of Table~\ref{tab:walls} at length at most $3$, and, among the other rules, only for $133$ and $161$, at the lengths $4$, $7$, $8$ and every $k$ with $10\le k\le40$. A brute-force search over all of $\A^k$, $k\le4$, agrees with the reduction to $\{0,3\}^k$. For each rule of Table~\ref{tab:walls}, the set $C$ of Theorem~\ref{thm:walls} for the listed word has at least two elements.

\subsection{Computer-assisted statements}\label{app:stated}

The following results are stated in the text as computations. For each we give the method in brief.

\emph{Walls and searches} (Sections~\ref{hw:sec:59} and~\ref{hw:sec:27}). Words fixed or time-periodic in every context, and invariant families of type (W1), are found by computing, for each word, its images under $F^{\pm p}$ in all contexts. The search for invariant families with a hidden involution satisfying (W1) and (H1), with windows of length at most $6$, considers the closure of each word under $F$ and $F^{-1}$ in all contexts when that closure has at most $300$ words. For $59$R and $p\in\{6,12,30,60,120\}$, a search finds, for every word of $X_A\cup X_B$, a context in which $F^p$ changes the protector. The overlaps of walls of rule $59$ at distances $2$ and $3$ were found on all words of length at most $9$.

\emph{The marker example} (Example~\ref{lin:ex:S3}). The statements about frames are checked exhaustively on all binary words of length $31$ and $21$, as stated there.

\emph{Rule $219$} (after Proposition~\ref{hw:219D}). The radius-$6$ variant of $\Psi_2$ is checked to commute with $F$ on all $4^{15}$ windows and to be an involution.

\emph{Linear rules} (Remark~\ref{lin:rem:window}). The graphs $\Lambda_M$ are built directly from the local rule and tested for strong connectivity, for $M\le4$ for the affine left-permutive rules listed there, and for $M\le8$ for rule $90$.

\emph{Rule $54$} (after Theorem~\ref{thm:54} and Table~\ref{tab:spread}). The orders $o_L$ of $F_L$ and the periods of the content $0^L$ were computed by iterating $F_L$ on $S_L$. The spreading data of Table~\ref{tab:spread} come from simulation on random configurations, as stated in its caption.

\subsection{Confirmations of results proved by hand}\label{app:confirm}

The following computations check results that the paper proves by hand. None of them is used in a proof.

\emph{Walls.} The construction of Theorem~\ref{thm:finite} was checked on $3{,}000$ random cyclic configurations of length about $60$ containing trains of $2$ to $4$ cars: $\Psi_m(\sigma)$ commuted with $F$, was inverted by $\Psi_m(\sigma^{-1})$ and composed as a homomorphism, and it acted nontrivially on $476$ of them. For each rule of Table~\ref{tab:walls}, the construction of Theorem~\ref{thm:walls} was checked on $1{,}500$ random cyclic configurations. The table of $\tau$ (Theorem~\ref{thm:tau}) differs from that of the identity on exactly $24$ of the $16{,}384$ entries, commutes with $F$ on all $4^9$ windows, composes with itself to the identity on all $4^{13}$ windows of length $13$, and differs from $D\tau$ on $147{,}456$ of them; the pair-graph test reports $\tau$ injective, and $\tau$ is injective on all spatially periodic configurations of period at most $10$. The automorphism $D$ of Proposition~\ref{cl:prop:133} commutes with $F$ on all $4^{11}$ windows of length $11$.

\emph{Hidden and rotating walls.} For Proposition~\ref{hw:census}, the elements of $T_R$ were composed exactly on windows of radius $4$, and the eight generators were checked to commute with $F$ exhaustively at radius $3$. For $59$R (Theorem~\ref{hw:59}), on rings of length $2000$ seeded with trains of length up to $6$, with all $m=2,\dots,6$ permuted at once, $F\Psi=\Psi F$, $F^{-1}\Psi=\Psi F^{-1}$ and the homomorphism law held for $60$ steps, and the radius bound $2m+2$ held when all cells outside the window were re-randomized. For $127$R, $179$R and $223$R (Proposition~\ref{hw:tp}), the flip of all walls was composed exactly at radius $4$ and passed the pair-graph test, and rings of length $1500$ seeded with trains at gaps $5$, $6$ and $8$, with walls of other classes mixed in, showed no failure of commutation with $F$ and $F^{-1}$ or of the homomorphism law. For $27$R (Theorem~\ref{hw:27}), (W1) and (H1) were checked for the $12$ words in all $16$ contexts, for $F$ and $F^{-1}$; (H2) and (H4) on all pairs of occurrences at distances $3$, $4$, $5$, that is, on words of length at most $11$; over all $16$ contexts, $\beta$ and $\beta+1$ are the only maps $W_{27}\to\Z/2$ with $\beta\circ\iota=\beta+1$ that satisfy (H3) for some $\varepsilon$; and simulations on rings with trains at gaps $3$, $6$ and $7$ showed no failure of commutation or of the homomorphism law. For $219$R and $126$R (Theorem~\ref{hw:219} and Proposition~\ref{hw:126}), the hypotheses were checked on all contexts by two implementations; the bond invariance, the $\pi$-evolution of bonded cells, the maps $\varphi_s$ and $R_s$, and the commutation of $D$ and $E$ with $F$ on all $4^7$ windows were checked, $D$ and $E$ are injective by the pair-graph test, and $\Psi_1$, $\Psi_2$, $\Psi_3$ were tested on planted rings.

\emph{Markers for $90$R and $165$R} (Theorem~\ref{lin:movers}). The maps $\Phi$ and $\Phi'$ were checked on all cyclic configurations of periods $10$ and $11$ and on planted random rings for commutation with $F$, the group law and the commutation of the two movers, with no failures; as a negative control, for $165$R the frames $111\,u\,111$ alone, whose integral is not complement-invariant, give commutation failures. The construction of Lemma~\ref{lin:markers} with $N=2$, three words of length $19$, was checked in the same way.

\emph{Theorem~\ref{cl:thm:A} and Lemma~\ref{cl:lem:stair}.} Parts (1) and (2) were checked by exhaustion for all $256$ rules and $n\le4$, and the construction of Lemma~\ref{cl:lem:stair} with random staircase data for all $240$ rules reading $c$ and $n\le31$.

\subsection{Running times and reproducibility}\label{app:repro}

The censuses were run on one core of a machine with 32\,GB of memory under Linux, using Python~3.12 or later, NumPy and PySAT~1.9 with CaDiCaL~1.5.3 and MiniCard. Table~\ref{tab:times} lists instance sizes and times. The $54$R radius-$2$ invertible census used MiniCard's native cardinality constraints for the balance condition. The code for all computations of this appendix accompanies this paper, with its logs and with the solution tables of the complete enumerations for $54$R at radius $1$, for the $23$ balanced maps of $54$R at radius $2$, and for $59$R; the larger solution sets are regenerated by the programs.

\begin{table}[h]
\caption{Instance sizes and running times (encode / enumerate / classify, in seconds).}
\label{tab:times}
\centering
\begin{tabular}{lrrrl}
\toprule
instance & variables & clauses & solutions & time (s)\\
\midrule
54R, $R=2$, balanced & $2{,}048+4{,}096$ & $2.9\times10^5$+card. & 23 & $3$ / $11$ / --\\
90R, $R=2$, pair-graph test & --- & --- & $2^{17}$ maps tested & $5{,}851$\\
90R, $R=3$, elimination & $4{,}096$ classes & $2.6\times10^5$ eqs & $2^{65}$ & $4$\\
\bottomrule
\end{tabular}
\end{table}


\medskip
\noindent\textbf{Use of AI tools.} Large language models (Anthropic's Claude models) were used extensively in this work: to propose and develop proofs, to design, write and run the computer searches and their independent re-implementations, to check proofs and computations, and to draft the manuscript. The models are not authors. The author has reviewed the content and takes full responsibility for it. All code and computational output accompany the paper (Appendix~\ref{app:comp}).

\begin{thebibliography}{99}

\bibitem{BRY2018} M.~Baake, J.~A.~G.~Roberts and R.~Yassawi, Reversing and extended symmetries of shift spaces, \emph{Discrete Contin. Dyn. Syst.} \textbf{38} (2018), 835--866.

\bibitem{CaDiCaL} A.~Biere, K.~Fazekas, M.~Fleury and M.~Heisinger, CaDiCaL, Kissat, Paracooba, Plingeling and Treengeling entering the SAT Competition 2020, in \emph{Proc. SAT Competition 2020}, Univ. Helsinki, 2020, 51--53.

\bibitem{BT2000} F.~Blanchard and P.~Tisseur, Some properties of cellular automata with equicontinuity points, \emph{Ann. Inst. H. Poincar\'e Probab. Statist.} \textbf{36} (2000), 569--582.

\bibitem{BNR1991} N.~Boccara, J.~Nasser and M.~Roger, Particlelike structures and their interactions in spatiotemporal patterns generated by one-dimensional deterministic cellular-automaton rules, \emph{Phys. Rev. A} \textbf{44} (1991), 866--875.

\bibitem{BLR1988} M.~Boyle, D.~Lind and D.~Rudolph, The automorphism group of a shift of finite type, \emph{Trans. Amer. Math. Soc.} \textbf{306} (1988), 71--114.

\bibitem{CFK2019} V.~Cyr, J.~Franks and B.~Kra, The spacetime of a shift endomorphism, \emph{Trans. Amer. Math. Soc.} \textbf{371} (2019), 461--488.

\bibitem{Hedlund1969} G.~A.~Hedlund, Endomorphisms and automorphisms of the shift dynamical system, \emph{Math. Systems Theory} \textbf{3} (1969), 320--375.

\bibitem{Hochman2010} M.~Hochman, On the automorphism groups of multidimensional shifts of finite type, \emph{Ergodic Theory Dynam. Systems} \textbf{30} (2010), 809--840.

\bibitem{PySAT} A.~Ignatiev, A.~Morgado and J.~Marques-Silva, PySAT: a Python toolkit for prototyping with SAT oracles, in \emph{Theory and Applications of Satisfiability Testing (SAT 2018)}, Lecture Notes in Comput. Sci. 10929, Springer, 2018, 428--437.

\bibitem{KS1989} B.~Kitchens and K.~Schmidt, Automorphisms of compact groups, \emph{Ergodic Theory Dynam. Systems} \textbf{9} (1989), 691--735.

\bibitem{KS1992} B.~Kitchens and K.~Schmidt, Markov subgroups of $(\Z/2\Z)^{\Z^2}$, in \emph{Symbolic Dynamics and its Applications}, Contemp. Math. 135, Amer. Math. Soc., 1992, 265--283.

\bibitem{KS2000} B.~Kitchens and K.~Schmidt, Isomorphism rigidity of irreducible algebraic $\mathbb{Z}^d$-actions, \emph{Invent. Math.} \textbf{142} (2000), 559--577.

\bibitem{Kopra2022} J.~Kopra, Glider automata on all transitive sofic shifts, \emph{Ergodic Theory Dynam. Systems} \textbf{42} (2022), 3716--3744.

\bibitem{Kurka1997} P.~K\r{u}rka, Languages, equicontinuity and attractors in cellular automata, \emph{Ergodic Theory Dynam. Systems} \textbf{17} (1997), 417--433.

\bibitem{Ledrappier1978} F.~Ledrappier, Un champ markovien peut \^etre d'entropie nulle et m\'elangeant, \emph{C. R. Acad. Sci. Paris S\'er. A-B} \textbf{287} (1978), A561--A563.

\bibitem{Ryan1972} J.~P.~Ryan, The shift and commutativity, \emph{Math. Systems Theory} \textbf{6} (1972), 82--85.

\bibitem{Salo2020} V.~Salo, Conjugacy of reversible cellular automata and one-head machines, preprint, arXiv:2011.07827 (2020).

\bibitem{Schmidt2002} K.~Schmidt, Algebraic $\mathbb{Z}^d$-actions, PIMS Distinguished Chair Lectures, University of Victoria, 2002.

\bibitem{Silverman2009} J.~H.~Silverman, \emph{The Arithmetic of Elliptic Curves}, 2nd ed., Grad. Texts in Math. 106, Springer, New York, 2009.

\bibitem{Takesue1987} S.~Takesue, Reversible cellular automata and statistical mechanics, \emph{Phys. Rev. Lett.} \textbf{59} (1987), 2499--2502.

\bibitem{Takesue1989} S.~Takesue, Ergodic properties and thermodynamic behavior of elementary reversible cellular automata. I. Basic properties, \emph{J. Stat. Phys.} \textbf{56} (1989), 371--402.

\bibitem{TM1990} T.~Toffoli and N.~Margolus, Invertible cellular automata: a review, \emph{Phys. D} \textbf{45} (1990), 229--253.

\bibitem{Ward1994} T.~Ward, Automorphisms of $\mathbb{Z}^d$-subshifts of finite type, \emph{Indag. Math. (N.S.)} \textbf{5} (1994), 495--504.

\bibitem{Salo2019Ryan} V.~Salo, Transitive action on finite points of a full shift and a finitary Ryan's theorem, \emph{Ergodic Theory Dynam. Systems} \textbf{39} (2019), 1637--1667.

\bibitem{KimRoush1990} K.~H.~Kim and F.~W.~Roush, On the automorphism groups of subshifts, \emph{Pure Math. Appl. Ser. B} \textbf{1} (1990), no.~4, 203--230.

\bibitem{MAM2013} G.~J.~Mart\'inez, A.~Adamatzky, F.~Chen and L.~Chua, On soliton collisions between localizations in complex elementary cellular automata: rules 54 and 110 and beyond, \emph{Complex Systems} \textbf{21} (2012), 117--142; arXiv:1301.6258.

\bibitem{BBGP1993} A.~Bobenko, M.~Bordemann, C.~Gunn and U.~Pinkall, On two integrable cellular automata, \emph{Comm. Math. Phys.} \textbf{158} (1993), 127--134.

\bibitem{BKP2021} B.~Bu\v{c}a, K.~Klobas and T.~Prosen, Rule 54: exactly solvable model of nonequilibrium statistical mechanics, \emph{J. Stat. Mech.} (2021), 074001.

\bibitem{Wolfram1984} S.~Wolfram, Universality and complexity in cellular automata, \emph{Phys. D} \textbf{10} (1984), 1--35.

\bibitem{Wolfram2002} S.~Wolfram, \emph{A New Kind of Science}, Wolfram Media, Champaign, IL, 2002.

\bibitem{PaperA} M.~Dayah, Reversible rule 110 has no hidden symmetries: six families of second-order cellular automata settled by edge moves, in preparation.

\bibitem{PaperB} M.~Dayah, Fronts at the speed of light: rigidity of permutive second-order cellular automata, in preparation.

\bibitem{PaperD} M.~Dayah, The 88 second-order elementary cellular automata: a classification by automorphism group, in preparation.

\bibitem{PaperE} M.~Dayah, The cellular automata that commute with Rule 54, and rigidity of the second-order automaton 94R, in preparation.

\end{thebibliography}


\end{document}
